% =====================================================================
% LEE'S COMPETITION MATHEMATICS WORKSHEET TEMPLATE
% =====================================================================
% HOW TO USE THIS FILE
% - Fill in the title fields below the packages. The author is Lee.
% - Set the two switches (solutions, working space). Usual editions:
%     Teacher copy: \showsolutionstrue  and \showspacefalse
%     Student copy: \showsolutionsfalse and \showspacetrue
% - Copy a unit block (notes -> example -> practice) for each theme.
% - Write every solution at the back with \solitem{label}{...}, in the
%   same order as the questions. Numbers (B1, P1, Example 1) are automatic.
% - Compiling (the worksheet uses Asymptote diagrams):
%     pdflatex file.tex
%     asy file-*.asy        (one .asy file per diagram)
%     pdflatex file.tex
%     pdflatex file.tex
%   Overleaf: choose pdfLaTeX and add a file named latexmkrc containing the
%   three lines below, each with the leading "% " removed:
% sub asy_diagram { return system('asy', '-o', $_[0], $_[0]); }
% add_cus_dep('asy', 'pdf', 0, 'asy_diagram');
% add_cus_dep('asy', 'tex', 0, 'asy_diagram');
%
% =====================================================================
% LEE'S PREFERENCES
% =====================================================================
% 1. AUDIENCE AND SCOPE
% - The student is strong at school mathematics but new to competitions.
%   Explain new reasoning carefully; skip elementary busywork.
% - Every question must really be about the worksheet's topic. Remove a
%   question that only touches the topic.
% - List (in a comment, near the start of the document) the topics saved
%   for later worksheets, and keep them out of this one.
% - No sources anywhere: no competition names, years or textbook references.
%
% 2. STRUCTURE
% - Sections are split into themed units (subsections). Each unit runs
%   concise notes -> titled worked example -> practice.
% - A section may open with "Getting started" (a worked Foundation box)
%   and a basic practice question (B1, B2, ...). Getting started uses
%   normal, representative calculations, not trivial ones.
% - The theme of a unit must be the most efficient way to solve each of
%   its questions. If another method is quicker, the question belongs in
%   that unit, or nowhere.
% - Every example's practice set always includes at least one variation of
%   the example (same type and form), so the student practises the
%   example's method at least once. The other practice questions vary the
%   structure, context or direction; do not just change the numbers.
% - Order units so that a skill is taught before it is needed (e.g. solving
%   equations before units whose questions need equation solving; the
%   algebra before the word problems that use it).
% - Easy items move earlier: a too-easy example becomes basic practice; a
%   too-easy practice question moves to an earlier unit or is dropped.
% - Put each question in the practice set of the unit whose method it
%   uses, rather than in a separate mixed Exercises section.
% - Add intermediate problems whenever the reasoning jumps too far.
% - Detailed solutions to every basic practice and practice question go at
%   the back, in the same order as the questions.
%
% 3. NOTES
% - Only essential definitions, notation, formulas and forms.
% - Proofs of formulas live in the notes, not in the examples:
%     "Why it works: picture"  - a diagram;
%     "Why it works: algebra"  - the cases, each with concrete numbers;
%     and, when it applies, how the idea builds on the previous unit
%     (e.g. "the same fact as the max formula, applied to the sizes").
% - When there are several methods (e.g. two cases, or squaring both
%   sides), list them and say which is usually quicker.
% - Define unfamiliar notation when it is introduced. For number theory,
%   state the positive-divisor convention if relevant.
%
% 4. EXAMPLES
% - A short title naming the theme or method.
% - A concise question with minimal set-up.
% - The solution goes in its own box. Add a takeaway box inside it only
%   when there is a reusable insight; do not force one.
% - When the notes already prove a formula, the example just applies it.
%
% 5. PRACTICE QUESTIONS
% - Levels, labelled on each question individually:
%     Foundation: direct use or straightforward variations.
%     Developing: an extra inference or a meaningful combination of ideas.
%     Extension:  a less obvious method or substantially deeper reasoning.
%   Large numbers or long arithmetic alone do not raise the level.
% - Mark a question that is harder than the rest with a star
%   (\starquestion), normally at Extension. Do NOT add any sentence
%   explaining the star.
% - Favour questions that APPLY the tools rather than just ask for them:
%   e.g. an unusual definition or set-up that looks complicated but turns
%   out to be something already taught. (Model: d(2021,n)+d(12345,n) with
%   a digit-by-digit "distance" d, which becomes "x between two points" at
%   every place.)
% - Prefer natural final-answer and word problems with real contexts
%   (e.g. lampposts beside a road) and natural wording (e.g. "the distance
%   from a on the number line"). Keep proof and impossibility questions
%   when appropriate. Do not force zero-answer counts.
% - Include questions that make the student check answers (e.g. drop an
%   outer absolute value so that one candidate fails).
% - Choose the working space for each question: short (the default), long
%   (several parts or a longer solution) or page (proofs).
% - Multiple-choice options: \fourchoices (one row) or \twochoices (two rows).
%
% 6. SOLUTION STYLE (worked examples and solutions at the back)
% - Point form: bold step headings (\steph{Step 1: ...}), arrows (\imp)
%   linking each deduction, short bullets, separate paragraphs and
%   displayed equations, plenty of space. Practice solutions are as
%   detailed as worked examples.
% - Write the way a person thinks: concrete numbers, diagrams and
%   "Without loss of generality, ..." rather than abstract algebra.
% - Justify every assumption the question does not state (e.g. putting B
%   to the right of A: look from the other side; distances do not change).
% - If the student has not met a tool yet (e.g. combinatorics), list the
%   concrete cases and count with simple multiplication.
% - Use a diagram instead of a counterexample or a long paragraph.
% - A rigorous proof that would interrupt the flow goes in a side note.
% - State a key fact first, then prove it, then use it.
% - For a sign decision, give both explanations:
%   \why{in words}{with inequalities}.
% - Check answers by substituting back, and show rejected candidates.
% - A slower alternative method appears only briefly, saying why the main
%   method is better.
% - Finish with a clear final answer.
%
% 7. FORMAT AND LAYOUT
% - Author: Lee (no surname).
% - Elegant black-and-white layout; standard Computer Modern fonts;
%   title, author, Name/Date lines and a contents page.
% - Fractions: \dfrac{a}{b}; in powers or superscripts: x^{\frac{1}{2}}.
%   Do not redefine \frac globally. Digit strings: \overline{12x34}.
% - Diagrams in Asymptote (helpers below: axis, tk, pt, gap).
% - No overfull lines: move long inline maths into displays or align*.
% - Every heading stays with the box that follows it. Reserve room BEFORE
%   the heading (\Needspace{...} just above \subsection), never between a
%   heading and its first box. A unit that starts straight with practice
%   needs \Needspace{\ifshowspace 0.4\textheight\else 14\baselineskip\fi}.
%   Check both editions for a heading left alone at the bottom of a page.
% - Allow long solutions to break across pages instead of forcing an
%   oversized block.
% - Generous space; there is no need to minimise the page count.
%
% 8. WHEN EDITING (for whoever prepares or revises the worksheet)
% - Do only what was asked. Do not add sentences, notes or explanations
%   to the worksheet that Lee did not ask for.
% - Keep authoring notes in comments, never in student-facing text.
% - When a question moves, its solution moves with it. Refer to questions
%   by label (\ref), never by a typed number.
% - A question cut because it needs another chapter goes into the question
%   bank (claude/question_bank.tex in the Wilson project), with its full
%   solution, under the chapter it needs. Check the bank first whenever a
%   worksheet on a new chapter is started.
% - Check every answer by calculation (e.g. a short computer search)
%   before it goes in.
% - Compile both editions with no overfull-box or undefined-reference
%   warnings, and look at the pages.
% =====================================================================

\documentclass[11pt,a4paper]{article}
\usepackage[utf8]{inputenc}
\usepackage[margin=0.65in,headheight=14pt]{geometry}
\usepackage{amsmath,amssymb,mathtools}
\usepackage{booktabs,tabularx,array,enumitem,needspace,etoolbox}
\usepackage[most]{tcolorbox}
\usepackage{xcolor,fancyhdr}
\usepackage[hidelinks,bookmarksopen,bookmarksnumbered]{hyperref}
\usepackage{asymptote}
% Compile the generated .asy files with Asymptote, then rerun pdfLaTeX.

% EDIT THESE FIELDS FOR EACH WORKSHEET.
\newcommand{\worksheettitle}{AP Calculus Unit 2: Differentiation:\texorpdfstring{\\}{ } Definition and Fundamental Properties}
\newcommand{\worksheetshorttitle}{Unit 2: Differentiation}
\newcommand{\worksheetauthor}{Lee}
\hypersetup{pdftitle={\worksheettitle},pdfauthor={\worksheetauthor}}

% THE TWO SWITCHES
% 1. Solutions. \showsolutionstrue prints the detailed solutions at the back;
%    \showsolutionsfalse hides them. Worked example solutions always stay.
% 2. Working space. \showspacetrue leaves room to write after every basic
%    practice and practice question; \showspacefalse gives a compact copy.
%    The size is chosen per question with the optional argument:
%      \question{...}         short: about 1/3 of a page (the default)
%      \question[long]{...}   long:  about 1/2 of a page
%      \question[page]{...}   page:  the question starts a new page and gets
%                             the whole page (for proofs)
% Usual editions: teacher copy = solutions on, space off;
%                 student copy = solutions off, space on.
\newif\ifshowsolutions
\showsolutionstrue
\newif\ifshowspace
\showspacetrue

\setlength{\parindent}{0pt}
\setlength{\parskip}{5pt}
\setlength{\emergencystretch}{2em}
\setlength{\jot}{5pt}
\setlist[enumerate]{leftmargin=*,itemsep=0.5em,topsep=0.35em}
\setlist[itemize]{leftmargin=1.5em,itemsep=3pt,topsep=3pt}
\renewcommand{\labelitemi}{$\bullet$}
\definecolor{notegray}{RGB}{245,245,245}
\definecolor{rulegray}{RGB}{232,232,232}
\definecolor{darkgray}{RGB}{70,70,70}
\definecolor{brightorange}{RGB}{255,128,0}
\pagestyle{fancy}
\fancyhf{}
\lhead{\worksheetshorttitle}
\rhead{\hyperlink{contents}{Contents}\quad\thepage}
\renewcommand{\headrulewidth}{0.4pt}
\setcounter{secnumdepth}{2}
\setcounter{tocdepth}{2}
\raggedbottom

\newtcolorbox{notesbox}[1]{
 enhanced,breakable,colback=notegray,colframe=darkgray,
 boxrule=0.6pt,arc=2mm,title={#1},fonttitle=\bfseries,coltitle=black,
 attach boxed title to top left={xshift=2mm,yshift=-1.5mm},
 boxed title style={colback=white,colframe=darkgray,boxrule=0.6pt,arc=1mm},
 top=2.5mm,before skip=10pt,after skip=10pt
}
\newtcolorbox{questionbox}[2][]{
 enhanced,breakable,colback=notegray,colframe=black,
 colbacktitle=notegray,coltitle=black,title={#2},fonttitle=\bfseries,
 boxrule=0.65pt,titlerule=0.4pt,arc=1.2mm,
 left=3mm,right=3mm,top=2mm,bottom=2mm,
 before skip=5pt,after skip=10pt,#1
}
\newtcolorbox{examplesolution}[1][]{
 enhanced,breakable,colback=white,colframe=darkgray,
 boxrule=0.5pt,arc=1mm,left=2mm,right=2mm,top=3mm,bottom=2mm,
 title={Solution},fonttitle=\small,coltitle=black,colbacktitle=white,
 attach boxed title to top left={xshift=2mm,yshift=-1.5mm},
 boxed title style={colback=white,colframe=darkgray,boxrule=0.5pt,arc=0.7mm},
 title after break={Solution (continued)},
 before skip=5pt,after skip=7pt,#1
}
\newtcolorbox{takeaway}{
 colback=notegray,colframe=darkgray,boxrule=0.6pt,arc=1mm,
 left=2mm,right=2mm,top=1.5mm,bottom=1.5mm,
 before skip=8pt,after skip=4pt
}
\newcommand{\leveltag}[1]{{\normalfont\footnotesize\textsc{#1}}}

% Automatic numbering; use unique labels to keep solutions matched when moved.
\newcounter{example}
\newcounter{basic}
\newcounter{practice}
\newcounter{exercise}
\renewcommand{\thebasic}{B\arabic{basic}}
\renewcommand{\thepractice}{P\arabic{practice}}
\renewcommand{\theexercise}{\arabic{section}.\arabic{exercise}}
\makeatletter\@addtoreset{exercise}{section}\makeatother
\newcommand{\exampletitle}[2]{%
 \Needspace{8\baselineskip}\refstepcounter{example}\label{#1}%
 \begin{questionbox}{Example \theexample\quad #2}}
% Close an \exampletitle question with \end{questionbox}.
% \question[size]{counter}{label}{level}{text}, then the parts, then
% \workingspace. size = short (default), long or page; see the switches.
% \starquestion: same as \question, with a star marking a harder question.
\makeatletter
\newlength{\ws@short}
\newlength{\ws@long}
\newlength{\ws@need}
\AtBeginDocument{%
 \setlength{\ws@short}{0.33\textheight}%
 \setlength{\ws@long}{0.5\textheight}}
\newcommand{\ws@kind}{short}
\newcommand{\ws@pending}{}
% The Practice heading is printed together with the first question after it,
% so it is never left alone at the bottom of a page.
\newcommand{\practiceheading}{\par\gdef\ws@pending{%
 \medskip{\large\bfseries Practice}\par\smallskip}}
\newcommand{\ws@start}[1]{%
 \par\gdef\ws@kind{#1}%
 \ifx\ws@pending\@empty\setlength{\ws@need}{0pt}%
 \else\setlength{\ws@need}{3\baselineskip}\fi
 \ifshowspace
  \ifdefstring{\ws@kind}{page}{\clearpage}{%
   \ifdefstring{\ws@kind}{long}%
    {\addtolength{\ws@need}{\ws@long}}{\addtolength{\ws@need}{\ws@short}}%
   \addtolength{\ws@need}{6\baselineskip}%
   \Needspace{\ws@need}}%
 \else
  \addtolength{\ws@need}{5\baselineskip}%
  \Needspace{\ws@need}%
 \fi
 \ws@pending\gdef\ws@pending{}}
\newcommand{\question}[5][short]{%
 \ws@start{#1}\refstepcounter{#2}\label{#3}%
 \textbf{\csname the#2\endcsname}\showsrc{#3}\quad #5\par}
\newcommand{\starquestion}[5][short]{%
 \ws@start{#1}\refstepcounter{#2}%
 \protected@edef\@currentlabel{\csname the#2\endcsname$^{\bigstar}$}\label{#3}%
 \textbf{\csname the#2\endcsname}$^{\bigstar}$\showsrc{#3}\quad #5\par}
\newcommand{\workingspace}{\par
 \ifshowspace
  \ifdefstring{\ws@kind}{page}{\clearpage}{%
   \ifdefstring{\ws@kind}{long}{\vspace{\ws@long}}{\vspace{\ws@short}}}%
 \else\vspace{8mm}\fi}
\makeatother
\newcommand{\solitem}[2]{%
 \par\addvspace{16pt}\Needspace{10\baselineskip}%
 {\large\bfseries \ref{#1}}\par\nobreak\vspace{2pt}%
 {\color{rulegray}\hrule height 0.5pt}\nobreak\vspace{7pt}%
 #2\par}


% Extra helpers for this worksheet.
% \DNE: "does not exist" in tables and answers.
\newcommand{\DNE}{\text{DNE}}
\newcommand{\DNEinf}{\ensuremath{\text{DNE}/\infty}}
\newcommand{\DNEninf}{\ensuremath{\text{DNE}/{-\infty}}}
% Written-solution helpers:
%  \why{...}      side note in an align* (use  &&\why{...}  at the end of a line)
%  \scratch{...}  a line of scratch-paper thinking (not part of the written answer)
%  \twosides{left}{right}  two columns, e.g. x->a^- on the left and x->a^+ on the right
\newcommand{\why}[1]{\text{\small\color{darkgray}#1}}
\newcommand{\scratch}[1]{{\small\color{darkgray}\textit{Scratch:} #1}\par}
\newcommand{\ontest}{{\small\textit{On the test:}}\ }
\newcommand{\twosides}[2]{\par\smallskip\noindent\begin{minipage}[t]{0.47\linewidth}#1\end{minipage}\hfill{\color{rulegray}\vrule width 0.6pt}\hfill\begin{minipage}[t]{0.47\linewidth}#2\end{minipage}\par\smallskip}
% Empty cell of a fill-in table.
\newcommand{\blankcell}{\rule{0pt}{18pt}}
% Calculator marker at the start of a question.
\newcommand{\calcicon}{\tikz[baseline={(0,-0.6pt)},x=1pt,y=1pt]{%
  \draw[line width=0.6pt,rounded corners=0.8pt] (0,0) rectangle (8,10);
  \fill (1.5,6.7) rectangle (6.5,8.6);
  \foreach \i in {0,1,2} {\foreach \j in {0,1,2} {\fill (1.5+\i*1.95,1.3+\j*1.75) rectangle ++(1.1,1.05);}}}}
\newcommand{\calc}{\calcicon\enspace}
% Source tags: \defsrc{label}{AB 2013} prints a small tag after the level.
\newcommand{\defsrc}[2]{\expandafter\def\csname qsrc@#1\endcsname{#2}}
\newcommand{\srctag}[1]{{\normalfont\footnotesize\color{darkgray}[#1]}}
\newcommand{\showsrc}[1]{\ifcsname qsrc@#1\endcsname\quad\srctag{\csname qsrc@#1\endcsname}\fi}
% Roman-numeral statements I, II, III.
\newlist{statements}{enumerate}{1}
\setlist[statements]{label=\Roman*.,leftmargin=3em,itemsep=2pt,topsep=4pt}

% Helpers for questions and solutions.
% \ph{...}: placeholder text in this template; replace it when writing.
\newcommand{\ph}[1]{\textit{[#1]}}
% Solution layout: \steph{...} starts a bold step heading on its own line;
% \imp is a spaced "implies" arrow for point-form reasoning.
\newcommand{\steph}[1]{\par\addvspace{7pt}{\bfseries #1}\par\nobreak\vspace{2pt}}
\newcommand{\imp}{\ \Longrightarrow\ }
% Multiple-choice options in one row, or in two rows for longer options.
\newcommand{\fourchoices}[4]{\par\smallskip\noindent
 \begin{tabularx}{\linewidth}{@{}XXXX@{}}
 (A)\ #1 & (B)\ #2 & (C)\ #3 & (D)\ #4
 \end{tabularx}\par}
\newcommand{\twochoices}[4]{\par\smallskip\noindent
 \begin{tabularx}{\linewidth}{@{}XX@{}}
 (A)\ #1 & (B)\ #2\\[10pt]
 (C)\ #3 & (D)\ #4
 \end{tabularx}\par}

% Shared Asymptote helpers for number-line diagrams.
\begin{asydef}
defaultpen(fontsize(10pt)+linewidth(0.6pt));
// Axis from lo to hi at height y, with an arrow on the right.
void axis(real lo, real hi, real y=0) {draw((lo,y)--(hi,y), EndArrow(4));}
// Tick mark at x with optional labels below and above.
void tk(real x, string below="", string above="", real y=0) {
  draw((x,y-0.12)--(x,y+0.12));
  if (below != "") label(below, (x,y-0.12), S);
  if (above != "") label(above, (x,y+0.12), N);
}
// Point at x with optional labels below and above.
void pt(real x, string below="", string above="", real y=0) {
  dot((x,y), linewidth(3.6pt));
  if (below != "") label(below, (x,y-0.1), S);
  if (above != "") label(above, (x,y+0.1), N);
}
// Double-headed measuring arrow from a to b at height h above the axis.
void gap(real a, real b, real h, string s, real y=0) {
  draw((a,y+0.15)--(a,y+h+0.08), gray(0.55));
  draw((b,y+0.15)--(b,y+h+0.08), gray(0.55));
  draw((a,y+h)--(b,y+h), Arrows(3));
  if (s != "") label(s, ((a+b)/2,y+h), N);
}
// ---------------------------------------------------------------
// Graph helpers (coordinate graphs).
// Light grid on integer points.
void grid(real x0, real x1, real y0, real y1, real step=1) {
  for (real x=ceil(x0/step)*step; x<=x1+1e-9; x+=step) draw((x,y0)--(x,y1), gray(0.8)+linewidth(0.4pt));
  for (real y=ceil(y0/step)*step; y<=y1+1e-9; y+=step) draw((x0,y)--(x1,y), gray(0.8)+linewidth(0.4pt));
}
// x- and y-axes with arrows and labels.
void axes(real x0, real x1, real y0, real y1, string xl="$x$", string yl="$y$") {
  draw((x0,0)--(x1,0), EndArrow(4)); draw((0,y0)--(0,y1), EndArrow(4));
  label(xl,(x1,0),E); label(yl,(0,y1),N);
}
// Tick marks with number labels on the axes.
void xnums(real[] xs) {for (real x: xs) {draw((x,0),(0,2.5)--(0,-2.5)); label("$"+string(x)+"$",(x,0),2S,fontsize(8pt));}}
void ynums(real[] ys) {for (real y: ys) {draw((0,y),(2.5,0)--(-2.5,0)); label("$"+string(y)+"$",(0,y),2W,fontsize(8pt));}}
// Tick mark with a text label (e.g. "$a$") on the x-axis.
void xlab(real x, string s) {draw((x,0),(0,2.5)--(0,-2.5)); label(s,(x,0),2S,fontsize(9pt));}
// Open and filled dots.
void opendot(pair p) {dot(p, linewidth(4.4pt)); dot(p, white+linewidth(2.8pt));}
void closeddot(pair p) {dot(p, linewidth(4.4pt));}
// Graph of y = f(x) for a <= x <= b, and a line segment.
void curve(real f(real), real a, real b, int n=100) {
  guide g; for (int i=0; i<=n; ++i) {real x=a+(b-a)*i/n; g=g..(x,f(x));}
  draw(g, linewidth(0.9pt));
}
void seg(pair p, pair q) {draw(p--q, linewidth(0.9pt));}
// Polyline graph with many points (for steep curves; clip afterwards).
void pcurve(real f(real), real a, real b, int n=600) {
  guide g; for (int i=0; i<=n; ++i) {real x=a+(b-a)*i/n; g=g--(x,f(x));}
  draw(g, linewidth(0.9pt));
}
// Dashed asymptotes.
void vasym(real a, real y0, real y1) {draw((a,y0)--(a,y1), dashed+linewidth(0.5pt));}
void hasym(real b, real x0, real x1) {draw((x0,b)--(x1,b), dashed+linewidth(0.5pt));}
\end{asydef}


% Source tags for practice and exercise questions (from the coverage map).
% Section 1
\defsrc{pr:fifths}{BC 2013 Alt}
\defsrc{pr:negpower}{AB 2014 Intl}
\defsrc{pr:quotcubic}{AB 2017 Alt}
\defsrc{pr:quotlinear}{BC 2008 Form B}
\defsrc{pr:quotfg}{BC 2021}
\defsrc{ex1:cubic}{AB 2013 Intl}
\defsrc{ex1:frq}{BC 2016 Practice FRQ}
\defsrc{ex1:threefns}{BC 2024}
% Section 2
% Section 2 (helper A: 2.1-2.3)
\defsrc{a:pr:arocpiece}{BC 2021}
\defsrc{a:pr:plane}{BC 2023}
\defsrc{a:pr:defexpr}{BC 2017 Alt}
\defsrc{a:pr:dqcube}{BC 2018 Alt}
\defsrc{a:pr:dqgraph}{BC 2022 Intl}
\defsrc{a:pr:dqtable}{BC 2024 Alt FRQ}
\defsrc{a:ex2:count}{AB 2009 Form B FRQ}
% Section 2 (helper B): tangent lines, slopes with the rules, estimating derivatives
\defsrc{b:pr:twopoints}{AB 2003}
\defsrc{b:pr:tanlinegraph}{BC 2019 Alt}
\defsrc{b:pr:limitslope}{BC 2024 Alt}
\defsrc{b:pr:tangentk}{BC 2008}
\defsrc{b:pr:parallelb}{AB 2018 Practice}
\defsrc{b:pr:avgrate}{BC 2003}
\defsrc{b:pr:tickets}{AB 2008 FRQ}
\defsrc{b:pr:tea}{AB 2011 FRQ}
\defsrc{b:pr:orderslopes}{AB 2017 Alt}
\defsrc{b:ex2:slopefield}{Based on AB 1998 FRQ}
\defsrc{b:ex2:xn}{AB 2004 Form B FRQ}
\defsrc{b:ex2:order}{AB 2013 Intl}
\defsrc{b:ex2:fastest}{BC 2022 Intl}
% Section 2 (helper C): rates in context, derivatives on the calculator
\defsrc{c:tankw}{AB 2017 Alt}
\defsrc{c:wave}{AB 2015 Intl}
\defsrc{c:etan}{AB 2012 Intl}
\defsrc{c:rainpiece}{BC 2024 Alt}
\defsrc{c:count}{BC 2021}
\defsrc{c:origin}{BC 2014 Alt}
\defsrc{c:parallel}{AB 2017 Alt}
\defsrc{c:xshoppers}{Based on AB 2014 Intl}
\defsrc{c:xwindchill}{AB 2007 Form B}
\defsrc{c:xexpquartic}{BC 2012}
\defsrc{a:pr:costwox}{AB 2015 Intl}
\defsrc{a:pr:fprimerow}{BC 2024 Alt FRQ}
\defsrc{a:pr:arocgraph}{AB 2016 FRQ}
% Section 3
% Section 3 (helper D: 3.1 Where a derivative does not exist)
\defsrc{d:pr:vtan}{AB 1998}
\defsrc{d:pr:contnotdiff}{AB 2003}
\defsrc{d:pr:rootabs}{AB 2012 Intl}
\defsrc{d:pr:whichgraph}{BC 2021}
\defsrc{d:pr:limfprime}{BC 2003}
\defsrc{d:pr:gsemicircle}{BC 2018 Alt FRQ}
\defsrc{d:pr:onesided}{AB 2009 Form B FRQ}
\defsrc{d:pr:velocity}{AB 2005 FRQ}
\defsrc{f:pr:mustbe}{AB 2017 Alt}
\defsrc{f:pr:couldfalse}{BC 2008 Form B}
\defsrc{f:pr:notdiff}{BC 2019 Alt}
\defsrc{f:pr:cont3}{BC 2012}
\defsrc{f:pr:graphhole}{BC 2017 Alt}
\defsrc{f:pr:limtable}{AB 2026 FRQ}
\defsrc{f:pr:ivtW}{AB 2013 Intl FRQ}
\defsrc{f:pr:ivtr}{AB 2022 FRQ}
% Section 3 (helper E: 3.3)
\defsrc{e:pr:ab2003}{AB 2003}
\defsrc{e:pr:removable}{BC 2008}
\defsrc{e:pr:samefn}{AB 2002 FRQ}
% Section 3 (helper G: 3.4)
\defsrc{g:pr:cd}{BC 2008}
\begin{document}
\begin{center}
 {\LARGE\bfseries \worksheettitle\par}\vspace{5pt}
 {\large \worksheetauthor}\\[9pt]
 \begin{tabularx}{\textwidth}{@{}X X@{}}
 Name: \rule{2.2in}{0.4pt} & Date: \rule{1.8in}{0.4pt}
 \end{tabularx}
\end{center}
\vspace{12pt}
\hypertarget{contents}{}%
{\setlength{\parskip}{0pt}\tableofcontents}
\clearpage

% Topics deliberately left for later worksheets (keep them out of this one):
% - Chain rule, implicit differentiation, derivatives of inverse functions
%   and inverse trig functions, higher-order derivatives (Unit 3)
% - Straight-line motion (velocity, acceleration, at rest), the meaning of
%   a derivative in context, "increasing because f' > 0", related rates,
%   linearization, L'Hospital's rule (Unit 4)
% - Mean Value Theorem, increasing/decreasing, extrema, concavity, and
%   matching the graphs of f and f' (Unit 5)
% - Integrals and accumulation functions (Unit 6 onwards)
% Questions cut from this worksheet for these units, with solutions:
% claude/question_bank.tex in the Calculus BC project.
%
% Order (Lee): procedure first. Section 1 = the rules (CED 2.5-2.10);
% Section 2 = what the derivative means (CED 2.1-2.3); Section 3 =
% differentiability and continuity (CED 2.4). Each section ends with
% Exercises.
% =====================================================================
\section{Derivative Rules}\label{sec:rules}
% =====================================================================
% REMOVED AP QUESTIONS (Lee, Nov 2026). These past-paper questions were cut because similar questions
% already do the job. They are NOT anywhere else: not in this worksheet and not in the question bank.
% Similar questions that stay are shown after the arrow.
% - AB 2015 Intl MC 2 (sqrt x + 3/sqrt x, f'(4)); AB 1988 MC 15 (sqrt(2x), f'(2));
%   the question based on BC 2022 Intl MC 9 (12x^{1/2} + 3x + 32/x + 5, f'(4))  ->  Example 1, P1-P3.
% - AB 1988 MC 12 (sin x, f'(pi/3)); AB 2012 Intl MC 4 (7x - 3 + ln x, f'(1)); AB 2019 Alt MC 7 (x^5 - 5^x)
%   ->  Example 2, the basic-functions practice.
% - AB 1988 MC 1 (x^2 e^x); AB 2012 Intl MC 1 (x sin x); AB 2015 Intl MC 10 (sin x cos x at pi/3)
%   ->  Example 3, the product-rule practice.
% - AB 1988 MC 6 ((ln x)/x); BC 2022 Intl MC 2 (x^3/(x + 2))  ->  Example 4, the quotient-rule practice.
% Example parts removed (new questions, not past papers): Example 1(d), Example 3(e), Example 4(a) and their practice twins.
% =====================================================================
% =====================================================================
% CED 2.5-2.10. Procedure first: the rules are practised here, before the
% definition and the meaning of the derivative (Section 2).
\Needspace{20\baselineskip}
\subsection{Power rule, sums and constant multiples}
\begin{notesbox}{Notes}
\begin{itemize}
 \item The derivative of $f$ is a new function $f'$. The number $f'(a)$ is the slope of the graph of $f$ at $x=a$ (the next section explains why). For $y=f(x)$ these all mean the derivative:
 \[
 f'(x),\qquad y',\qquad \frac{dy}{dx},\qquad \frac{d}{dx}f(x).
 \]
 \item The rules ($c$ and $n$ are constants):
 \[
 \begin{aligned}
 &\frac{d}{dx}(c)=0, &\qquad &\frac{d}{dx}\bigl(x^n\bigr)=nx^{n-1},\\[4pt]
 &\frac{d}{dx}\bigl(cf(x)\bigr)=cf'(x), &\qquad &\frac{d}{dx}\bigl(f(x)\pm g(x)\bigr)=f'(x)\pm g'(x).
 \end{aligned}
 \]
\end{itemize}
\end{notesbox}

\exampletitle{ex:power}{The power rule}
Find the derivative of each function.
\begin{enumerate}[label=(\alph*)]
 \item $y=6\sqrt x-\dfrac{4}{x^2}+\dfrac{5}{3x}+\pi^2$
 \item $y=\bigl(3x^2-1\bigr)^2$
 \item $f(x)=\dfrac{x^2-4x+3}{\sqrt x}$.\quad Find $f'(1)$.
\end{enumerate}
\end{questionbox}
\begin{examplesolution}
\relax
\begin{enumerate}[label=(\alph*)]
 \item
 \begin{align*}
 y&=6x^{\frac12}-4x^{-2}+\frac53x^{-1}+\pi^2 &&\why{rewrite every term as $cx^n$}\\
 \frac{dy}{dx}&=6\cdot\frac12x^{-\frac12}-4(-2)x^{-3}+\frac53(-1)x^{-2}+0 &&\why{power rule; $\pi^2$ is a constant}\\
 &=\frac{3}{\sqrt x}+\frac{8}{x^3}-\frac{5}{3x^2}.
 \end{align*}
 \item A bracket is not $x$, so the power rule does not apply to it. Expand first.
 \begin{align*}
 y&=9x^4-6x^2+1 &&\why{expand}\\
 \frac{dy}{dx}&=36x^3-12x.
 \end{align*}
 The trap answer $2\bigl(3x^2-1\bigr)$ treats the bracket as $x$.
 \item One term on the bottom: split. Then differentiate, then substitute.
 \begin{align*}
 f(x)&=\frac{x^2}{x^{\frac12}}-\frac{4x}{x^{\frac12}}+\frac{3}{x^{\frac12}}=x^{\frac32}-4x^{\frac12}+3x^{-\frac12} &&\why{split; subtract the powers}\\
 f'(x)&=\frac32x^{\frac12}-2x^{-\frac12}-\frac32x^{-\frac32}\\
 f'(1)&=\frac32-2-\frac32=-2. &&\why{every power of $1$ is $1$}
 \end{align*}
 The trap answer is $f(1)=\dfrac{1-4+3}{1}=0$.
\end{enumerate}
\begin{takeaway}
Rewrite first, then differentiate: every term must look like $cx^n$ before the power rule is used.
\end{takeaway}
\end{examplesolution}

\practiceheading
\question[long]{practice}{pr:power}{}{Find the derivative of each function.
\begin{enumerate}[label=(\alph*)]
 \item $y=8\sqrt[3]{x}+\dfrac{2}{x^3}-\dfrac{7}{4x}+\sqrt3$
 \item $y=\bigl(2x^2-1\bigr)^2$
 \item $f(x)=\dfrac{x^3-2x+5}{x^2}$.\quad Find $f'(2)$.
\end{enumerate}}
\workingspace

\question{practice}{pr:fifths}{}{If $y=\dfrac12x^{\frac45}-\dfrac{3}{x^5}$, then $\dfrac{dy}{dx}=$}
\twochoices{$\dfrac{2}{5x^{\frac15}}+\dfrac{15}{x^6}$}{$\dfrac{2}{5x^{\frac15}}+\dfrac{15}{x^4}$}{$\dfrac{2}{5x^{\frac15}}-\dfrac{3}{5x^4}$}{$\dfrac{2x^{\frac15}}{5}+\dfrac{15}{x^6}$}
\workingspace

\question{practice}{pr:negpower}{}{If $f(x)=4x^{-2}+\dfrac14x^2+4$, then $f'(2)=$}
\fourchoices{$-58$}{$-3$}{$0$}{$1$}
\workingspace

% AB 2003 MC 1, (x^3+1)^2, moved to the question bank (3.1 chain rule): its answers are in chain-rule form (Lee).
% ---------------------------------------------------------------------
\Needspace{20\baselineskip}
\subsection{\texorpdfstring{Derivatives of $\sin x$, $\cos x$, $e^x$, $\ln x$ and $b^x$}{Derivatives of sin x, cos x, e\textasciicircum x, ln x and b\textasciicircum x}}
\begin{notesbox}{Notes}
\begin{itemize}
 \item The derivatives (angles in radians; $b>0$ is a constant):
 \[
 \begin{aligned}
 &\frac{d}{dx}(\sin x)=\cos x, &\qquad &\frac{d}{dx}(\cos x)=-\sin x,\\[4pt]
 &\frac{d}{dx}\bigl(e^x\bigr)=e^x, &\qquad &\frac{d}{dx}(\ln x)=\frac1x, &\qquad &\frac{d}{dx}\bigl(b^x\bigr)=b^x\ln b.
 \end{aligned}
 \]
 \item $\ln$ of a constant times $x$: split it first, $\ln(ax)=\ln a+\ln x$. Then $\ln a$ is just a constant, so $\dfrac{d}{dx}\ln(4x)=\dfrac{d}{dx}(\ln4+\ln x)=\dfrac1x$.
 \item Numbers such as $e^2$, $\ln7$ and $\sin\dfrac{\pi}{6}$ are constants: their derivative is $0$.
\end{itemize}
\steph{Traps.}
\begin{itemize}
 \item The minus sign is on the derivative of $\cos x$, not of $\sin x$.
 \item $b^x$ is not a power of $x$: $\dfrac{d}{dx}\bigl(5^x\bigr)=5^x\ln5$, not $x\cdot5^{x-1}$.
 \item $f'(a)$, not $f(a)$: for $f(x)=\sin x$, $f'\Bigl(\dfrac{\pi}{3}\Bigr)=\cos\dfrac{\pi}{3}=\dfrac12$, while $f\Bigl(\dfrac{\pi}{3}\Bigr)=\dfrac{\sqrt3}{2}$.
\end{itemize}
\end{notesbox}

\Needspace{0.45\textheight}
\exampletitle{ex:basicfns}{The basic functions}
Find the derivative of each function.
\begin{enumerate}[label=(\alph*)]
 \item $f(x)=3\sin x-2\cos x+4e^x$.\quad Find $f'(0)$.
 \item $y=5\ln x-\dfrac2x+\ln7$
 \item $y=3^x+x^3+e^3$
 \item $f(x)=\ln(4x)-\sqrt x$.\quad Find $f'(1)$.
 \item $f(x)=2\sin x+\cos x$.\quad Find $f'\Bigl(\dfrac{\pi}{3}\Bigr)$.
\end{enumerate}
\end{questionbox}
\begin{examplesolution}
\relax
\begin{enumerate}[label=(\alph*)]
 \item Differentiate first, then substitute.
 \begin{align*}
 f'(x)&=3\cos x-2(-\sin x)+4e^x=3\cos x+2\sin x+4e^x\\
 f'(0)&=3(1)+2(0)+4(1)=7. &&\why{$\cos0=1$, $\sin0=0$, $e^0=1$}
 \end{align*}
 The trap answer is $f(0)=0-2+4=2$.
 \item
 \begin{align*}
 y&=5\ln x-2x^{-1}+\ln7 &&\why{$\ln7$ is a constant}\\
 \frac{dy}{dx}&=\frac5x+2x^{-2}+0=\frac5x+\frac{2}{x^2}.
 \end{align*}
 \item $3^x$ is an exponential, $x^3$ is a power and $e^3$ is a constant.
 \[
 \frac{dy}{dx}=3^x\ln3+3x^2+0=3^x\ln3+3x^2.
 \]
 \item Split the $\ln$ first.
 \begin{align*}
 f(x)&=\ln4+\ln x-x^{\frac12} &&\why{$\ln(4x)=\ln4+\ln x$}\\
 f'(x)&=0+\frac1x-\frac12x^{-\frac12}=\frac1x-\frac{1}{2\sqrt x} &&\why{$\ln4$ is a constant}\\
 f'(1)&=1-\frac12=\frac12.
 \end{align*}
 The trap $\dfrac{d}{dx}\ln(4x)=\dfrac{1}{4x}$ gives $\dfrac14-\dfrac12=-\dfrac14$.
 \item Differentiate first, then substitute.
 \begin{align*}
 f'(x)&=2\cos x-\sin x\\
 f'\Bigl(\frac{\pi}{3}\Bigr)&=2\cdot\frac12-\frac{\sqrt3}{2}=1-\frac{\sqrt3}{2}.
 \end{align*}
\end{enumerate}
\end{examplesolution}

\practiceheading
\question[long]{practice}{pr:basicfns}{}{Find the derivative of each function.
\begin{enumerate}[label=(\alph*)]
 \item $f(x)=4\cos x-5\sin x+2e^x-e^2$.\quad Find $f'(0)$.
 \item $y=3\ln x+\dfrac4x-\ln2$
 \item $y=2^x-x^2+e$
 \item $f(x)=\ln(5x)+4\sqrt x$.\quad Find $f'(4)$.
 \item $f(x)=3\cos x-2\sin x$.\quad Find $f'\Bigl(\dfrac{\pi}{6}\Bigr)$.
\end{enumerate}}
\workingspace

% ---------------------------------------------------------------------
\Needspace{20\baselineskip}
\subsection{The product rule}
\begin{notesbox}{Notes}
\begin{itemize}
 \item Product rule: differentiate one factor and leave the other alone; then differentiate the other factor and leave the first alone; add.
 \[
 \frac{d}{dx}\bigl(f(x)g(x)\bigr)=\frac{d}{dx}\bigl(f(x)\bigr)\cdot g(x)+\frac{d}{dx}\bigl(g(x)\bigr)\cdot f(x)=f'(x)g(x)+g'(x)f(x).
 \]
 \item Three factors (or more): the same idea. Differentiate one factor at a time and leave the others alone, once for each factor, then add. There is one term for each factor.
 \[
 \frac{d}{dx}\bigl(f(x)g(x)h(x)\bigr)=\frac{d}{dx}\bigl(f(x)\bigr)\cdot g(x)h(x)+\frac{d}{dx}\bigl(g(x)\bigr)\cdot f(x)h(x)+\frac{d}{dx}\bigl(h(x)\bigr)\cdot f(x)g(x).
 \]
 \item Write the first line of the working in this form, so that you can see which factor is being differentiated. For $y=x^5\ln x$:
 \[
 \frac{dy}{dx}=\frac{d}{dx}\bigl(x^5\bigr)\cdot\ln x+\frac{d}{dx}(\ln x)\cdot x^5=5x^4\ln x+\frac1x\cdot x^5=5x^4\ln x+x^4.
 \]
 \item Tidy the answer by factoring out a common factor; answer choices are often factored, e.g.\ $2xe^x+x^2e^x=xe^x(x+2)$.
 \item Two polynomials: expanding first is just as good, and often quicker. With $e^x$, $b^x$, $\sin x$, $\cos x$ or $\ln x$, use the product rule.
 \item A constant times a function is not a product rule question: $\dfrac{d}{dx}(5\sin x)=5\cos x$.
\end{itemize}
\steph{Trap.}
\begin{itemize}
 \item The derivative of a product is not the product of the derivatives: $\dfrac{d}{dx}\bigl(x^2e^x\bigr)\ne2x\cdot e^x$.
\end{itemize}
\end{notesbox}

\Needspace{0.45\textheight}
\exampletitle{ex:product}{The product rule}
Find the derivative of each function.
\begin{enumerate}[label=(\alph*)]
 \item $y=x^3\cos x$
 \item $f(x)=e^x\bigl(x^2-3x\bigr)$.\quad Find $f'(1)$.
 \item $y=\sqrt x\,\ln x$
 \item $y=\bigl(x^2+1\bigr)(3x-5)$
 \item $f(x)=x^2e^x\ln x$.\quad Find $f'(1)$.
\end{enumerate}
\end{questionbox}
\begin{examplesolution}
\relax
\begin{enumerate}[label=(\alph*)]
 \item
 \begin{align*}
 \frac{dy}{dx}&=\frac{d}{dx}\bigl(x^3\bigr)\cdot\cos x+\frac{d}{dx}(\cos x)\cdot x^3 &&\why{product rule}\\
 &=3x^2\cos x+(-\sin x)\,x^3\\
 &=3x^2\cos x-x^3\sin x.
 \end{align*}
 \item
 \begin{align*}
 f'(x)&=\frac{d}{dx}\bigl(e^x\bigr)\cdot\bigl(x^2-3x\bigr)+\frac{d}{dx}\bigl(x^2-3x\bigr)\cdot e^x &&\why{product rule}\\
 &=e^x\bigl(x^2-3x\bigr)+(2x-3)e^x\\
 &=e^x\bigl(x^2-3x+2x-3\bigr)=e^x\bigl(x^2-x-3\bigr) &&\why{factor out $e^x$}\\
 f'(1)&=e(1-1-3)=-3e.
 \end{align*}
 The trap answer $-e$ comes from the product of the derivatives, $e^x(2x-3)$.
 \item
 \begin{align*}
 \frac{dy}{dx}&=\frac{d}{dx}\bigl(x^{\frac12}\bigr)\cdot\ln x+\frac{d}{dx}(\ln x)\cdot\sqrt x &&\why{product rule}\\
 &=\frac{1}{2\sqrt x}\cdot\ln x+\frac1x\cdot\sqrt x\\
 &=\frac{\ln x}{2\sqrt x}+\frac{1}{\sqrt x} &&\why{$\dfrac{\sqrt x}{x}=\dfrac{1}{\sqrt x}$}\\
 &=\frac{\ln x+2}{2\sqrt x}.
 \end{align*}
 \item Two polynomials: either method works.
 \twosides{\textbf{Product rule}
 \begin{align*}
 \frac{dy}{dx}&=\frac{d}{dx}\bigl(x^2+1\bigr)\cdot(3x-5)\\
 &\qquad+\frac{d}{dx}(3x-5)\cdot\bigl(x^2+1\bigr)\\
 &=2x(3x-5)+3\bigl(x^2+1\bigr)\\
 &=9x^2-10x+3.
 \end{align*}}{\textbf{Expand first}
 \begin{align*}
 y&=3x^3-5x^2+3x-5\\
 \frac{dy}{dx}&=9x^2-10x+3.
 \end{align*}}
 \item Three factors: three terms. Each time, differentiate one factor and leave the other two alone.
 \begin{align*}
 f'(x)&=\frac{d}{dx}\bigl(x^2\bigr)\cdot e^x\ln x+\frac{d}{dx}\bigl(e^x\bigr)\cdot x^2\ln x+\frac{d}{dx}(\ln x)\cdot x^2e^x &&\why{product rule, three factors}\\
 &=2xe^x\ln x+e^x\cdot x^2\ln x+\frac1x\cdot x^2e^x\\
 &=2xe^x\ln x+x^2e^x\ln x+xe^x\\
 f'(1)&=0+0+e=e. &&\why{$\ln1=0$}
 \end{align*}
\end{enumerate}
\begin{takeaway}
Differentiate one factor at a time and leave the others alone: one term for each factor. Factor the answer to match the choices.
\end{takeaway}
\end{examplesolution}

\practiceheading
\question[long]{practice}{pr:product}{}{Find the derivative of each function.
\begin{enumerate}[label=(\alph*)]
 \item $y=x^4\sin x$
 \item $f(x)=e^x\bigl(x^2+2x\bigr)$.\quad Find $f'(1)$.
 \item $y=\sqrt[3]{x}\,\ln x$
 \item $y=(2x-1)\bigl(x^2+3\bigr)$
 \item $f(x)=x^3\,2^x\ln x$.\quad Find $f'(1)$.
\end{enumerate}}
\workingspace

% Unit 2 bank Q180, (x-1)^2 sin x, moved to the question bank (3.1 chain rule) (Lee).
% AB 1998 MC 8 = Unit 2 bank Q38 (h' = f g', so f is constant) moved to the question bank (6.8):
% its last step, f' = 0 for all x so f is constant, is an antiderivative fact.

% ---------------------------------------------------------------------
\Needspace{20\baselineskip}
\subsection{The quotient rule}
\begin{notesbox}{Notes}
\begin{itemize}
 \item Quotient rule:
 \[
 \frac{d}{dx}\left(\frac{f(x)}{g(x)}\right)=\frac{\dfrac{d}{dx}\bigl(f(x)\bigr)\cdot g(x)-\dfrac{d}{dx}\bigl(g(x)\bigr)\cdot f(x)}{\bigl(g(x)\bigr)^2}=\frac{f'(x)g(x)-g'(x)f(x)}{\bigl(g(x)\bigr)^2}.
 \]
 In words: differentiate the top and leave the bottom alone, minus differentiate the bottom and leave the top alone, all over the bottom squared. The top comes first.
 \item Simplify the numerator; leave the bottom as a square, not expanded.
 \item Quicker without the quotient rule:
 \begin{itemize}
  \item one term on the bottom ($x^2$, $\sqrt x$, $3x$): split and use the power rule, e.g.\ $\dfrac{x^2+3\sqrt x}{x}=x+3x^{-\frac12}$;
  \item a constant on the bottom: $\dfrac{x^2+\sin x}{5}=\dfrac15\bigl(x^2+\sin x\bigr)$.
 \end{itemize}
\end{itemize}
\steph{Traps.}
\begin{itemize}
 \item The order in the numerator matters: $g'(x)f(x)-f'(x)g(x)$ gives the negative of the answer.
 \item Square the bottom. And $\dfrac{f'(x)}{g'(x)}$ is not the derivative.
\end{itemize}
\end{notesbox}

\Needspace{0.45\textheight}
% The old part (a) (removed) had a second method (split to 2+7/(x-3), then the chain rule): it is in the question bank, 3.1 (Lee).
\exampletitle{ex:quotient}{The quotient rule}
Find the derivative of each function.
\begin{enumerate}[label=(\alph*)]
 \item $y=\dfrac{e^x}{x^2+1}$
 \item $y=\dfrac{\sin x}{1+\cos x}$
 \item $y=\dfrac{x^2+3\sqrt x}{x}$
 \item $f(x)=\dfrac{3^x}{x}$.\quad Find $f'(1)$.
 \item $f(x)=\dfrac{\ln x}{x+1}$.\quad Find $f'(1)$.
\end{enumerate}
\end{questionbox}
\begin{examplesolution}
\relax
\begin{enumerate}[label=(\alph*)]
 \item
 \begin{align*}
 \frac{dy}{dx}&=\frac{\dfrac{d}{dx}\bigl(e^x\bigr)\cdot\bigl(x^2+1\bigr)-\dfrac{d}{dx}\bigl(x^2+1\bigr)\cdot e^x}{\bigl(x^2+1\bigr)^2} &&\why{quotient rule}\\
 &=\frac{e^x\bigl(x^2+1\bigr)-2x\,e^x}{\bigl(x^2+1\bigr)^2}\\
 &=\frac{e^x\bigl(x^2-2x+1\bigr)}{\bigl(x^2+1\bigr)^2}=\frac{e^x(x-1)^2}{\bigl(x^2+1\bigr)^2}. &&\why{factor out $e^x$}
 \end{align*}
 \item
 \begin{align*}
 \frac{dy}{dx}&=\frac{\dfrac{d}{dx}(\sin x)\cdot(1+\cos x)-\dfrac{d}{dx}(1+\cos x)\cdot\sin x}{(1+\cos x)^2} &&\why{quotient rule}\\
 &=\frac{\cos x(1+\cos x)-(-\sin x)\sin x}{(1+\cos x)^2}\\
 &=\frac{\cos x+\cos^2x+\sin^2x}{(1+\cos x)^2}=\frac{\cos x+1}{(1+\cos x)^2} &&\why{$\sin^2x+\cos^2x=1$}\\
 &=\frac{1}{1+\cos x}.
 \end{align*}
 \item One term on the bottom: split instead.
 \begin{align*}
 y&=\frac{x^2}{x}+\frac{3x^{\frac12}}{x}=x+3x^{-\frac12} &&\why{split}\\
 \frac{dy}{dx}&=1-\frac32x^{-\frac32}.
 \end{align*}
 \item The bottom is one term, but $3^x$ is not a power of $x$, so splitting does not help. Use the quotient rule.
 \begin{align*}
 f'(x)&=\frac{\dfrac{d}{dx}\bigl(3^x\bigr)\cdot x-\dfrac{d}{dx}(x)\cdot3^x}{x^2} &&\why{quotient rule}\\
 &=\frac{3^x\ln3\cdot x-3^x}{x^2}=\frac{3^x(x\ln3-1)}{x^2}\\
 f'(1)&=\frac{3(\ln3-1)}{1}=3\ln3-3.
 \end{align*}
 \item
 \begin{align*}
 f'(x)&=\frac{\dfrac{d}{dx}(\ln x)\cdot(x+1)-\dfrac{d}{dx}(x+1)\cdot\ln x}{(x+1)^2} &&\why{quotient rule}\\
 &=\frac{\dfrac1x(x+1)-\ln x}{(x+1)^2}\\
 f'(1)&=\frac{1\cdot2-0}{2^2}=\frac12. &&\why{$\ln1=0$}
 \end{align*}
\end{enumerate}
\begin{takeaway}
Bottom-squared denominator, top-first numerator. Before using the rule, check whether splitting is quicker.
\end{takeaway}
\end{examplesolution}

\practiceheading
\question[long]{practice}{pr:quotient}{}{Find the derivative of each function.
\begin{enumerate}[label=(\alph*)]
 \item $y=\dfrac{x}{e^x}$
 \item $y=\dfrac{\cos x}{1-\sin x}$
 \item $y=\dfrac{4x^3-\sqrt x}{x^2}$
 \item $f(x)=\dfrac{2^x}{x}$.\quad Find $f'(2)$.
 \item $f(x)=\dfrac{\ln x}{2x-1}$.\quad Find $f'(1)$.
\end{enumerate}}
\workingspace

\question{practice}{pr:quotcubic}{}{If $f(x)=\dfrac{5-x}{x^3+2}$, then $f'(x)=$}
\twochoices{$\dfrac{-4x^3+15x^2-2}{\bigl(x^3+2\bigr)^2}$}{$\dfrac{-2x^3+15x^2+2}{\bigl(x^3+2\bigr)^2}$}{$\dfrac{2x^3-15x^2-2}{\bigl(x^3+2\bigr)^2}$}{$\dfrac{4x^3-15x^2+2}{\bigl(x^3+2\bigr)^2}$}
\workingspace

\question{practice}{pr:quotlinear}{}{If $f(x)=\dfrac{3x-2}{2x+3}$, then $f'(x)=$}
\fourchoices{$-\dfrac{13}{(2x+3)^2}$}{$\dfrac{5}{(2x+3)^2}$}{$\dfrac{13}{(2x+3)^2}$}{$\dfrac{12x+5}{(2x+3)^2}$}
\workingspace

\question{practice}{pr:quotfg}{}{If $f(x)=x^2+4x-2$, $g(x)=4x+2$ and $h(x)=\dfrac{f(x)}{g(x)}$, then $h'(1)=$}
\fourchoices{$-\dfrac23$}{$\dfrac23$}{$\dfrac43$}{$\dfrac32$}
\workingspace

% ---------------------------------------------------------------------
\Needspace{20\baselineskip}
\subsection{\texorpdfstring{Derivatives of $\tan x$, $\cot x$, $\sec x$ and $\csc x$}{Derivatives of tan x, cot x, sec x and csc x}}
\begin{notesbox}{Notes}
\begin{itemize}
 \item The derivatives:
 \[
 \begin{aligned}
 &\frac{d}{dx}(\tan x)=\sec^2x, &\qquad &\frac{d}{dx}(\sec x)=\sec x\tan x,\\[4pt]
 &\frac{d}{dx}(\cot x)=-\csc^2x, &\qquad &\frac{d}{dx}(\csc x)=-\csc x\cot x.
 \end{aligned}
 \]
 The ``co'' functions ($\cot$, $\csc$) get a minus sign, like $\cos$.
 \item Where they come from: write the function with $\sin$ and $\cos$, then use the quotient rule (\mbox{Example~\ref{ex:trigsix}}).
 \item Simplify first when you can: $\cos x\tan x=\sin x$,\quad $\sin x\sec x=\tan x$,\quad $\sin x\cot x=\cos x$.
 \item Values: $\sec x=\dfrac{1}{\cos x}$, $\csc x=\dfrac{1}{\sin x}$, $\cot x=\dfrac{1}{\tan x}$; e.g.\ $\sec\dfrac{\pi}{4}=\sqrt2$, $\csc\dfrac{\pi}{6}=2$.
\end{itemize}
\end{notesbox}

\Needspace{0.45\textheight}
\exampletitle{ex:trigsix}{The other trig functions}
\begin{enumerate}[label=(\alph*)]
 \item Use the quotient rule to show that $\dfrac{d}{dx}(\tan x)=\sec^2x$ and $\dfrac{d}{dx}(\sec x)=\sec x\tan x$.
 \item Find $\dfrac{dy}{dx}$ for $y=3\tan x-2\sec x+\csc x+4\cos x-6\sqrt x+5\ln x$.
 \item Find $\dfrac{dy}{dx}$ for $y=e^x\tan x$.
 \item Let $f(x)=\sin x\sec x$. Find $f'\Bigl(\dfrac{\pi}{3}\Bigr)$.
 \item Find $\dfrac{dy}{dx}$ for $y=2^x\cot x$.
 \item Let $f(x)=\sec x+\tan x$. Find $f'\Bigl(\dfrac{\pi}{4}\Bigr)$.
\end{enumerate}
\end{questionbox}
\begin{examplesolution}
\relax
\begin{enumerate}[label=(\alph*)]
 \item Write each with $\sin$ and $\cos$.
 \begin{align*}
 \dfrac{d}{dx}(\tan x)&=\dfrac{d}{dx}\left(\frac{\sin x}{\cos x}\right)=\frac{\dfrac{d}{dx}(\sin x)\cdot\cos x-\dfrac{d}{dx}(\cos x)\cdot\sin x}{\cos^2x} &&\why{quotient rule}\\
 &=\frac{\cos x\cos x-(-\sin x)\sin x}{\cos^2x}=\frac{\cos^2x+\sin^2x}{\cos^2x}\\
 &=\frac{1}{\cos^2x}=\sec^2x. &&\why{$\sin^2x+\cos^2x=1$}\\[4pt]
 \dfrac{d}{dx}(\sec x)&=\dfrac{d}{dx}\left(\frac{1}{\cos x}\right)=\frac{\dfrac{d}{dx}(1)\cdot\cos x-\dfrac{d}{dx}(\cos x)\cdot1}{\cos^2x} &&\why{quotient rule}\\
 &=\frac{0\cdot\cos x-(-\sin x)}{\cos^2x}=\frac{\sin x}{\cos^2x} &&\why{the top $1$ has derivative $0$}\\
 &=\frac{1}{\cos x}\cdot\frac{\sin x}{\cos x}=\sec x\tan x.
 \end{align*}
 \item One rule for each term.
 \[
 \frac{dy}{dx}=3\sec^2x-2\sec x\tan x-\csc x\cot x-4\sin x-\frac{3}{\sqrt x}+\frac5x.
 \]
 \item
 \begin{align*}
 \frac{dy}{dx}&=\frac{d}{dx}\bigl(e^x\bigr)\cdot\tan x+\frac{d}{dx}(\tan x)\cdot e^x &&\why{product rule}\\
 &=e^x\tan x+e^x\sec^2x=e^x\bigl(\tan x+\sec^2x\bigr).
 \end{align*}
 \item Simplify first.
 \begin{align*}
 f(x)&=\sin x\cdot\frac{1}{\cos x}=\tan x &&\why{$\sec x=\dfrac{1}{\cos x}$}\\
 f'(x)&=\sec^2x\\
 f'\Bigl(\frac{\pi}{3}\Bigr)&=2^2=4. &&\why{$\sec\dfrac{\pi}{3}=\dfrac{1}{\cos\frac{\pi}{3}}=2$}
 \end{align*}
 (The product rule gives $\cos x\cdot\sec x+\sec x\tan x\cdot\sin x=1+\tan^2x$, which is the same, $\sec^2x$, but takes longer.)
 \item
 \begin{align*}
 \frac{dy}{dx}&=\frac{d}{dx}\bigl(2^x\bigr)\cdot\cot x+\frac{d}{dx}(\cot x)\cdot2^x &&\why{product rule}\\
 &=2^x\ln2\,\cot x-2^x\csc^2x=2^x\bigl(\ln2\,\cot x-\csc^2x\bigr).
 \end{align*}
 \item
 \begin{align*}
 f'(x)&=\sec x\tan x+\sec^2x\\
 f'\Bigl(\frac{\pi}{4}\Bigr)&=\sqrt2\cdot1+\bigl(\sqrt2\bigr)^2=2+\sqrt2. &&\why{$\sec\dfrac{\pi}{4}=\sqrt2$, $\tan\dfrac{\pi}{4}=1$}
 \end{align*}
\end{enumerate}
\end{examplesolution}

\practiceheading
\question[long]{practice}{pr:trigsix}{}{(a)\enspace Use the quotient rule to show that $\dfrac{d}{dx}(\cot x)=-\csc^2x$ and $\dfrac{d}{dx}(\csc x)=-\csc x\cot x$.
\begin{enumerate}[label=(\alph*),start=2]
 \item Find $\dfrac{dy}{dx}$ for $y=4\cot x+3\csc x-\sec x+2\sin x+8\sqrt x-3\ln x$.
 \item Find $\dfrac{dy}{dx}$ for $y=e^x\sec x$.
 \item Let $f(x)=\cot x\sin x$. Find $f'\Bigl(\dfrac{\pi}{6}\Bigr)$.
 \item Find $\dfrac{dy}{dx}$ for $y=5^x\tan x$.
 \item Let $f(x)=\csc x-\cot x$. Find $f'\Bigl(\dfrac{\pi}{6}\Bigr)$.
\end{enumerate}}
\workingspace

\question{practice}{pr:tanminuscot}{}{If $y=\tan x-\cot x$, then $\dfrac{dy}{dx}=$}
\fourchoices{$\sec x\csc x$}{$\sec x+\csc x$}{$\sec^2x-\csc^2x$}{$\sec^2x+\csc^2x$}
\workingspace

\question[long]{practice}{pr:derivecot}{}{Which of the following correctly shows the derivation of $\dfrac{d}{dx}(\cot x)$?
\begin{enumerate}[label=(\Alph*),leftmargin=2.2em,itemsep=8pt]
 \item $\dfrac{d}{dx}(\cot x)=\dfrac{d}{dx}\left(\dfrac{1}{\tan x}\right)=\dfrac{1}{\frac{d}{dx}(\tan x)}=\dfrac{-1}{\sec^2x}$
 \item $\dfrac{d}{dx}(\cot x)=\dfrac{d}{dx}\left(\dfrac{1}{\tan x}\right)=\dfrac{1}{\frac{d}{dx}(\tan x)}=\dfrac{1}{\sec^2x}$
 \item $\dfrac{d}{dx}(\cot x)=\dfrac{d}{dx}\left(\dfrac{1}{\tan x}\right)=\dfrac{\frac{d}{dx}(1)\cdot\tan x-1\cdot\frac{d}{dx}(\tan x)}{\tan^2x}=\dfrac{0\cdot\tan x-\sec^2x}{\tan^2x}=-\dfrac{\sec^2x}{\tan^2x}$
 \item $\dfrac{d}{dx}(\cot x)=\dfrac{d}{dx}\left(\dfrac{1}{\tan x}\right)=\dfrac{\frac{d}{dx}(1)\cdot\tan x+1\cdot\frac{d}{dx}(\tan x)}{\tan^2x}=\dfrac{0\cdot\tan x+\sec^2x}{\tan^2x}=\dfrac{\sec^2x}{\tan^2x}$
\end{enumerate}}
\workingspace

\question{practice}{pr:costan}{}{$\dfrac{d}{dx}(\cos x\tan x)=$}
\fourchoices{$\sec x+\sin x\tan x$}{$\cos x$}{$-\sin x\sec^2x$}{$\sin x$}
\workingspace

% ---------------------------------------------------------------------
\Needspace{20\baselineskip}
\subsection{Rules with values from a table}
\begin{notesbox}{Notes}
\begin{itemize}
 \item When $f$ and $g$ are known only through a table or a few values, write the rule with letters first, then substitute $x=a$:
 \[
 (fg)'(a)=f'(a)g(a)+g'(a)f(a),\qquad \left(\frac fg\right)'(a)=\frac{f'(a)g(a)-g'(a)f(a)}{\bigl(g(a)\bigr)^2}.
 \]
 \item Terms with $x$ are differentiated as usual, e.g.
 \[
 \frac{d}{dx}\bigl(x^3f(x)\bigr)=\frac{d}{dx}\bigl(x^3\bigr)\cdot f(x)+\frac{d}{dx}\bigl(f(x)\bigr)\cdot x^3=3x^2f(x)+f'(x)\,x^3
 \]
 and $\dfrac{d}{dx}\bigl(x-g(x)\bigr)=1-g'(x)$.
 \item A product or quotient needs both the values ($f(a)$, $g(a)$) and the derivatives ($f'(a)$, $g'(a)$). A sum needs only the derivatives.
 \item Read which row to use: a table often has several values of $x$, and each part may use a different row.
\end{itemize}
\steph{Trap.}
\begin{itemize}
 \item Read the table header: the columns are not always in the order $f$, $f'$, $g$, $g'$.
\end{itemize}
\end{notesbox}

\Needspace{0.45\textheight}
\exampletitle{ex:tablerules}{Rules with values from a table}
The table gives values of the differentiable functions $f$ and $g$ and their derivatives at selected values of $x$.
\begin{center}
\renewcommand{\arraystretch}{1.3}
\begin{tabular}{|c|c|c|c|c|}\hline
$x$ & $f(x)$ & $f'(x)$ & $g(x)$ & $g'(x)$\\ \hline
$0$ & $3$ & $2$ & $-1$ & $4$\\ \hline
$1$ & $2$ & $-3$ & $4$ & $5$\\ \hline
\end{tabular}
\end{center}
\begin{enumerate}[label=(\alph*)]
 \item $h(x)=3f(x)-e^xg(x)\cos x+2^x$.\quad Find $h'(0)$.
 \item $h(x)=\dfrac{f(x)g(x)}{\sqrt x+\ln x}$.\quad Find $h'(1)$.
 \item $h(x)=\dfrac{g(x)+\sin x}{\tan x+f(x)}$.\quad Find $h'(0)$.
\end{enumerate}
\end{questionbox}
\begin{examplesolution}
\relax
\begin{enumerate}[label=(\alph*)]
 \item A sum; the middle term is a product of three factors. Use the row $x=0$.
 \begin{align*}
 h'(x)&=3f'(x)-\Bigl[\frac{d}{dx}\bigl(e^x\bigr)\cdot g(x)\cos x+\frac{d}{dx}\bigl(g(x)\bigr)\cdot e^x\cos x\\
 &\qquad\qquad+\frac{d}{dx}(\cos x)\cdot e^xg(x)\Bigr]+2^x\ln2\\
 &=3f'(x)-\Bigl[e^xg(x)\cos x+g'(x)e^x\cos x-\sin x\,e^xg(x)\Bigr]+2^x\ln2\\
 h'(0)&=3(2)-\bigl[(1)(-1)(1)+(4)(1)(1)-0\bigr]+\ln2 &&\why{$e^0=1$, $\cos0=1$, $\sin0=0$}\\
 &=6-3+\ln2=3+\ln2.
 \end{align*}
 \item A quotient whose top is a product. Use the row $x=1$.
 \begin{itemize}
  \item Top: $f(x)g(x)$. At $x=1$: $(2)(4)=8$. Its derivative is $f'(x)g(x)+g'(x)f(x)$; at $x=1$: $(-3)(4)+(5)(2)=-2$.
  \item Bottom: $\sqrt x+\ln x$. At $x=1$: $1+0=1$. Its derivative is $\dfrac{1}{2\sqrt x}+\dfrac1x$; at $x=1$: $\dfrac12+1=\dfrac32$.
 \end{itemize}
 \begin{align*}
 h'(1)&=\frac{\dfrac{d}{dx}(\text{top})\cdot\text{bottom}-\dfrac{d}{dx}(\text{bottom})\cdot\text{top}}{(\text{bottom})^2} &&\why{quotient rule}\\
 &=\frac{(-2)(1)-\frac32(8)}{1^2}=-2-12=-14.
 \end{align*}
 \item A quotient. Use the row $x=0$.
 \begin{itemize}
  \item Top: $g(x)+\sin x$. At $x=0$: $-1+0=-1$. Its derivative is $g'(x)+\cos x$; at $x=0$: $4+1=5$.
  \item Bottom: $\tan x+f(x)$. At $x=0$: $0+3=3$. Its derivative is $\sec^2x+f'(x)$; at $x=0$: $1+2=3$.
 \end{itemize}
 \begin{align*}
 h'(0)&=\frac{(5)(3)-(3)(-1)}{3^2}=\frac{15+3}{9}=2. &&\why{quotient rule}
 \end{align*}
\end{enumerate}
\begin{takeaway}
Find the value and the derivative of each piece first (top, bottom, each factor); then put them into the rule.
\end{takeaway}
\end{examplesolution}

\practiceheading
\question[long]{practice}{pr:tablerules}{}{The table gives values of the differentiable functions $f$ and $g$ and their derivatives at selected values of $x$.
\begin{center}
\renewcommand{\arraystretch}{1.3}
\begin{tabular}{|c|c|c|c|c|}\hline
$x$ & $f(x)$ & $f'(x)$ & $g(x)$ & $g'(x)$\\ \hline
$0$ & $-2$ & $3$ & $1$ & $-4$\\ \hline
$1$ & $-1$ & $4$ & $3$ & $-2$\\ \hline
\end{tabular}
\end{center}
\begin{enumerate}[label=(\alph*)]
 \item $h(x)=4g(x)-e^xf(x)\cos x-3^x$.\quad Find $h'(0)$.
 \item $h(x)=\dfrac{f(x)g(x)}{2\sqrt x+\ln x}$.\quad Find $h'(1)$.
 \item $h(x)=\dfrac{f(x)-\sin x}{\tan x+g(x)}$.\quad Find $h'(0)$.
\end{enumerate}}
\workingspace

% ---------------------------------------------------------------------
\Needspace{\ifshowspace 0.4\textheight\else 14\baselineskip\fi}
\subsection*{Exercise \thesection}\addcontentsline{toc}{subsection}{Exercise \thesection}
\question{exercise}{ex1:cubic}{}{If $f(x)=x^3-x^2+x-1$, then $f'(2)=$}
\fourchoices{$5$}{$7$}{$9$}{$10$}
\workingspace

\question{exercise}{ex1:dfthree}{}{If $f(x)=x^3+2x$, then $\dfrac{d}{dx}\bigl(f(2)\bigr)=$}
\fourchoices{$0$}{$12$}{$14$}{$16$}
\workingspace

\question{exercise}{ex1:expsin}{}{If $f(x)=e^x\sin x$, then $f'(x)=$}
\twochoices{$e^x\cos x$}{$-e^x\cos x$}{$e^x(\sin x+\cos x)$}{$e^x(\sin x-\cos x)$}
\workingspace

\question{exercise}{ex1:cosln}{}{If $f(x)=\dfrac{\cos x}{\ln x}$, then $f'(x)=$}
\twochoices{$\dfrac{x\sin x\ln x-\cos x}{x(\ln x)^2}$}{$\dfrac{-x\sin x\ln x-\cos x}{x(\ln x)^2}$}{$-x\sin x$}{$\dfrac{-x\sin x\ln x+\cos x}{x(\ln x)^2}$}
\workingspace

\question{exercise}{ex1:cot}{}{$\dfrac{d}{dx}(\cot x)=$}
\fourchoices{$-\tan x$}{$-\csc^2x$}{$\csc x$}{$\sec^2x$}
\workingspace

\question[long]{exercise}{ex1:big}{}{
\begin{enumerate}[label=(\alph*)]
 \item Find $\dfrac{dy}{dx}$ for $y=4\sqrt x-\dfrac{3}{x^2}+2\sin x-\tan x+5e^x-\ln x+2^x+e^2+\ln3$.
 \item Find $\dfrac{dy}{dx}$ for $y=x^3e^x\cos x$.
 \item Let $f(x)=\dfrac{e^x+\sin x}{x^2+\cos x}$. Find $f'(0)$.
 \item Let $f(x)=x\ln x-3^x$. Find $f'(1)$ and $\dfrac{d}{dx}\bigl(f(1)\bigr)$.
\end{enumerate}}
\workingspace

\question{exercise}{ex1:frq}{}{The table gives values of a function $f$ and its derivative $f'$ for selected values of $x$. Let $g$ be the function given by $g(x)=(2x-1)f(x)$. Find $g'(3)$.
\begin{center}
\renewcommand{\arraystretch}{1.3}
\begin{tabular}{|c|c|c|c|c|c|}\hline
$x$ & $1$ & $2$ & $3$ & $4$ & $5$\\ \hline
$f(x)$ & $4$ & $8$ & $7$ & $3$ & $-1$\\ \hline
$f'(x)$ & $5$ & $-2$ & $-1$ & $6$ & $2$\\ \hline
\end{tabular}
\end{center}}
\workingspace

\question{exercise}{ex1:threefns}{}{The table gives values of the differentiable functions $f$, $g$ and $h$ and their derivatives at $x=0$. If $M(x)=\dfrac{f(x)+g(x)}{h(x)}$, what is the value of $M'(0)$?
\begin{center}
\renewcommand{\arraystretch}{1.3}
\begin{tabular}{|c|c|c|c|c|c|c|}\hline
$x$ & $f(x)$ & $f'(x)$ & $g(x)$ & $g'(x)$ & $h(x)$ & $h'(x)$\\ \hline
$0$ & $-1$ & $-2$ & $3$ & $5$ & $2$ & $7$\\ \hline
\end{tabular}
\end{center}}
\fourchoices{$-2$}{$-\dfrac{8}{49}$}{$\dfrac37$}{$5$}
\workingspace

\question[long]{exercise}{ex1:bigtable}{}{The table gives values of the differentiable functions $f$ and $g$ and their derivatives at selected values of $x$.
\begin{center}
\renewcommand{\arraystretch}{1.3}
\begin{tabular}{|c|c|c|c|c|}\hline
$x$ & $f(x)$ & $f'(x)$ & $g(x)$ & $g'(x)$\\ \hline
$0$ & $2$ & $-1$ & $3$ & $1$\\ \hline
$1$ & $-1$ & $2$ & $4$ & $-3$\\ \hline
\end{tabular}
\end{center}
\begin{enumerate}[label=(\alph*)]
 \item $h(x)=\dfrac{f(x)g(x)}{e^x+x}+g(1)$.\quad Find $h'(0)$.
 \item $k(x)=3^xf(x)-\sqrt x\,g(x)$.\quad Find $k'(1)$.
 \item Find $\dfrac{d}{dx}\bigl(f(0)g(0)\bigr)$.
\end{enumerate}}
\workingspace

\question{exercise}{ex1:cuberoot}{}{Let $f$ be a differentiable function with $f(8)=2$ and $f'(8)=5$. If $g(x)=\dfrac{f(x)}{\sqrt[3]{x}}$, what is the value of $g'(8)$?}
\fourchoices{$\dfrac{59}{24}$}{$\dfrac52$}{$\dfrac{61}{24}$}{$60$}
\workingspace

\question[long]{exercise}{ex1:whichrule}{}{Find $\dfrac{dy}{dx}$. Use the quickest method each time.
\begin{enumerate}[label=(\alph*)]
 \item $y=\bigl(x^2+1\bigr)\bigl(x^2-1\bigr)$
 \item $y=\dfrac{x^3-4\sqrt x}{x}$
 \item $y=\dfrac{x}{x^2+1}$
 \item $y=\sin^2x$
 \item $y=\dfrac{5}{2x^3}$
\end{enumerate}}
\workingspace
\clearpage
\section{What the Derivative Means}\label{sec:meaning}
% =====================================================================
% REMOVED AP QUESTIONS (Lee, Nov 2026). These past-paper / AP Classroom bank questions were cut to keep
% the section short. They are NOT anywhere else: not in this worksheet and not in the question bank.
% Similar questions that stay in the worksheet serve the same purpose (shown after the arrow).
% - AB 2015 Intl MC 77 = Q3 (average rate of change from a table)  ->  P37 (BC 2024 Alt FRQ 4(a)).
% - Q181 (greatest average rate of change from a table), Q182, Q183, Q184 (average rate of change of an
%   expression)  ->  Example 11 and its practice (average rate of change), BC 2024 Alt FRQ 4(a).
% - Q85 (e^{2+h}), Q108 ((16+h)^{1/4} - 2, a hidden f(a)), Q107 (top in the reverse order),
%   BC 2015 MC 12 ((e^x - 1)/x), AB 2014 MC 20 ((ln(x+3) - ln 5)/(x - 2), a shifted inside)
%   ->  Example 8 and P28 (limits that are derivatives).
% - AB 2012 MC 18 = Q226 (ln(4+h)), Q225 (ln(e+h) - 1), Q106 ((5^x - 1)/x)  ->  P28.
% - AB 2019 MC 2 = Q162 and AB 1998 MC 18 = Q210 (tangent line of cos x, x + cos x),
%   AB 2003 MC 24 = Q207 and AB 2011 FRQ 3(a) (tangent line of a polynomial)  ->  Example 9(a), P32(a)-(b).
% - Q170 (x-intercept of the tangent to 1/x at x = 3)  ->  P33 (AB 1999 FRQ 6, intercepts of a tangent).
% - AB 2009 FRQ 5(a) and BC 2016 Practice FRQ 1(a) (derivative estimates from a table)
%   ->  Example 10(a)-(b), P40, the AB 2008 and AB 2011 FRQ questions.
% - Not chosen for Exercise 2: AB/BC 2022 MC 1 Part B (where f changes fastest, from a graph) and
%   BC 2012 MC 77 (calculator, tangent slopes)  ->  Example 10(d), the calculator section.
% =====================================================================
% =====================================================================
% CED 2.1-2.3 (and 2.4 tangent lines with the rules). Built from helpers A, B, C.
% ---------------------------------------------------------------------
\Needspace{20\baselineskip}
\subsection{The algebraic definition of the derivative}
\begin{notesbox}{Notes}
The definition of the derivative (first principles):
\begin{align*}
f'(a)&=\lim_{x\to a}\frac{f(x)-f(a)}{x-a}\\
f'(a)&=\lim_{h\to0}\frac{f(a+h)-f(a)}{h}\\
f'(x)&=\lim_{h\to0}\frac{f(x+h)-f(x)}{h}
\end{align*}
\end{notesbox}

\Needspace{0.45\textheight}
% (a)-(d) new.
\exampletitle{a:ex:algdef}{The definition of the derivative}
\begin{enumerate}[label=(\alph*)]
 \item Let $f(x)=x^3-4x$. Write $f'(x)$ as a limit. Then write $f'(2)$ as a limit in two ways, with $h\to0$ and with $x\to2$. Do not evaluate the limits.
 \item Using first principles, find $f'(4)$ for $f(x)=\sqrt x$.
 \item Use the definition to find $f'(x)$ for $f(x)=\dfrac1x$. Then find $f'(2)$.
\end{enumerate}
\end{questionbox}
\begin{examplesolution}
\relax
\begin{enumerate}[label=(\alph*)]
 \item Replace every $x$ by $x+h$. For $f'(2)$, note that $f(2)=8-8=0$.
 \begin{align*}
 f'(x)&=\lim_{h\to0}\frac{(x+h)^3-4(x+h)-\bigl(x^3-4x\bigr)}{h},\\
 f'(2)&=\lim_{h\to0}\frac{(2+h)^3-4(2+h)-0}{h}=\lim_{x\to2}\frac{x^3-4x-0}{x-2}.
 \end{align*}
 (Check with the rules: $f'(x)=3x^2-4$, so all three limits give $f'(2)=8$.)
 \item Use $f'(a)=\displaystyle\lim_{x\to a}\frac{f(x)-f(a)}{x-a}$ with $a=4$.
 \begin{align*}
 f'(4)&=\lim_{x\to4}\frac{\sqrt x-2}{x-4} &&\why{$f(4)=\sqrt4=2$}\\
 &=\lim_{x\to4}\frac{\sqrt x-2}{\bigl(\sqrt x-2\bigr)\bigl(\sqrt x+2\bigr)} &&\why{$x-4=\bigl(\sqrt x\bigr)^2-2^2$}\\
 &=\lim_{x\to4}\frac{1}{\sqrt x+2}=\frac14.
 \end{align*}
 (Check: $f'(x)=\dfrac{1}{2\sqrt x}$, $f'(4)=\dfrac14$.)
 \item Combine the fractions on the top.
 \begin{align*}
 f'(x)&=\lim_{h\to0}\frac{\dfrac{1}{x+h}-\dfrac1x}{h}=\lim_{h\to0}\frac{\dfrac{x-(x+h)}{x(x+h)}}{h} &&\why{common denominator}\\
 &=\lim_{h\to0}\frac{-h}{h\,x(x+h)}=\lim_{h\to0}\frac{-1}{x(x+h)} &&\why{cancel $h$}\\
 &=-\frac{1}{x^2}, &&\why{$h=0$}\\
 f'(2)&=-\frac14.
 \end{align*}
 (Check: $\dfrac{d}{dx}\bigl(x^{-1}\bigr)=-x^{-2}$.)
\end{enumerate}
\end{examplesolution}

\practiceheading
% (a)-(e) new.
\question[long]{practice}{a:pr:defvar}{}{%
\begin{enumerate}[label=(\alph*)]
 \item Let $f(x)=2x^3+x$. Write $f'(x)$ as a limit. Then write $f'(-1)$ as a limit in two ways, with $h\to0$ and with $x\to-1$. Do not evaluate the limits.
 \item Using first principles, find $f'(9)$ for $f(x)=\sqrt x$.
 \item Use the definition to find $f'(x)$ for $f(x)=\dfrac{1}{x+2}$. Then find $f'(2)$.
\end{enumerate}}
\workingspace

\question{practice}{a:pr:defexpr}{}{If $f(x)=3x^2+2x$, then $f'(x)=$}
\twochoices{$\displaystyle\lim_{h\to0}\frac{\bigl(3x^2+2x+h\bigr)-\bigl(3x^2+2x\bigr)}{h}$}{$\displaystyle\lim_{x\to0}\frac{\bigl(3x^2+2x+h\bigr)-\bigl(3x^2+2x\bigr)}{h}$}{$\displaystyle\lim_{h\to0}\frac{\bigl(3(x+h)^2+2(x+h)\bigr)-\bigl(3x^2+2x\bigr)}{h}$}{$\displaystyle\lim_{x\to0}\frac{\bigl(3(x+h)^2+2(x+h)\bigr)-\bigl(3x^2+2x\bigr)}{h}$}
\workingspace

% ---------------------------------------------------------------------
% Helper B: Tangent lines (classification 1.3), Slopes and tangent lines with
% the rules (2.4), Estimating derivatives from tables and graphs (1.4, 1.5,
% plus Q160 and AB 2026 FRQ 4(A)).

% ---------------------------------------------------------------------
\Needspace{20\baselineskip}
\subsection{Limits that are derivatives}
\begin{notesbox}{Notes}
\begin{itemize}
 \item A limit in the form of the definition is a derivative: find $f$ and $a$, then evaluate the derivative with the rules.
 \item The step does not have to be $h$. Anything that goes to $0$ works, as long as the top and the bottom use the same step:
 \[
 f'(a)=\lim_{h\to0}\frac{f(a+2h)-f(a)}{2h}=\lim_{h\to0}\frac{f\bigl(a+h^2\bigr)-f(a)}{h^2}=\lim_{h\to0}\frac{f(a-3h)-f(a)}{-3h}.
 \]
\end{itemize}
\end{notesbox}

\Needspace{0.45\textheight}
% Parts and sources: (a) AB 2015 MC 25 = BC 2015 MC 25 = Q224 (P), rewritten in the x -> a form (Lee); (b) AB 2018 MC 19 = BC 2016P MC 19 = Q52 (P);
% (c) new; (d) AB 2016 MC 25 = BC 2016 MC 25 = Q222 (P); (e) new (step h^2 over h); (f) new, the method of BC 2024 Alt FRQ 5(a);
% (g) BC 2008B MC 4 = Q47 (P) (the value at pi/4 added, the idea of Q126).
% Dropped on Lee's request (similar parts): Q85, Q108, Q107, BC 2015 MC 12, AB 2014 MC 20.
\exampletitle{a:ex:dqlimits}{Limits that are derivatives}
\begin{enumerate}[label=(\alph*)]
 \item Evaluate $\displaystyle\lim_{x\to\frac{\pi}{3}}\frac{\sin x-\sin\dfrac{\pi}{3}}{x-\dfrac{\pi}{3}}$
 \item Evaluate $\displaystyle\lim_{x\to e}\frac{\bigl(x^{20}-3x\bigr)-\bigl(e^{20}-3e\bigr)}{x-e}$
 \item Evaluate $\displaystyle\lim_{h\to0}\frac{\tan\Bigl(\dfrac{\pi}{4}+h\Bigr)-1}{h}$
 \item Evaluate $\displaystyle\lim_{h\to0}\frac{e^{-1-h}-e^{-1}}{h}$
 \item Evaluate $\displaystyle\lim_{h\to0}\frac{\cos\Bigl(\dfrac{\pi}{2}+h^2\Bigr)-\cos\Bigl(\dfrac{\pi}{2}\Bigr)}{h}$
 \item The differentiable function $f$ has the values in the table. Find $\displaystyle\lim_{x\to3}\frac{8-2f(x)}{3-x}$.
 \begin{center}
 \renewcommand{\arraystretch}{1.3}
 \begin{tabular}{|c|c|c|}\hline
 $x$ & $1$ & $3$\\ \hline
 $f(x)$ & $6$ & $4$\\ \hline
 $f'(x)$ & $5$ & $-2$\\ \hline
 \end{tabular}
 \end{center}
 \item For values of $h$ very close to $0$, which function best approximates $f(x)=\dfrac{\tan(x+h)-\tan x}{h}$? When $x=\dfrac{\pi}{4}$, what is the limit of $\dfrac{\tan(x+h)-\tan x}{h}$ as $h\to0$?
\end{enumerate}
\end{questionbox}
\begin{examplesolution}
\relax
\begin{enumerate}[label=(\alph*)]
 \item $f(x)=\sin x$ and $a=\dfrac{\pi}{3}$:
 \[
 \lim_{x\to\frac{\pi}{3}}\frac{\sin x-\sin\frac{\pi}{3}}{x-\frac{\pi}{3}}=f'\Bigl(\frac{\pi}{3}\Bigr)=\cos\frac{\pi}{3}=\frac12.
 \]
 The trap answer $\dfrac{\sqrt3}{2}$ is $f\bigl(\frac{\pi}{3}\bigr)$.
 \item $f(x)=x^{20}-3x$ and $a=e$; the subtracted bracket is $f(e)$.
 \[
 \lim_{x\to e}\frac{f(x)-f(e)}{x-e}=f'(e)=20e^{19}-3. \qquad\why{$f'(x)=20x^{19}-3$}
 \]
 The trap answer $e^{20}-3e$ is $f(e)$.
 \item The $1$ is a hidden $f(a)$: $1=\tan\dfrac{\pi}{4}$. So $f(x)=\tan x$ and $a=\dfrac{\pi}{4}$:
 \[
 \lim_{h\to0}\frac{\tan\bigl(\frac{\pi}{4}+h\bigr)-\tan\frac{\pi}{4}}{h}=f'\Bigl(\frac{\pi}{4}\Bigr)=\sec^2\frac{\pi}{4}=\bigl(\sqrt2\bigr)^2=2.
 \]
 \item The step is $-h$, so the bottom must be $-h$ too. Divide top and bottom by $-1$:
 \begin{align*}
 \lim_{h\to0}\frac{e^{-1-h}-e^{-1}}{h}&=-\lim_{h\to0}\frac{e^{-1+(-h)}-e^{-1}}{-h} &&\why{$f(x)=e^x$, $a=-1$, step $-h$}\\
 &=-f'(-1)=-e^{-1}=-\frac1e.
 \end{align*}
 The trap answer $\dfrac1e$ forgets the minus sign.
 \item The step on the top is $h^2$, but the bottom is only $h$. Multiply the top and the bottom by $h$ to make the bottom $h^2$:
 \begin{align*}
 \lim_{h\to0}\frac{\cos\bigl(\frac{\pi}{2}+h^2\bigr)-\cos\frac{\pi}{2}}{h}&=\lim_{h\to0}h\cdot\frac{\cos\bigl(\frac{\pi}{2}+h^2\bigr)-\cos\frac{\pi}{2}}{h^2} &&\why{$f(x)=\cos x$, $a=\dfrac{\pi}{2}$, step $h^2$}\\
 &=0\cdot f'\Bigl(\frac{\pi}{2}\Bigr)=0\cdot\Bigl(-\sin\frac{\pi}{2}\Bigr)=0.
 \end{align*}
 \item $f(3)=4$, so $8-2f(x)=-2\bigl(f(x)-f(3)\bigr)$ and $3-x=-(x-3)$.
 \begin{align*}
 \lim_{x\to3}\frac{8-2f(x)}{3-x}&=\lim_{x\to3}\frac{-2\bigl(f(x)-f(3)\bigr)}{-(x-3)}=2\lim_{x\to3}\frac{f(x)-f(3)}{x-3}\\
 &=2f'(3)=2(-2)=-4.
 \end{align*}
 (The values at $x=1$ are not needed.)
 \item As $h\to0$ the quotient tends to $\dfrac{d}{dx}(\tan x)=\sec^2x$, so for $h$ very close to $0$ it is approximately $\sec^2x$. At $x=\dfrac{\pi}{4}$ the limit is
 \[
 \sec^2\frac{\pi}{4}=\bigl(\sqrt2\bigr)^2=2.
 \]
\end{enumerate}
\end{examplesolution}

\practiceheading
% (a) BC 2022 MC 21 = Q116 (P), rewritten in the x -> a form; (b) AB 2017 MC 21 = Q36 (P); (c)-(f) new; (g) Q23 with the "small h" question added.
% Dropped on Lee's request (similar parts): AB 2012 MC 18 = Q226, Q225, Q106.
\question[long]{practice}{a:pr:derivlimits}{}{%
\begin{enumerate}[label=(\alph*)]
 \item Evaluate $\displaystyle\lim_{x\to\frac{\pi}{5}}\frac{\cos x-\cos\dfrac{\pi}{5}}{x-\dfrac{\pi}{5}}$
 \item If $f(x)=\ln x$, find $\displaystyle\lim_{x\to3}\frac{f(x)-f(3)}{x-3}$.
 \item Evaluate $\displaystyle\lim_{h\to0}\frac{\sqrt{4+h}-2}{h}$
 \item Evaluate $\displaystyle\lim_{h\to0}\frac{\ln(3-h)-\ln3}{h}$
 \item Evaluate $\displaystyle\lim_{h\to0}\frac{\sin\bigl(\pi+h^2\bigr)-\sin\pi}{h}$
 \item For values of $h$ very close to $0$, which function best approximates $\dfrac{\ln(x+h)-\ln x}{h}$? When $x=2e$, what is the limit of $\dfrac{\ln(x+h)-\ln x}{h}$ as $h\to0$?
\end{enumerate}}
\workingspace

% BC 2024 Alt FRQ 5(a): the stem's "twice differentiable" and "f, f', f'' are continuous" shortened (f'' is Unit 3).
\question{practice}{a:pr:dqtable}{}{The function $f$ is differentiable. Values of $f$ and $f'$ at selected values of $x$ are given in the table. Evaluate $\displaystyle\lim_{x\to4}\left(\frac{6-2f(x)}{4-x}\right)$ or determine that the limit does not exist. Justify your answer.
\begin{center}
\renewcommand{\arraystretch}{1.3}
\begin{tabular}{|c|c|c|c|c|}\hline
$x$ & $0$ & $1$ & $2$ & $4$\\ \hline
$f(x)$ & $11$ & $6$ & $-1$ & $3$\\ \hline
$f'(x)$ & $-2$ & $-7$ & $-6$ & $14$\\ \hline
\end{tabular}
\end{center}}
\workingspace

% Unit 2 bank Q109 and Q115 (new numbers).
\question{practice}{a:pr:evallim}{}{
\begin{enumerate}[label=(\alph*)]
 \item Evaluate $\displaystyle\lim_{h\to0}\frac{(x+h)^2+4(x+h)-x^2-4x}{h}$. Which derivative is it?
 \item Evaluate $\displaystyle\lim_{h\to0}\frac{(2+h)^3+2(2+h)-(2)^3-2(2)}{h}$.
\end{enumerate}}
\workingspace

\Needspace{22\baselineskip}
\question{practice}{a:pr:dqgraph}{}{Which of the following could be the graph of $y=\dfrac{\cos(x+0.001)-\cos(x)}{0.001}$?}
\par\noindent\begin{minipage}{\linewidth}
\begin{center}
\begin{minipage}[t]{0.45\linewidth}\centering
(A)\par
\begin{asy}
unitsize(0.6cm);
axes(-3.9,4,-1.6,1.7);
xlab(-pi,"$-\pi$"); xlab(pi,"$\pi$"); ynums(new real[]{-1,1});
curve(new real(real x){return cos(x);}, -pi, pi);
\end{asy}
\end{minipage}\hfill
\begin{minipage}[t]{0.45\linewidth}\centering
(B)\par
\begin{asy}
unitsize(0.6cm);
axes(-3.9,4,-1.6,1.7);
xlab(-pi,"$-\pi$"); xlab(pi,"$\pi$"); ynums(new real[]{-1,1});
curve(new real(real x){return -cos(x);}, -pi, pi);
\end{asy}
\end{minipage}\par\medskip
\begin{minipage}[t]{0.45\linewidth}\centering
(C)\par
\begin{asy}
unitsize(0.6cm);
axes(-3.9,4,-1.6,1.7);
xlab(-pi,"$-\pi$"); xlab(pi,"$\pi$"); ynums(new real[]{-1,1});
curve(new real(real x){return sin(x);}, -pi, pi);
\end{asy}
\end{minipage}\hfill
\begin{minipage}[t]{0.45\linewidth}\centering
(D)\par
\begin{asy}
unitsize(0.6cm);
axes(-3.9,4,-1.6,1.7);
xlab(-pi,"$-\pi$"); xlab(pi,"$\pi$"); ynums(new real[]{-1,1});
curve(new real(real x){return -sin(x);}, -pi, pi);
\end{asy}
\end{minipage}
\end{center}
\end{minipage}
\workingspace

% ---------------------------------------------------------------------
\Needspace{20\baselineskip}
\subsection{The graphical meaning of the derivative}
\begin{notesbox}{Notes}
\begin{center}
\renewcommand{\arraystretch}{1.6}
\begin{tabular}{>{\centering\arraybackslash}p{0.41\linewidth}|>{\centering\arraybackslash}p{0.1\linewidth}|>{\centering\arraybackslash}p{0.41\linewidth}}
\begin{asy}
unitsize(0.5cm);
real f(real x){return 0.5+0.12*x^2;}
real fp(real x){return 0.24*x;}
pen org=rgb(1,0.5,0)+linewidth(1.1pt);
real a=2.6, b=5.4;
pair A=(a,f(a)), B=(b,f(b));
axes(-0.4,6.8,-0.4,5.4);
curve(f,0,6.2);
pair d=unit(B-A);
draw((A-1.0*d)--(B+0.5*d), org);
closeddot(A); closeddot(B);
label("$(a,f(a))$",A,NW,fontsize(8pt));
label("$(x,f(x))$",B,SE,fontsize(8pt));
\end{asy}
& &
\begin{asy}
unitsize(0.5cm);
real f(real x){return 0.5+0.12*x^2;}
real fp(real x){return 0.24*x;}
pen org=rgb(1,0.5,0)+linewidth(1.1pt);
real a=2.6, b=3.3;
pair A=(a,f(a)), B=(b,f(b));
axes(-0.4,6.8,-0.4,5.4);
curve(f,0,6.2);
draw((a-2.2,f(a)-2.2*fp(a))--(a+3.2,f(a)+3.2*fp(a)), org);
closeddot(A); closeddot(B);
label("$(a,f(a))$",A,NW,fontsize(8pt));
label("$(x,f(x))$",B+(0.15,-0.1),SE,fontsize(8pt));
\end{asy}
\\ \cline{1-1}\cline{3-3}
\textcolor{brightorange}{\textbf{secant line}} & $\xrightarrow{\ x\to a\ }$ & \textcolor{brightorange}{\textbf{tangent line}}\\ \cline{1-1}\cline{3-3}
\begin{tabular}[t]{@{}c@{}}slope of the \textcolor{brightorange}{secant line}\\
$\displaystyle=\dfrac{f(x)-f(a)}{x-a}$
\end{tabular}
& &
\begin{tabular}[t]{@{}c@{}}slope of the \textcolor{brightorange}{tangent line}\\
$\displaystyle=\lim_{x\to a}\dfrac{f(x)-f(a)}{x-a}$\\
$=f'(a)$
\end{tabular}\\
\end{tabular}
\end{center}
\medskip
\begin{center}
\renewcommand{\arraystretch}{1.6}
\begin{tabular}{>{\centering\arraybackslash}p{0.41\linewidth}|>{\centering\arraybackslash}p{0.1\linewidth}|>{\centering\arraybackslash}p{0.41\linewidth}}
\begin{asy}
unitsize(0.5cm);
real f(real x){return 0.5+0.12*x^2;}
real fp(real x){return 0.24*x;}
pen org=rgb(1,0.5,0)+linewidth(1.1pt);
real a=2.6, b=5.4;
pair A=(a,f(a)), B=(b,f(b));
axes(-0.4,6.8,-0.4,5.4);
curve(f,0,6.2);
pair d=unit(B-A);
draw((A-1.0*d)--(B+0.5*d), org);
closeddot(A); closeddot(B);
label("$(a,f(a))$",A,NW,fontsize(8pt));
draw(A--(b,f(a)), gray(0.4)+dashed);
draw((b,f(a))--B, gray(0.6)+dotted);
label("$h$",((a+b)/2,f(a)),S,fontsize(8pt));
label("$(a+h,f(a+h))$",B,SE,fontsize(8pt));
\end{asy}
& &
\begin{asy}
unitsize(0.5cm);
real f(real x){return 0.5+0.12*x^2;}
real fp(real x){return 0.24*x;}
pen org=rgb(1,0.5,0)+linewidth(1.1pt);
real a=2.6, b=3.3;
pair A=(a,f(a)), B=(b,f(b));
axes(-0.4,6.8,-0.4,5.4);
curve(f,0,6.2);
draw((a-2.2,f(a)-2.2*fp(a))--(a+3.2,f(a)+3.2*fp(a)), org);
closeddot(A); closeddot(B);
label("$(a,f(a))$",A,NW,fontsize(8pt));
label("$(a+h,f(a+h))$",B+(0.15,-0.1),SE,fontsize(8pt));
\end{asy}
\\ \cline{1-1}\cline{3-3}
\textcolor{brightorange}{\textbf{secant line}} & $\xrightarrow{\ h\to0\ }$ & \textcolor{brightorange}{\textbf{tangent line}}\\ \cline{1-1}\cline{3-3}
\begin{tabular}[t]{@{}c@{}}slope of the \textcolor{brightorange}{secant line}\\
$\displaystyle=\dfrac{f(a+h)-f(a)}{h}$
\end{tabular}
& &
\begin{tabular}[t]{@{}c@{}}slope of the \textcolor{brightorange}{tangent line}\\
$\displaystyle=\lim_{h\to0}\dfrac{f(a+h)-f(a)}{h}$\\
$=f'(a)$
\end{tabular}\\
\end{tabular}
\end{center}
\begin{center}
\large $f'(a)=$ slope of the \textcolor{brightorange}{tangent line}
\end{center}
\end{notesbox}

\Needspace{0.45\textheight}
% (a) new (basic); (b) AB 1999 FRQ 4(a), 4(c); (c) Unit 2 bank Q164 (graph and data);
% (d) AB 2006 Form B FRQ 3(a), (b), (d), the skateboard ramp (moved here from Section 3; FRQ part (c) and "(iii)" need "increasing", Unit 5). The old (a), AB 2001 FRQ 4(c), is now P32(b).
% Moved to practice only on Lee's request: the DE slope (BC 2013 FRQ 6(a)), the backwards parts (Q166),
% intercepts in terms of w, "show tangent", parallel/perpendicular, unknown constants.
% Removed on Lee's request: the old (a) (new), (g) (Unit 2 bank Q162 = AB 2019 MC 2; AB 1998 MC 18 = Q210), (h) (new).
\exampletitle{b:ex:tangent}{Tangent lines}
\begin{enumerate}[label=(\alph*)]
 \item Let $h(x)=x^3-2x+1$. Write an equation of the line tangent to the graph of $h$ at $x=2$.
 \item $f(0)=2$ and $f'(0)=-3$. The function $g$ has derivative $g'(x)=e^{-2x}\bigl(3f(x)+2f'(x)\bigr)$, and $g(0)=4$. Write equations of the lines tangent to the graphs of $f$ and of $g$ at $x=0$.
 \item $f(1)=3$, and the graph of $f'$, the derivative of $f$, is shown. Write an equation of the line tangent to the graph of $f$ at $x=1$.
 \begin{center}
 \begin{asy}
 unitsize(1cm);
 grid(-3.3,3.3,-2.2,2.3,0.5);
 axes(-3.4,3.7,-2.3,2.7);
 xnums(new real[]{-3,-2,-1,1,2,3}); ynums(new real[]{-2,-1,2});
 curve(new real(real x){return 1.5-0.5*cos(pi*x);}, -3.3, 2);
 seg((2,1),(3.3,1));
 label("$1$",(0,1),2W,fontsize(8pt),Fill(white));
 label("$O$",(0,0),SW,fontsize(8pt));
 \end{asy}
\par\vspace{2pt}
 {\small Graph of $f'$}
 \end{center}
 \item \srctag{AB 2006 Form B FRQ}\enspace The graph below models the height of a skateboard ramp as a function of $x$. The function meets the following requirements.
 \begin{enumerate}[label=\arabic*.]
  \item At $x=0$, the value of the function is $0$, and the slope of the graph of the function is $0$.
  \item At $x=4$, the value of the function is $1$, and the slope of the graph of the function is $1$.
  \item Between $x=0$ and $x=4$, the function is increasing.
 \end{enumerate}
 \begin{center}
 \begin{asy}
 unitsize(1cm,2cm);
 axes(-0.3,5,-0.15,1.35,"$x$","height");
 xnums(new real[]{1,2,3,4}); ynums(new real[]{1});
 curve(new real(real x){return x^4/256;}, 0, 4.1);
 \end{asy}
 \end{center}
 \begin{enumerate}[label=(\roman*)]
  \item Let $f(x)=ax^2$, where $a$ is a nonzero constant. Show that it is not possible to find a value for $a$ so that $f$ meets requirement 2 above.
  \item Let $g(x)=cx^3-\dfrac{x^2}{16}$, where $c$ is a nonzero constant. Find the value of $c$ so that $g$ meets requirement 2 above. Show the work that leads to your answer.
  \item Let $h(x)=\dfrac{x^n}{k}$, where $k$ is a nonzero constant and $n$ is a positive integer. Find the values of $k$ and $n$ so that $h$ meets requirement 2 above. Show that $h$ also meets requirement 1 above.
 \end{enumerate}
\end{enumerate}
\end{questionbox}
\begin{examplesolution}
\relax
\begin{enumerate}[label=(\alph*)]
 \item The point and the slope at $x=2$:
 \begin{align*}
 h(2)&=8-4+1=5 &&\why{the point is $(2,5)$}\\
 h'(x)&=3x^2-2\imp h'(2)=12-2=10 &&\why{the slope}\\
 y-5&=10(x-2).
 \end{align*}
 \item The line tangent to $f$: point $(0,2)$, slope $-3$, so $y-2=-3x$.

 For $g$: substitute $x=0$, $f(0)=2$ and $f'(0)=-3$ into the formula for $g'$.
 \begin{align*}
 g'(0)&=e^0\bigl(3f(0)+2f'(0)\bigr)=1\cdot\bigl(3(2)+2(-3)\bigr)=0 &&\why{$e^0=1$}\\
 y-4&=0(x-0),\quad\text{i.e. } y=4. &&\why{a horizontal tangent}
 \end{align*}
 \item Read the height of the graph of $f'$ at $x=1$: $f'(1)=2$. The point is $\bigl(1,f(1)\bigr)=(1,3)$.
 \[
 y-3=2(x-1).
 \]
 The graph of $f'$ gives slopes, not values of $f$: the point $(1,2)$ is not on the graph of $f$.
 \item Requirement 2 is two equations: value $1$ and slope $1$ at $x=4$.
 \begin{enumerate}[label=(\roman*)]
  \item $f'(x)=2ax$.
  \begin{align*}
  f(4)&=16a=1\imp a=\frac1{16} &&\why{the value}\\
  f'(4)&=8a=1\imp a=\frac18 &&\why{the slope}
  \end{align*}
  The two equations need different values of $a$, so no value of $a$ meets requirement 2.
  \item $g'(x)=3cx^2-\dfrac{x}{8}$.
  \begin{align*}
  g(4)&=64c-1=1\imp c=\frac1{32} &&\why{the value}\\
  g'(4)&=48c-\frac12=1\imp c=\frac1{32} &&\why{the slope}
  \end{align*}
  Both equations give $c=\dfrac1{32}$.
  \item $h'(x)=\dfrac{nx^{n-1}}{k}$.
  \begin{align*}
  h(4)&=\frac{4^n}{k}=1\imp k=4^n &&\why{the value}\\
  h'(4)&=\frac{n\cdot4^{n-1}}{k}=\frac{n\cdot4^{n-1}}{4^n}=\frac n4=1\imp n=4 &&\why{the slope}
  \end{align*}
  So $n=4$ and $k=4^4=256$, i.e. $h(x)=\dfrac{x^4}{256}$.

  Requirement 1: $h(0)=0$ and $h'(0)=\dfrac{4\cdot0^3}{256}=0$.
 \end{enumerate}
\end{enumerate}
\begin{takeaway}
Every tangent line needs a point $\bigl(a,f(a)\bigr)$ and a slope $f'(a)$. ``Tangent'' means same point and same slope; ``parallel'' means same slope; ``perpendicular'' means slope $-\dfrac1m$.
\end{takeaway}
\end{examplesolution}

\practiceheading
% (a) new (basic); (b) AB 2001 FRQ 4(c), moved from Example 9; (d) AB 2003 FRQ 4(c) = BC 2003 FRQ 4(c); (e) AB 2005 Form B FRQ 6(b).
\question[long]{practice}{b:pr:tangent}{}{
\begin{enumerate}[label=(\alph*)]
 \item Let $g(x)=2x^3+x^2-5$. Write an equation of the line tangent to the graph of $g$ at $x=-1$.
 \item $h(4)=-3$ and $h'(x)=\dfrac{x^2-2}{x}$ for $x\ne0$. Write an equation of the line tangent to the graph of $h$ at $x=4$.
 \item $f(0)=-1$ and $f'(0)=3$. The function $g$ has derivative $g'(x)=\bigl(3f(x)+f'(x)\bigr)\cos x$, and $g(0)=5$. Write equations of the lines tangent to the graphs of $f$ and of $g$ at $x=0$.
 \item The function $f$ is defined on $-3\le x\le4$, and $f(0)=3$. The graph of $f'$, the derivative of $f$, consists of one line segment and a semicircle, as shown. Write an equation of the line tangent to the graph of $f$ at the point $(0,3)$.
\end{enumerate}}
\begin{center}
\begin{asy}
unitsize(0.9cm);
grid(-3,4,-2,2);
axes(-3.5,4.7,-2.5,2.5);
xnums(new real[]{-3,-1}); ynums(new real[]{-2,-1,1,2});
draw((-2,0),(0,2.5)--(0,-2.5)); label("$-2$",(-2,0),2SW,fontsize(8pt));
for (int i=1; i<=4; ++i) {draw((i,0),(0,2.5)--(0,-2.5)); label("$"+string(i)+"$",(i,0),2N,fontsize(8pt));}
seg((-3,1),(0,-2));
draw(arc((2,-2),2,180,0), linewidth(0.9pt));
closeddot((-3,1)); closeddot((-2,0)); closeddot((0,-2)); closeddot((2,0)); closeddot((4,-2));
label("$(-3,1)$",(-3,1),N,fontsize(8pt),Fill(white));
label("$(4,-2)$",(4,-2),S,fontsize(8pt),Fill(white));
\end{asy}
\par\vspace{2pt}
{\small Graph of $f'$}
\end{center}
\begin{enumerate}[label=(\alph*),start=5]
 \item The graph below models the height of a slide as a function of $x$. The function meets the following requirements.
 \begin{enumerate}[label=\arabic*.]
  \item At $x=0$, the value of the function is $0$, and the slope of the graph of the function is $0$.
  \item At $x=2$, the value of the function is $1$, and the slope of the graph of the function is $2$.
  \item Between $x=0$ and $x=2$, the function is increasing.
 \end{enumerate}
 \begin{center}
 \begin{asy}
 unitsize(1.4cm,1.4cm);
 axes(-0.3,2.8,-0.15,1.35,"$x$","height");
 xnums(new real[]{1,2}); ynums(new real[]{1});
 curve(new real(real x){return x^4/16;}, 0, 2.05);
 \end{asy}
 \end{center}
 \begin{enumerate}[label=(\roman*)]
  \item Let $f(x)=ax^2$, where $a$ is a nonzero constant. Show that it is not possible to find a value for $a$ so that $f$ meets requirement 2 above.
  \item Let $g(x)=cx^3-\dfrac{x^2}{4}$, where $c$ is a nonzero constant. Find the value of $c$ so that $g$ meets requirement 2 above.
  \item Let $h(x)=\dfrac{x^n}{k}$, where $k$ is a nonzero constant and $n$ is a positive integer. Find the values of $k$ and $n$ so that $h$ meets requirement 2 above. Show that $h$ also meets requirement 1 above.
 \end{enumerate}
 \item Let $y=f(x)$ be the solution of the differential equation $\dfrac{dy}{dx}=\dfrac{-xy^2}{2}$ with $f(-1)=2$. Write an equation of the line tangent to the graph of $f$ at $x=-1$.
 \item
 \begin{enumerate}[label=(\roman*)]
  \item The line tangent to the graph of $f$ at $x=-2$ is $y=-3x+1$. Find $f(-2)$ and $f'(-2)$. Must $f(0)$ be $1$?
  \item The line tangent to the graph of $f$ at the point $(-1,3)$ passes through the point $(3,11)$. Find $f'(-1)$, and write an equation of this tangent line.
 \end{enumerate}
\end{enumerate}
\workingspace

% (g) AB 1999 FRQ 6(a), 6(b); (h) AB 2003 Form B FRQ 1(a); (i)(ii) AB 2011 Form B FRQ 3(c).
% Removed on Lee's request (matching the example): AB 2003 MC 24 = Q207 and AB 2011 FRQ 3(a) (old (g)); the old (h).
\question[long]{practice}{b:pr:tanrules}{}{
\begin{enumerate}[label=(\alph*),start=8]
 \item The line $\ell$ is tangent to the graph of $y=\dfrac{1}{x^2}$ at the point $P\Bigl(w,\dfrac{1}{w^2}\Bigr)$, where $w>0$. The line $\ell$ crosses the $x$-axis at the point $R(k,0)$. Find $k$ in terms of $w$, and the value of $k$ when $w=3$.
 \item Let $f(x)=4x^2-x^3$. Show that the line $y=18-3x$ is tangent to the graph of $f$ at $x=3$.
 \item (i) Find the points on the graph of $y=x^3+3x^2-5$ at which the tangent line is parallel to the line $y=9x-1$.\\
 (ii) Let $f(x)=\sqrt x$ and $g(x)=6-x$. There is a point $P$ on the graph of $f$ at which the line tangent to the graph of $f$ is perpendicular to the graph of $g$. Find the coordinates of $P$.
 \item (i) The line $y=3x+k$ is tangent to the graph of $y=x^2-5x+1$. Find $k$.\\
 (ii) Let $g(x)=ax^2+3x$, where $a$ is a constant. The line tangent to the graph of $g$ at $x=1$ is parallel to the line through $(0,1)$ and $(2,-3)$. Find $a$.
\end{enumerate}}
\workingspace

\question{practice}{b:pr:twopoints}{}{If the line tangent to the graph of the function $f$ at the point $(1,7)$ passes through the point $(-2,-2)$, then $f'(1)$ is}
\fourchoices{$1$}{$3$}{$7$}{undefined}
\workingspace

\question{practice}{b:pr:drawntangent}{}{Shown is the graph of the differentiable function $f$ along with the line tangent to the graph of $f$ at $x=3$. What is the value of $f'(3)$?}
\begin{center}
\begin{asy}
unitsize(0.62cm);
grid(0,6.4,0,6.4);
axes(-0.4,6.9,-0.4,6.9);
xnums(new real[]{1,2,3,4,5,6}); ynums(new real[]{1,2,3,4,5,6});
seg((0,5.5),(6.4,2.3));
curve(new real(real x){return -0.0036*x^4-0.076378*x^3+0.055474*x^2+1.618184*x+1;}, 0, 4.732);
label("$y=f(x)$",(4.95,1.45),E,fontsize(9pt),Fill(white));
draw((4.95,1.45)--(4.62,1.75),linewidth(0.4pt));
label("$O$",(0,0),SW,fontsize(8pt));
\end{asy}
\end{center}
\fourchoices{$-\dfrac12$}{$-2$}{$4$}{$\dfrac{11}{2}$}
\workingspace

\question{practice}{b:pr:tanlinegraph}{}{Let $f$ be a differentiable function. The figure shows the graph of the line tangent to the graph of $f$ at $x=0$. Of the following, which must be true?}
\begin{center}
\begin{asy}
unitsize(0.7cm);
axes(-2.7,2.8,-2.7,2.8);
for (real t: new real[]{-2,-1,1,2}) {draw((t,0),(0,2.5)--(0,-2.5)); draw((0,t),(2.5,0)--(-2.5,0));}
label("$1$",(1,0),2S,fontsize(8pt)); label("$1$",(0,1),2W,fontsize(8pt));
label("$O$",(0,0),SW,fontsize(8pt));
seg((-2.32,2.64),(0.30,-2.6));
\end{asy}
\end{center}
\fourchoices{$f'(0)=-f(0)$}{$f'(0)<f(0)$}{$f'(0)=f(0)$}{$f'(0)>f(0)$}
\workingspace

\question{practice}{b:pr:limitslope}{}{The function $f$ satisfies the properties shown. What is the slope of the line tangent to the graph of $f$ at $x=-1$?
\[
f(0)=0\qquad\quad \lim_{x\to-1}f(x)=\frac25\qquad\quad \lim_{x\to-1}\frac{f(x)-f(-1)}{x+1}=-\frac35
\]}
\fourchoices{$-\dfrac35$}{$-\dfrac25$}{$\dfrac25$}{$1$}
\workingspace

\question{practice}{b:pr:tangentk}{}{In the $xy$-plane, the line $x+y=k$, where $k$ is a constant, is tangent to the graph of $y=x^2+3x+1$. What is the value of $k$?}
\fourchoices{$-3$}{$-2$}{$-1$}{$1$}
\workingspace

\question{practice}{b:pr:parallelb}{}{The function $g$ is defined by $g(x)=x^2+bx$, where $b$ is a constant. If the line tangent to the graph of $g$ at $x=-1$ is parallel to the line that contains the points $(0,-2)$ and $(3,4)$, what is the value of $b$?}
\fourchoices{$-1$}{$2$}{$\dfrac52$}{$4$}
\workingspace

\Needspace{0.45\textheight}
% (a) AB 2012 FRQ 1(a) = BC 2012 FRQ 1(a), estimate only (the interpretation is banked, 4.1);
% (b) AB/BC 2022 FRQ 1(a) (moved here from Example 12);
% (c)(i) AB 2015 FRQ 1(a) = BC 2015 FRQ 1(a), R'(45) only (the interpretation is banked, 4.1);
% (d) graph of BC 2014 MC 76 = Unit 2 bank Q12, with u moved to x = 4.10 (where the
% tangent is parallel to the secant; the source drew it at 4.17) and a grid added; (d)(iii) is Q12.
\exampletitle{b:ex:estimate}{Derivatives from tables and graphs}
\begin{enumerate}[label=(\alph*)]
\item The temperature of water in a tub is modeled by a differentiable function $W$, where $W(t)$ is measured in degrees Fahrenheit and $t$ in minutes. Selected values are given in the table.
\begin{center}
\renewcommand{\arraystretch}{1.3}
\begin{tabular}{|c|c|c|c|c|c|}\hline
$t$ (minutes) & $0$ & $4$ & $9$ & $15$ & $20$\\ \hline
$W(t)$ (degrees Fahrenheit) & $55.0$ & $57.1$ & $61.8$ & $67.9$ & $71.0$\\ \hline
\end{tabular}
\end{center}
\begin{enumerate}[label=(\roman*)]
 \item Use the data in the table to estimate $W'(12)$. Indicate units of measure.
 \item Estimate $W'(4)$.
\end{enumerate}
\item The radius of the circular cross section of a juice bottle at height $h$ above its base is given by a differentiable function $r$, where $h$ and $r(h)$ are measured in centimeters. Selected values are given in the table. Approximate $r'(3.1)$ using the average rate of change of $r$ over $0\le h\le6.2$. Indicate units of measure.
 \begin{center}
 \renewcommand{\arraystretch}{1.3}
 \begin{tabular}{|c|c|c|c|c|}\hline
 $h$ (cm) & $0$ & $6.2$ & $13$ & $20$\\ \hline
 $r(h)$ (cm) & $2.5$ & $2.7$ & $4.3$ & $2.6$\\ \hline
 \end{tabular}
 \end{center}
\item The piecewise-linear graph shows the rate $R(t)$, in liters per minute, at which water is pumped into a tank.
\begin{center}
\begin{asy}
unitsize(0.12cm,0.1cm);
for (real x=10; x<=55; x+=10) draw((x,0)--(x,32), gray(0.8)+linewidth(0.4pt));
for (real y=10; y<=30; y+=10) draw((0,y)--(57,y), gray(0.8)+linewidth(0.4pt));
axes(0,62,0,37,"$t$","$R(t)$");
xnums(new real[]{10,20,30,40,50}); ynums(new real[]{10,20,30});
seg((0,10),(20,30)); seg((20,30),(35,30)); seg((35,30),(55,0));
closeddot((0,10)); closeddot((20,30)); closeddot((35,30)); closeddot((55,0));
label("$(0,10)$",(0,10),SE,fontsize(8pt),Fill(white));
label("$(20,30)$",(20,30),N,fontsize(8pt));
label("$(35,30)$",(35,30),N,fontsize(8pt));
label("$(55,0)$",(55,0),NE,fontsize(8pt));
label("Time (minutes)",(28,-6),S,fontsize(9pt));
\end{asy}
\end{center}
\begin{enumerate}[label=(\roman*)]
 \item Find $R'(45)$ and $R'(10)$. Indicate units of measure.
 \item Let $k(t)=t^2R(t)$. Find $k'(10)$.
\end{enumerate}
\item The graph of the differentiable function $f$ is shown for $1\le x\le5$.
\begin{center}
\begin{asy}
unitsize(1.05cm);
real f(real x){real s=x-3; return 2.27374*s-0.494575*s^3+0.0704253*s^5-0.0082797*s^7;}
grid(0,5.5,-2.5,2.5,0.5);
axes(-0.3,5.9,-2.7,2.8);
ynums(new real[]{-2,-1,1,2});
curve(f,1,5);
real[] xs={1,1.4755,3,4.10,4.815,5}; string[] ls={"$1$","$r$","$t$","$u$","$w$","$5$"};
pair[] dirs={2N,2N,2NW,2S,2SW,2SE};
for (int i=0; i<6; ++i) {
  draw((xs[i],0),(0,2.5)--(0,-2.5));
  label(ls[i],(xs[i],0),dirs[i],fontsize(9pt),Fill(white));
  closeddot((xs[i],f(xs[i])));
}
label("$O$",(0,0),SW,fontsize(8pt));
\end{asy}
\par\vspace{2pt}
{\small Graph of $f$}
\end{center}
\begin{enumerate}[label=(\roman*)]
 \item Put $f'(r)$, $f'(t)$ and $f'(w)$ in order, from least to greatest.
 \item Which is greatest: $f(t)$, $f'(t)$, $f'(w)$ or $\dfrac{f(5)-f(1)}{5-1}$?
 \item Which of $r$, $t$, $u$, $w$ could be the $x$-coordinate of a point at which the line tangent to the graph of $f$ is parallel to the secant line through $\bigl(1,f(1)\bigr)$ and $\bigl(5,f(5)\bigr)$?
 \item How many such points are there on $1<x<5$?
\end{enumerate}
\end{enumerate}
\end{questionbox}
\begin{examplesolution}
\relax
\begin{enumerate}[label=(\alph*)]
 \item
 \begin{enumerate}[label=(\roman*)]
  \item $W'(12)$ is the slope of the tangent line at $t=12$. The closest table values on each side of $12$ are $t=9$ and $t=15$; the secant line through these two points has almost the same slope as the tangent line.
  \begin{center}
\begin{asy}
unitsize(0.36cm,0.2cm);
real[] T={0,4,9,15,20}; real[] V={55.0,57.1,61.8,67.9,71.0};
guide g; for (int i=0; i<5; ++i) g=g..(T[i],V[i]-50);
pen org=rgb(1,0.5,0)+linewidth(1.1pt);
draw((-0.8,0)--(22,0), EndArrow(4)); draw((0,-0.8)--(0,24), EndArrow(4));
label("$t$",(22,0),E); label("$W(t)$",(0,24),N);
draw((-0.4,1.0)--(0.4,1.8)); draw((-0.4,2.0)--(0.4,2.8));
for (real v: new real[]{55,60,65,70}) {draw((0,v-50),(2.5,0)--(-2.5,0)); label("$"+string(v)+"$",(0,v-50),2W,fontsize(8pt));}
draw(g, linewidth(0.9pt));
for (int i=0; i<5; ++i) {closeddot((T[i],V[i]-50)); draw((T[i],0),(0,2.5)--(0,-2.5)); label("$"+string(T[i])+"$",(T[i],0),2S,fontsize(8pt));}
void tangat(real x, real h){real[] tt=times(g,x); pair P=point(g,tt[0]); pair d=dir(g,tt[0]); real m=d.y/d.x;
  draw((x-h,P.y-h*m)--(x+h,P.y+h*m), org); dot(P, linewidth(3pt));}
draw((9,11.8)--(15,17.9), gray(0.35)+dashed+linewidth(0.8pt));
tangat(12,3);
draw((12,0)--point(g,times(g,12)[0]), gray(0.6)+dotted); draw((12,0),(0,2.5)--(0,-2.5)); label("$12$",(12,0),2S,fontsize(8pt));
label("secant through $t=9$ and $t=15$",(17,6),S,gray(0.3)+fontsize(8pt)); draw((17,6.3)--(13.4,15.1),gray(0.5)+linewidth(0.4pt),EndArrow(3));
label("tangent at $t=12$",(4,19),N,rgb(0.85,0.4,0)+fontsize(8pt)); draw((4.8,19)--(10.2,13.6),rgb(0.85,0.4,0)+linewidth(0.4pt),EndArrow(3));
\end{asy}
  \par\small The table gives only the dots; the curve is one possible graph of $W$.
  \end{center}
  \[
  W'(12)\approx\frac{W(15)-W(9)}{15-9}=\frac{67.9-61.8}{6}=\frac{6.1}{6}\approx1.017\text{ degrees Fahrenheit per minute.}
  \]
  \item $t=4$ is in the table; use its neighbours $0$ and $9$. The secant line through them has almost the same slope as the tangent line at $t=4$.
  \begin{center}
\begin{asy}
unitsize(0.36cm,0.2cm);
real[] T={0,4,9,15,20}; real[] V={55.0,57.1,61.8,67.9,71.0};
guide g; for (int i=0; i<5; ++i) g=g..(T[i],V[i]-50);
pen org=rgb(1,0.5,0)+linewidth(1.1pt);
draw((-0.8,0)--(22,0), EndArrow(4)); draw((0,-0.8)--(0,24), EndArrow(4));
label("$t$",(22,0),E); label("$W(t)$",(0,24),N);
draw((-0.4,1.0)--(0.4,1.8)); draw((-0.4,2.0)--(0.4,2.8));
for (real v: new real[]{55,60,65,70}) {draw((0,v-50),(2.5,0)--(-2.5,0)); label("$"+string(v)+"$",(0,v-50),2W,fontsize(8pt));}
draw(g, linewidth(0.9pt));
for (int i=0; i<5; ++i) {closeddot((T[i],V[i]-50)); draw((T[i],0),(0,2.5)--(0,-2.5)); label("$"+string(T[i])+"$",(T[i],0),2S,fontsize(8pt));}
void tangat(real x, real h){real[] tt=times(g,x); pair P=point(g,tt[0]); pair d=dir(g,tt[0]); real m=d.y/d.x;
  draw((x-h,P.y-h*m)--(x+h,P.y+h*m), org); dot(P, linewidth(3pt));}
draw((0,5)--(9,11.8), gray(0.35)+dashed+linewidth(0.8pt));
tangat(4,3);
label("secant through $t=0$ and $t=9$",(14,4),S,gray(0.3)+fontsize(8pt)); draw((12,4.3)--(6.6,9.9),gray(0.5)+linewidth(0.4pt),EndArrow(3));
label("tangent at $t=4$",(3,16),N,rgb(0.85,0.4,0)+fontsize(8pt)); draw((3.2,16)--(5.6,8.5),rgb(0.85,0.4,0)+linewidth(0.4pt),EndArrow(3));
\end{asy}
  \end{center}
  \[
  W'(4)\approx\frac{W(9)-W(0)}{9-0}=\frac{61.8-55.0}{9}=\frac{6.8}{9}\approx0.756\text{ degrees Fahrenheit per minute.}
  \]
 \end{enumerate}
 \item $3.1$ is the midpoint of $[0,6.2]$.
 \[
 r'(3.1)\approx\frac{r(6.2)-r(0)}{6.2-0}=\frac{2.7-2.5}{6.2}=\frac{0.2}{6.2}=\frac{1}{31}\approx0.032\text{ centimeter per centimeter}.
 \]
 The two table points closest to $h=3.1$ on either side are $h=0$ and $h=6.2$. The secant line through them has almost the same slope as the tangent line at $h=3.1$:
 \begin{center}
\begin{asy}
unitsize(0.75cm);
real Y(real r){return (r-2.4)*16;}
real r(real h){return 2.5-0.03*h+0.0100417*h^2;}
real rp(real h){return -0.03+0.0200834*h;}
draw((-0.4,0)--(8,0), EndArrow(4)); draw((0,-0.3)--(0,Y(2.82)), EndArrow(4));
draw((-0.15,0.35)--(0.15,0.5)); draw((-0.15,0.5)--(0.15,0.65));
label("$h$",(8,0),E); label("$r(h)$",(0,Y(2.82)),N);
for (real h: new real[]{3.1,6.2}) {draw((h,0),(0,2.5)--(0,-2.5));}
label("$3.1$",(3.1,0),2S,fontsize(8pt)); label("$6.2$",(6.2,0),2S,fontsize(8pt));
for (real v: new real[]{2.5,2.6,2.7}) {draw((0,Y(v)),(2.5,0)--(-2.5,0)); label("$"+string(v)+"$",(0,Y(v)),2W,fontsize(8pt));}
guide g; for (int i=0; i<=100; ++i) {real h=7.4*i/100; g=g..(h,Y(r(h)));} draw(g, linewidth(0.9pt));
pair A=(0,Y(2.5)), B=(6.2,Y(2.7)); pair d=unit(B-A);
draw((A-0.3*d)--(B+1.0*d), gray(0.4)+dashed+linewidth(0.7pt));
pair T=(3.1,Y(r(3.1))); real m=16*rp(3.1);
draw((3.1-2.2,T.y-2.2*m)--(3.1+2.2,T.y+2.2*m), linewidth(1.2pt));
draw((3.1,0)--T, gray(0.6)+dotted); draw((6.2,0)--B, gray(0.6)+dotted);
closeddot(A); closeddot(B); dot(T, linewidth(3pt));
label("secant line",(1.0,A.y+1.0*(B.y-A.y)/6.2),SE,gray(0.3)+fontsize(8pt));
label("tangent line at $h=3.1$",(3.1+2.2,T.y+2.2*m),SE,fontsize(8pt));
\end{asy}
\par\small The table gives only the dots; the curve is one possible graph of $r$.\\ The two lines have almost the same slope, so the secant slope estimates $r'(3.1)$.
 \end{center}
 \item
 \begin{enumerate}[label=(\roman*)]
  \item $t=45$ is on the segment from $(35,30)$ to $(55,0)$; $t=10$ is on the segment from $(0,10)$ to $(20,30)$.
  \begin{align*}
  R'(45)&=\frac{0-30}{55-35}=-\frac{30}{20}=-1.5\text{ liters per minute per minute},\\
  R'(10)&=\frac{30-10}{20-0}=1\text{ liter per minute per minute}.
  \end{align*}
  Use the segment that contains the point: the slope of the middle segment, $0$, is not $R'(45)$.
  \item Product rule, with values read from the graph of $R$: $R(10)=20$ (halfway up the first segment) and $R'(10)=1$ from (i).
  \begin{align*}
  k'(t)&=\frac{d}{dt}\bigl(t^2\bigr)\cdot R(t)+\frac{d}{dt}\bigl(R(t)\bigr)\cdot t^2=2tR(t)+R'(t)\,t^2 &&\why{product rule}\\
  k'(10)&=2(10)(20)+(1)(100)=500.
  \end{align*}
 \end{enumerate}
 \item
 \begin{enumerate}[label=(\roman*)]
  \item At $r$ the graph has a valley, so $f'(r)=0$. At $t$ the graph rises steeply, so $f'(t)>0$. At $w$ (after the peak) the graph falls, so $f'(w)<0$.
  \[
  f'(w)<f'(r)<f'(t).
  \]
  \item Read each number.
  \begin{itemize}
   \item $f(t)=0$: the graph crosses the $x$-axis at $t$.
   \item $f'(w)<0$: the graph falls at $w$.
   \item $\dfrac{f(5)-f(1)}{5-1}\approx\dfrac{1.8-(-1.8)}{4}=0.9$: the slope of the secant.
   \item $f'(t)\approx2$: near $t$ the graph climbs about $2$ units from $x=2.5$ to $x=3.5$, much steeper than the secant.
  \end{itemize}
  The greatest is $f'(t)$.
  \item Draw the secant line (dashed), then the tangent line at each labelled point. Only the tangent at $u$ runs parallel to the secant: the tangent at $r$ is flat, the one at $t$ is much steeper, and the one at $w$ falls. The answer is $u$.
  \begin{center}
\begin{asy}
unitsize(1.15cm);
real f(real x){real s=x-3; return 2.27374*s-0.494575*s^3+0.0704253*s^5-0.0082797*s^7;}
real fp(real x){real s=x-3; return 2.27374-3*0.494575*s^2+5*0.0704253*s^4-7*0.0082797*s^6;}
pen org=rgb(1,0.5,0)+linewidth(1.1pt);
grid(0,5.5,-2.5,2.5,0.5);
axes(-0.3,5.9,-2.7,2.8);
ynums(new real[]{-2,-1,1,2});
curve(f,1,5);
pair P1=(1,f(1)), P5=(5,f(5));
draw(P1--P5, gray(0.35)+dashed+linewidth(0.8pt));
closeddot(P1); closeddot(P5);
void tang(real a, real h, pen p){draw((a-h,f(a)-h*fp(a))--(a+h,f(a)+h*fp(a)), p); closeddot((a,f(a)));}
tang(1.4755,0.4,org); tang(3,0.35,org); tang(4.10,0.32,org); tang(4.815,0.18,org);
real[] xs={1.4755,3,4.10,4.815}; string[] ls={"$r$","$t$","$u$","$w$"};
pair[] dirs={2N,2NW,2S,2SE};
for (int i=0; i<4; ++i) {draw((xs[i],0),(0,2.5)--(0,-2.5)); label(ls[i],(xs[i],0),dirs[i],fontsize(9pt),Fill(white));}
label("secant",(2.4,f(1)+1.4*0.8923),SE,gray(0.3)+fontsize(8pt),Fill(white));
\end{asy}
  \end{center}
  \item Slide a line parallel to the secant along the curve and mark where it just touches the curve as a tangent. That happens twice: at $u$, and at one point between $r$ and $t$. There are $2$ such points.
  \begin{center}
\begin{asy}
unitsize(1.15cm);
real f(real x){real s=x-3; return 2.27374*s-0.494575*s^3+0.0704253*s^5-0.0082797*s^7;}
real fp(real x){real s=x-3; return 2.27374-3*0.494575*s^2+5*0.0704253*s^4-7*0.0082797*s^6;}
pen org=rgb(1,0.5,0)+linewidth(1.1pt);
grid(0,5.5,-2.5,2.5,0.5);
axes(-0.3,5.9,-2.7,2.8);
ynums(new real[]{-2,-1,1,2});
curve(f,1,5);
pair P1=(1,f(1)), P5=(5,f(5));
draw(P1--P5, gray(0.35)+dashed+linewidth(0.8pt));
closeddot(P1); closeddot(P5);
void tang(real a, real h, pen p){draw((a-h,f(a)-h*fp(a))--(a+h,f(a)+h*fp(a)), p); closeddot((a,f(a)));}
tang(1.9005,0.8,org); tang(4.0995,0.8,org);
real[] xs={1.4755,3,4.10}; string[] ls={"$r$","$t$","$u$"};
pair[] dirs={2N,2NW,2S};
for (int i=0; i<3; ++i) {draw((xs[i],0),(0,2.5)--(0,-2.5)); label(ls[i],(xs[i],0),dirs[i],fontsize(9pt),Fill(white));}
\end{asy}
  \par\small The two tangent lines (orange) are parallel to the secant (dashed).
  \end{center}
 \end{enumerate}
\end{enumerate}
\begin{takeaway}
From a table: the nearest values on each side, divided by their real gap. From a graph: the slope of the segment or tangent at the point, never the height.
\end{takeaway}
\end{examplesolution}

\practiceheading
% (a) AB 2018 FRQ 4(a), estimate only (the interpretation is banked, 4.1);
% (a)(i) AB 2006 Form B FRQ 4(a) (f'(22)), with f'(14) and f'(2) added; (b) graph of BC 2008 Form B MC 89 = Unit 2 bank Q215 (grid added).
% The made-up table parts (tree, coffee) were removed: the real AB 2008 and AB 2011 FRQ table questions follow (Lee).
\question[long]{practice}{b:pr:estimate}{}{}
\begin{enumerate}[label=(\alph*)]
\item The rate, in calories per minute, at which a person using an exercise machine burns calories is modeled by the function $f$. In the figure, $f(t)=-\dfrac14t^3+\dfrac32t^2+1$ for $0\le t\le4$, and $f$ is piecewise linear for $4\le t\le24$.
\begin{center}
\begin{asy}
unitsize(0.26cm,0.26cm);
for (real x=4; x<=24; x+=4) draw((x,0)--(x,17), gray(0.8)+linewidth(0.4pt));
for (real y=3; y<=15; y+=3) draw((0,y)--(26,y), gray(0.8)+linewidth(0.4pt));
axes(0,28,0,19,"$t$","");
xnums(new real[]{4,8,12,16,20,24}); ynums(new real[]{3,6,9,12,15});
curve(new real(real t){return -t^3/4+1.5*t^2+1;},0,4);
seg((4,9),(12,9)); seg((12,9),(16,15)); seg((16,15),(20,15)); seg((20,15),(24,3));
closeddot((4,9)); closeddot((12,9)); closeddot((16,15)); closeddot((20,15)); closeddot((24,3));
label("$(4,9)$",(4,9),N,fontsize(8pt),Fill(white));
label("$(12,9)$",(12,9),SE,fontsize(8pt),Fill(white));
label("$(16,15)$",(16,15),NW,fontsize(8pt),Fill(white));
label("$(20,15)$",(20,15),NE,fontsize(8pt),Fill(white));
label("$(24,3)$",(24,3),E,fontsize(8pt),Fill(white));
label("Minutes",(13,-3.5),S,fontsize(9pt));
label(rotate(90)*"Calories per minute",(-4.5,9),W,fontsize(9pt));
\end{asy}
\end{center}
\begin{enumerate}[label=(\roman*)]
 \item Find $f'(22)$, $f'(14)$ and $f'(2)$. Indicate units of measure.
 \item Let $k(t)=t\,f(t)$. Find $k'(14)$.
\end{enumerate}
\item The figure shows the graph of a function $f$.
\begin{center}
\begin{asy}
unitsize(0.75cm);
real f(real x){return 2.185-0.716*(x-0.37)^2;}
grid(-2.5,3.2,-3.5,2.5);
axes(-2.8,3.7,-3.6,3.0);
pcurve(f,-2.40,3.15);
real[] xs={-2.27,-0.89,0.37,1.07}; string[] ls={"$a$","$b$","$c$","$d$"};
for (int i=0; i<4; ++i) {draw((xs[i],0),(0,2.5)--(0,-2.5)); label(ls[i],(xs[i],0),2S,fontsize(9pt),Fill(white)); closeddot((xs[i],f(xs[i])));}
label("$O$",(0,0),SW,fontsize(8pt));
\end{asy}
\par\vspace{2pt}
{\small Graph of $f$}
\end{center}
\begin{enumerate}[label=(\roman*)]
 \item Put $f'(b)$, $f'(c)$ and $f'(d)$ in order, from least to greatest.
 \item Which of the following has the greatest value: $f(a)$, $f'(a)$, $f'(c)$, $f(d)-f(c)$ or $\dfrac{f(b)-f(a)}{b-a}$?
 \item How many points on $a<x<d$ have a tangent line parallel to the secant line through $\bigl(a,f(a)\bigr)$ and $\bigl(d,f(d)\bigr)$? Between which two of the labelled points is each one?
\end{enumerate}
\end{enumerate}
\workingspace

\question{practice}{b:pr:tickets}{}{Concert tickets went on sale at noon ($t=0$). The number of people waiting in line to purchase tickets at time $t$ is modeled by a twice-differentiable function $L$ for $0\le t\le9$. Estimate the rate at which the number of people waiting in line was changing at 5:30 \textsc{p.m.} ($t=5.5$). Show the computations that lead to your answer. Indicate units of measure.
\begin{center}
\renewcommand{\arraystretch}{1.3}
\begin{tabular}{|c|c|c|c|c|c|c|c|}\hline
$t$ (hours) & $0$ & $1$ & $3$ & $4$ & $7$ & $8$ & $9$\\ \hline
$L(t)$ (people) & $120$ & $156$ & $176$ & $126$ & $150$ & $80$ & $0$\\ \hline
\end{tabular}
\end{center}}
\workingspace

\question{practice}{b:pr:tea}{}{As a pot of tea cools, its temperature $H(t)$, in degrees Celsius, is modeled by a differentiable function $H$ for $0\le t\le10$, where $t$ is measured in minutes. Approximate the rate at which the temperature of the tea is changing at time $t=3.5$. Show the computations that lead to your answer. Indicate units of measure.
\begin{center}
\renewcommand{\arraystretch}{1.3}
\begin{tabular}{|c|c|c|c|c|c|}\hline
$t$ (minutes) & $0$ & $2$ & $5$ & $9$ & $10$\\ \hline
$H(t)$ (degrees Celsius) & $66$ & $60$ & $52$ & $44$ & $43$\\ \hline
\end{tabular}
\end{center}}
\workingspace

\question{practice}{b:pr:orderslopes}{}{The graph of a differentiable function $f$ is shown. Which of the following is true?}
\begin{center}
\begin{asy}
unitsize(0.8cm);
pair[] P={(-4.23,1.74),(-4.065,1.479),(-3.884,1.216),(-3.702,0.962),(-3.521,0.753),(-3.339,0.590),(-3.158,0.472),(-2.976,0.426),(-2.795,0.481),(-2.613,0.590),(-2.432,0.735),(-2.25,0.917),(-2,1.18),(-1.797,1.452),(-1.615,1.661),(-1.434,1.878),(-1.252,2.060),(-1.071,2.223),(-0.889,2.332),(-0.708,2.414),(-0.526,2.468),(-0.35,2.492),(-0.17,2.499),(0,2.5),(0.17,2.501),(0.35,2.515),(0.563,2.577),(0.744,2.713),(0.926,2.940),(1.107,3.321),(1.289,3.721),(1.47,4.038),(1.652,4.229),(1.833,4.283),(2.015,4.192),(2.196,3.993),(2.377,3.721),(2.559,3.403),(2.74,3.049),(3,2.5),(3.194,2.069),(3.376,1.661),(3.557,1.234),(3.80,0.69)};
guide g; for (pair p: P) g=g..p; draw(g,linewidth(0.9pt));
axes(-4.5,5.1,-0.8,5.3);
xnums(new real[]{-4,-3,-2,-1,1,2,3,4});
closeddot((-2,1.18)); closeddot((0,2.5)); closeddot((3,2.5));
label("$O$",(0,0),SW,fontsize(8pt));
\end{asy}
\par\vspace{2pt}
{\small Graph of $f$}
\end{center}
\twochoices{$f'(-2)<f'(0)<f'(3)$}{$f'(-2)<f'(3)<f'(0)$}{$f'(3)<f'(-2)<f'(0)$}{$f'(3)<f'(0)<f'(-2)$}
\workingspace

\question{practice}{b:pr:countsecant}{}{The graph of a function $f$ with $f(b)>f(a)$ is shown for $a\le x\le b$. The derivative of $f$ exists for all $x$ in the interval $a<x<b$ except $x=0$. For how many values of $c$, for $a<c<b$, does $\displaystyle\lim_{x\to c}\frac{f(x)-f(c)}{x-c}=\frac{f(b)-f(a)}{b-a}$?}
\begin{center}
\begin{asy}
unitsize(0.62cm);
pair[] L={(-4.67,3.2),(-4.63,2.96),(-4.55,2.72),(-4.455,2.48),(-4.4,2.365),(-4.25,2.08),(-4.1,1.855),(-3.95,1.705),(-3.8,1.59),(-3.65,1.505),(-3.5,1.455),(-3.35,1.455),(-3.2,1.49),(-3.05,1.58),(-2.9,1.705),(-2.75,1.855),(-2.6,2.005),(-2.45,2.13),(-2.3,2.23),(-2.15,2.31),(-2.0,2.37),(-1.85,2.395),(-1.7,2.4),(-1.55,2.38),(-1.4,2.34),(-1.25,2.275),(-1.1,2.18),(-0.95,2.06),(-0.8,1.895),(-0.65,1.695),(-0.5,1.44),(-0.395,1.23),(-0.295,0.99),(-0.21,0.75),(-0.13,0.51),(-0.075,0.27),(0,0.03)};
pair[] R={(0,0.03),(0.075,0.27),(0.13,0.51),(0.2,0.75),(0.27,0.99),(0.365,1.23),(0.45,1.41),(0.55,1.595),(0.7,1.79),(0.85,1.925),(1.0,2.005),(1.15,2.045),(1.3,2.045),(1.45,2.015),(1.6,1.955),(1.75,1.86),(1.9,1.75),(2.05,1.625),(2.2,1.525),(2.35,1.46),(2.5,1.46),(2.65,1.51),(2.8,1.6),(2.95,1.725),(3.1,1.9),(3.25,2.12),(3.43,2.45),(3.54,2.7),(3.63,2.95),(3.72,3.2),(3.79,3.45),(3.86,3.7),(3.91,3.95),(3.965,4.2),(4.015,4.45),(4.05,4.7),(4.09,4.95),(4.12,5.2),(4.145,5.445)};
guide g; for (pair p: L) g=g..p; draw(g,linewidth(0.9pt));
guide h; for (pair p: R) h=h..p; draw(h,linewidth(0.9pt));
axes(-5.2,5.5,-0.9,6.6);
xlab(-4.67,"$a$"); xlab(4.145,"$b$");
closeddot((-4.67,3.2)); closeddot((4.145,5.445));
label("$(a,f(a))$",(-4.67,3.2),N,fontsize(9pt));
label("$(b,f(b))$",(4.145,5.445),NW,fontsize(9pt));
label("$O$",(0,0),SW,fontsize(8pt));
\end{asy}
\par\vspace{2pt}
{\small Graph of $f$}
\end{center}
\fourchoices{Zero}{Two}{Three}{Four}
\workingspace

\question[long]{practice}{b:pr:rulesgraph}{}{(a)\enspace The graphs of the linear functions $f$ and $g$ are shown. If $h(x)=f(x)+g(x)$, then $h'(x)=$}
\begin{center}
\begin{asy}
unitsize(0.45cm);
grid(-5.5,5.5,-1.6,9.5);
axes(-5.5,6,-1.6,10);
xnums(new real[]{-5,-4,-3,-2,-1,1,2,3,4,5}); ynums(new real[]{-1,1,2,3,4,5,6,7,8,9});
seg((-2.6,-1.6),(5.3,6.3)); seg((-5.4,5.6),(5.3,-1.5333));
label("$y=f(x)$",(2.6,4.6),NW,fontsize(9pt),Fill(white));
label("$y=g(x)$",(3.3,1.3),NE,fontsize(9pt),Fill(white));
label("$O$",(0,0),SW,fontsize(8pt));
\end{asy}
\end{center}
\fourchoices{$\dfrac53$}{$0$}{$\dfrac13$}{$\dfrac13x+3$}
\begin{enumerate}[label=(\alph*),start=2]
 \item \srctag{AB 2026 FRQ}\enspace Let $f$ be a twice-differentiable function on the closed interval $[-4,4]$ with $f(2)=3$. The graph of $f'$, the derivative of $f$, is shown. For $x>0$, the function $g$ is defined by $g(x)=f(x)-\ln x$. Find $g'(2)$. Show the work that leads to your answer.
\end{enumerate}
\begin{center}
\begin{asy}
unitsize(0.75cm,0.62cm);
axes(-4.7,5.0,-1.6,5.9);
for (int i=-4; i<=4; ++i) if (i!=0) draw((i,0),(0,2.5)--(0,-2.5));
for (int j=-1; j<=5; ++j) if (j!=0) draw((0,j),(2.5,0)--(-2.5,0));
xnums(new real[]{1,2,3,4});
for (int i=-4; i<=-1; ++i) label("$"+string(i)+"$",(i,0),2NW,fontsize(8pt));
ynums(new real[]{-1,1,2,3,4,5});
draw((-4,-0.5)..controls(-3.564,-0.875)and(-3.369,-1)..(-3,-1)..controls(-1.920,-1)and(-1.185,3)..(1,3)..controls(1.427,3)and(1.815,1.902)..(2,1.5)..controls(2.164,1.127)and(2.617,0)..(3,0)..controls(3.429,0)and(3.653,1.289)..(4,5), linewidth(0.9pt));
closeddot((-4,-0.5)); closeddot((-3,-1)); closeddot((1,3)); closeddot((2,1.5)); closeddot((4,5));
label("$(4,5)$",(4,5),W,fontsize(8pt));
label("$(1,3)$",(1,3),N,fontsize(8pt));
label("$(2,1.5)$",(2,1.5),NE,fontsize(8pt));
label("$(-4,-0.5)$",(-4,-0.5),SW,fontsize(8pt));
label("$(-3,-1)$",(-3,-1),SE,fontsize(8pt));
label("$O$",(0,0),SW,fontsize(8pt));
\end{asy}
 \par\nopagebreak\vspace{2pt}
{\small Graph of $f'$}
\end{center}
\workingspace

% ---------------------------------------------------------------------
\Needspace{20\baselineskip}
\subsection{Derivatives and rates of change}
\begin{notesbox}{Notes}
\begin{itemize}
 \item The average rate of change of $f$ on $[a,b]$ is the change in $f$ divided by the change in $x$:
 \[
 \text{average rate of change}=\frac{f(b)-f(a)}{b-a}.
 \]
 It is the slope of the secant line through $\bigl(a,f(a)\bigr)$ and $\bigl(b,f(b)\bigr)$. Its units are the units of $f$ per unit of $x$.
\end{itemize}
\begin{center}
\begin{asy}
unitsize(0.9cm);
real f(real x){return 1+0.5*x+0.6*sin(1.4*x);}
real a=1, b=5;
pair A=(a,f(a)), B=(b,f(b));
axes(-0.4,7,-0.4,5.2);
curve(f,0.5,6.2);
pair d=unit(B-A);
draw((A-0.9*d)--(B+0.9*d), linewidth(0.7pt)+dashed);
draw((a,0)--A, gray(0.5)+dashed); draw((b,0)--B, gray(0.5)+dashed);
draw((0,f(a))--A, gray(0.5)+dashed); draw((0,f(b))--B, gray(0.5)+dashed);
closeddot(A); closeddot(B);
xlab(a,"$a$"); xlab(b,"$b$");
label("$f(a)$",(0,f(a)),W,fontsize(9pt)); label("$f(b)$",(0,f(b)),W,fontsize(9pt));
label("$\bigl(a,f(a)\bigr)$",A+(0.15,0.5),fontsize(9pt)); label("$\bigl(b,f(b)\bigr)$",B+(-0.75,0.4),fontsize(9pt));
label("$y=f(x)$",(6.2,f(6.2)),E,fontsize(9pt));
label("secant line",A+2.3*d,NW,fontsize(9pt));
draw((b+0.35,f(a))--(b+0.35,f(b)),Arrows(3)); label("$f(b)-f(a)$",(b+0.35,(f(a)+f(b))/2),E,fontsize(8pt));
draw((a,f(a)-0.25)--(b,f(a)-0.25),Arrows(3)); label("$b-a$",((a+b)/2,f(a)-0.25),S,fontsize(8pt));
draw((a,f(a))--(b+0.35,f(a)), gray(0.6)+dotted);
\end{asy}
\par\small The average rate of change is the slope of the secant line: $\dfrac{f(b)-f(a)}{b-a}$.
\end{center}
Let the interval shrink to the point $a$:
\begin{center}
\renewcommand{\arraystretch}{2.4}
\begin{tabular}{c|c|c}
average rate of change & & instantaneous rate of change\\ \hline
$\displaystyle\frac{f(x)-f(a)}{x-a}$ & $\xrightarrow{\ x\to a\ }$ & $\displaystyle\lim_{x\to a}\frac{f(x)-f(a)}{x-a}=f'(a)$\\
$\displaystyle\frac{f(a+h)-f(a)}{h}$ & $\xrightarrow{\ h\to0\ }$ & $\displaystyle\lim_{h\to0}\frac{f(a+h)-f(a)}{h}=f'(a)$
\end{tabular}
\end{center}
\textbf{$f'(a)$ is the slope of the tangent line and the instantaneous rate of change, and you get it from secant slopes (average rates) by letting the second point move onto the first.}
\begin{itemize}
 \item The instantaneous rate of change equals the average rate of change on $[a,b]$: solve
 \[
 f'(c)=\frac{f(b)-f(a)}{b-a}
 \]
 and keep only the solutions with $a\le c\le b$.
\end{itemize}
\end{notesbox}

\Needspace{0.45\textheight}
% (a) AB 2018 FRQ 5(a); (b) AB 2004 FRQ 1(d); (c) AB 2019 FRQ 3(a) = BC 2019 FRQ 3(a); (d) new.
\exampletitle{a:ex:aroc}{Average rate of change}
\begin{enumerate}[label=(\alph*)]
 \item Find the average rate of change of $f(x)=e^x\cos x$ on the interval $0\le x\le\pi$.
 \item \calc The traffic flow at an intersection, in cars per minute, is modeled by $F(t)=82+4\sin\Bigl(\dfrac t2\Bigr)$ for $0\le t\le30$, where $t$ is measured in minutes. Find the average rate of change of the traffic flow over $10\le t\le15$. Indicate units of measure.
 \item Let $f(x)=\begin{cases}\sqrt{9-x^2} & \text{for } -3\le x\le0\\[2pt] -x+3\cos\Bigl(\dfrac{\pi x}{2}\Bigr) & \text{for } 0<x\le4.\end{cases}$\quad Find the average rate of change of $f$ on the interval $-3\le x\le4$.
 \item Let $f(x)=x^2+x$. The average rate of change of $f$ over $[1,b]$ is $7$, where $b>1$. Find $b$.
\end{enumerate}
\end{questionbox}
\begin{examplesolution}
\relax
\begin{enumerate}[label=(\alph*)]
 \item Only the values at the endpoints are needed (no product rule).
 \begin{align*}
 \frac{f(\pi)-f(0)}{\pi-0}&=\frac{e^\pi\cos\pi-e^0\cos0}{\pi}\\
 &=\frac{-e^\pi-1}{\pi}. &&\why{$\cos\pi=-1$, $\cos0=1$, $e^0=1$}
 \end{align*}
 \item The function is itself a rate (cars per minute), so the answer is in cars per minute per minute.
 \begin{align*}
 \frac{F(15)-F(10)}{15-10}&=\frac{85.752-78.164}{5} &&\why{calculator, radians}\\
 &=1.518\text{ cars per minute per minute (or }1.517).
 \end{align*}
 \item Put each endpoint into its own piece: $x=-3$ is in the first piece, $x=4$ in the second.
 \begin{align*}
 f(-3)&=\sqrt{9-9}=0, &\qquad f(4)&=-4+3\cos(2\pi)=-4+3=-1,
 \end{align*}
 \[
 \frac{f(4)-f(-3)}{4-(-3)}=\frac{-1-0}{7}=-\frac17.
 \]
 (The graph has a corner at $x=0$; that does not matter.)
 \item
 \begin{align*}
 \frac{f(b)-f(1)}{b-1}&=7 &&\why{average rate of change over $[1,b]$}\\
 \frac{b^2+b-2}{b-1}&=7 &&\why{$f(1)=2$}\\
 \frac{(b+2)(b-1)}{b-1}&=7\\
 b+2&=7 &&\why{$b\ne1$}\\
 b&=5.
 \end{align*}
\end{enumerate}
\end{examplesolution}

\practiceheading
\question[long]{practice}{a:pr:aroc}{}{
\begin{enumerate}[label=(\alph*)]
 \item Find the average rate of change of $f(x)=e^x\sin x$ on the interval $0\le x\le\dfrac{\pi}{2}$.
 \item \calc Water flows into a tank at the rate $R(t)=30+5\cos\Bigl(\dfrac t3\Bigr)$ gallons per hour, where $t$ is measured in hours. Find the average rate of change of $R$ over $2\le t\le8$. Indicate units of measure.
 \item Let $f(x)=x^2-3x$. The average rate of change of $f$ over $[2,b]$ is $6$, where $b>2$. Find $b$.
\end{enumerate}}
\workingspace

\question{practice}{a:pr:costwox}{}{What is the average rate of change of $y=\cos(2x)$ on the interval $\Bigl[0,\dfrac{\pi}{2}\Bigr]$?}
\fourchoices{$-\dfrac{4}{\pi}$}{$-1$}{$0$}{$\dfrac{4}{\pi}$}
\workingspace

\question[long]{practice}{a:pr:fprimerow}{}{Let $f$ be a differentiable function. Selected values of $f$ and its derivative $f'$ are shown in the table.
\begin{center}
\renewcommand{\arraystretch}{1.3}
\begin{tabular}{|c|c|c|c|c|c|}\hline
$x$ & $-2$ & $3$ & $5$ & $12$ & $15$\\ \hline
$f(x)$ & $-10$ & $-2$ & $-1$ & $15$ & $36$\\ \hline
$f'(x)$ & $12$ & $5$ & $0$ & $9$ & $2$\\ \hline
\end{tabular}
\end{center}
\begin{enumerate}[label=(\alph*)]
 \item Find the average rate of change of $f$ over the interval $3\le x\le12$.
 \item Find the average rate of change of $f'$ over the interval $3\le x\le12$.
 \item On which of the intervals $[-2,3]$, $[3,5]$, $[5,12]$, $[12,15]$ is the average rate of change of $f$ greatest?
\end{enumerate}}
\workingspace

\question{practice}{a:pr:arocgraph}{}{The function $f$ is defined on $-5\le x\le c$, where $c>0$ and $f(c)=0$. Its graph, shown below, consists of three line segments and a quarter of a circle with center $(-3,0)$ and radius $2$. Find the average rate of change of $f$ over the interval $[-5,0]$.}
\begin{center}
\begin{asy}
unitsize(0.62cm);
axes(-5.8,6.3,-2.6,2.9);
for (real x : new real[]{-5,-4,-3,-2,-1}) draw((x,0),(0,2.5)--(0,-2.5));
label("$-5$",(-5,0),3SW,fontsize(8pt)); label("$-3$",(-3,0),2S,fontsize(8pt)); label("$-1$",(-1,0),2S,fontsize(8pt));
xlab(5.27,"$c$");
ynums(new real[]{-2,-1,1,2});
draw(arc((-3,0),2,180,270), linewidth(0.9pt));
seg((-3,-2),(-2,0)); seg((-2,0),(0,2)); seg((0,2),(5.27,0));
\end{asy}
\par Graph of $f$
\end{center}
\workingspace

\question{practice}{a:pr:arocpiece}{}{Let $g$ be the function given by $g(x)=\begin{cases}x^2+1 & \text{for } x<3\\ 13-x & \text{for } x\ge3.\end{cases}$\\[4pt]
What is the average rate of change of $g$ on the closed interval $[1,5]$?}
\fourchoices{$\dfrac12$}{$\dfrac32$}{$5$}{Undefined}
\workingspace

\Needspace{0.45\textheight}
% (a) is Unit 2 bank Q161 (= Unit 1 bank Q161, P); (b) AB 2016 MC 26. (The juice bottle, 2022 FRQ 1(a), moved to Example 10.)
\exampletitle{a:ex:definition}{Instantaneous rate of change}
\begin{enumerate}[label=(\alph*)]
 \item The size of a population of bacteria is modeled by the function $P$, where $P(t)$ gives the number of bacteria and $t$ gives the number of hours after midnight for $0\le t\le10$. The graph of $P$ and the line tangent to $P$ at $t=8$ are shown. Which of the following gives the best estimate for the instantaneous rate of change of $P$ at $t=8$?
 \begin{center}
 \begin{asy}
 unitsize(0.6cm,0.2cm);
 axes(-0.6,11,-1.5,17,"$t$","$P(t)$");
 xnums(new real[]{2,4,6,8,10});
 string[] lab={"2{,}500","5{,}000","7{,}500","10{,}000","12{,}500","15{,}000"};
 for (int k=1; k<=6; ++k) {draw((0,2.5k),(2.5,0)--(-2.5,0)); label("$"+lab[k-1]+"$",(0,2.5k),2W,fontsize(8pt));}
 curve(new real(real t){return 5*(t/8)^4.5;}, 0, 10);
 draw((6.3,5+2.8125*(6.3-8))--(10,5+2.8125*2), linewidth(0.7pt));
 \end{asy}
 \end{center}
 \twochoices{$P(8)$}{The slope of the line joining $\bigl(0,P(0)\bigr)$ and $\bigl(8,P(8)\bigr)$}{The slope of the line joining $\bigl(0,P(0)\bigr)$ and $\bigl(10,P(10)\bigr)$}{The slope of the line joining $\bigl(7.9,P(7.9)\bigr)$ and $\bigl(8.1,P(8.1)\bigr)$}
 \item Let $f(x)=x^3+5x$. For what value of $x$ in $[1,3]$ does the instantaneous rate of change of $f$ equal the average rate of change of $f$ on $[1,3]$?
\end{enumerate}
\end{questionbox}
\begin{examplesolution}
\relax
\begin{enumerate}[label=(\alph*)]
 \item The instantaneous rate at $t=8$ is $P'(8)$. The best estimate is the average rate of change over the shortest interval containing $8$, here $[7.9,8.1]$. The answer is (D).\\
 (A) is a number of bacteria, not a rate; (B) and (C) are average rates over long intervals.
 \item Average rate of change: $\dfrac{f(3)-f(1)}{3-1}=\dfrac{42-6}{2}=18$. Instantaneous rate: $f'(x)=3x^2+5$.
 \[
 3x^2+5=18\imp x^2=\frac{13}{3}\imp x=\pm\sqrt{\frac{13}{3}}.
 \]
 Only $x=\sqrt{\dfrac{13}{3}}\approx2.08$ is in $[1,3]$; reject $-\sqrt{\dfrac{13}{3}}$.
\end{enumerate}
\end{examplesolution}

\practiceheading
% (a) is Unit 2 bank Q162 (= Unit 1 bank Q162, P); (b) BC 2016 MC 26. (The coffee table moved to the Example 10 practice.)
\question[long]{practice}{a:pr:instant}{}{(a)\enspace The size of a population of fish in a pond is modeled by the function $P$, where $P(t)$ gives the number of fish and $t$ gives the number of years after the fish were introduced to the pond, for $0\le t\le10$. The graph of $P$ and the line tangent to $P$ at $t=4$ are shown. Which of the following gives the best estimate for the instantaneous rate of change of $P$ at $t=4$?}
\begin{center}
\begin{asy}
unitsize(0.6cm,0.2cm);
axes(-0.6,11,-2,17,"$t$","$P(t)$");
xnums(new real[]{2,4,6,8,10});
string[] lab={"250","500","750","1{,}000","1{,}250","1{,}500"};
for (int k=1; k<=6; ++k) {draw((0,2.5k),(2.5,0)--(-2.5,0)); label("$"+lab[k-1]+"$",(0,2.5k),2W,fontsize(8pt));}
real P(real t) {return 15.5/(1+exp(-1.27*(t-5.25)));}
curve(P, 0, 10);
real m=1.27*P(4)*(1-P(4)/15.5);
draw((2.9,P(4)+m*(2.9-4))--(5.5,P(4)+m*(5.5-4)), linewidth(0.7pt));
\end{asy}
\end{center}
\twochoices{$P(4)$}{The slope of the line joining $\bigl(0,P(0)\bigr)$ and $\bigl(4,P(4)\bigr)$}{The slope of the line joining $\bigl(0,P(0)\bigr)$ and $\bigl(8,P(8)\bigr)$}{The slope of the line joining $\bigl(3.9,P(3.9)\bigr)$ and $\bigl(4.1,P(4.1)\bigr)$}
\begin{enumerate}[label=(\alph*),start=2]
 \item Let $f(x)=x^3-2x^2+5x-16$. For what value of $x$ in $[0,5]$ does the instantaneous rate of change of $f$ equal the average rate of change of $f$ over that interval?
\end{enumerate}
\workingspace

\question{practice}{a:pr:plane}{}{An airplane's height above sea level is modeled by the differentiable function $y=h(t)$, where $t$ is the number of minutes after takeoff and $h(t)$ is measured in feet. The graph of $h$ is shown for $0\le t\le15$. Of the following, which is the best estimate of the rate at which the plane is rising above sea level, in feet per minute, at time $t=9$?}
\begin{center}
\begin{asy}
unitsize(0.42cm,0.1cm);
axes(-0.8,16.6,-3,44,"$t$","$y$");
xnums(new real[]{5,10,15});
for (int k=1; k<=4; ++k) {draw((0,10k),(2.5,0)--(-2.5,0)); label("$"+string(10k)+"{,}000$",(0,10k),2W,fontsize(8pt));}
real h(real t) {return 0.150+0.495t^2-0.022t^3;}
curve(h, 0, 15);
closeddot((0,h(0))); closeddot((5,h(5))); closeddot((7,h(7))); closeddot((11,h(11))); closeddot((15,h(15)));
label("$(0, 150)$",(0,h(0)),SE,fontsize(8pt));
label("$(5, 9{,}775)$",(5,h(5)),SE,fontsize(8pt));
label("$(7, 16{,}859)$",(7,h(7)),NW,fontsize(8pt));
label("$(11, 30{,}763)$",(11,h(11)),NW,fontsize(8pt));
label("$(15, 37{,}275)$",(15,h(15)),SE,fontsize(8pt));
\end{asy}
\end{center}
\twochoices{$\dfrac{37{,}275-150}{15-0}$}{$\dfrac{37{,}275-9{,}775}{15-5}$}{$\dfrac{37{,}275-16{,}859}{15-7}$}{$\dfrac{30{,}763-16{,}859}{11-7}$}
\workingspace

\question{practice}{b:pr:avgrate}{}{\calc Let $f$ be the function defined by $f(x)=x+\ln x$. What is the value of $c$ for which the instantaneous rate of change of $f$ at $x=c$ is the same as the average rate of change of $f$ over $[1,4]$?}
\fourchoices{$0.456$}{$1.244$}{$2.164$}{$2.342$}
\workingspace

% ---------------------------------------------------------------------
\Needspace{20\baselineskip}
\subsection{Rates of change in context}
\begin{notesbox}{Notes}
The instantaneous rate of change of a quantity $Q$ with respect to $t$ at $t=a$ is the derivative $Q'(a)$. Its units are the units of $Q$ per unit of $t$:
\[
\text{rate of change of $Q$ at $t=a$}=Q'(a),\qquad \text{units of }Q'(a)=\frac{\text{units of }Q}{\text{units of }t}.
\]
\end{notesbox}

\Needspace{0.45\textheight}
% (a) AB 2011 Form B FRQ 1(c); (b) AB 2016 Alt FRQ 1(a), rate only (the "increasing or decreasing" sentence is Unit 5).
\exampletitle{c:ex:rates}{Rates of change in context}
\begin{enumerate}[label=(\alph*)]
 \item \srctag{AB 2011 Form B FRQ}\enspace\calc\ A cylindrical can of radius $10$ millimeters is used to measure rainfall in Stormville. The can is initially empty, and rain enters the can during a $60$-day period. The height of water in the can is modeled by the function $S$, where $S(t)$ is measured in millimeters and $t$ is measured in days for $0\le t\le60$. The rate at which the height of the water is rising in the can is given by $S'(t)=2\sin(0.03t)+1.5$. Assuming no evaporation occurs, at what rate is the volume of water in the can changing at time $t=7$? Indicate units of measure.
 \item \srctag{AB 2016 Alt FRQ}\enspace A company produces and sells chili powder. The company's weekly profit on the sale of $x$ kilograms of chili powder is modeled by the function $P$ given by $P(x)=48x+1.4x^2-0.05x^{2.8}-270$, where $P(x)$ is in dollars and $0\le x\le80$. Find the rate, in dollars per kilogram, at which the company's weekly profit is changing when it sells $32$ kilograms of chili powder.
\end{enumerate}
\end{questionbox}
\begin{examplesolution}
\relax
\begin{enumerate}[label=(\alph*)]
 \item
 \steph{Step 1: The volume in terms of the height.}
 ``At what rate is the volume changing'' asks for $V'(7)$. The radius is fixed at $10$ millimeters, so
 \[
 V=\pi(10)^2S=100\pi S\imp V'(t)=100\pi S'(t). \qquad\why{constant-multiple rule}
 \]
 \steph{Step 2: Substitute $t=7$.}
 \begin{align*}
 V'(7)&=100\pi\bigl(2\sin(0.21)+1.5\bigr)=602.218. &&\why{radian mode}
 \end{align*}
 The volume is changing at $602.218$ cubic millimeters per day.
 \item
 The rate is $P'(32)$.
 \begin{align*}
 P'(x)&=48+2.8x-0.14x^{1.8} &&\why{power rule: $0.05\cdot2.8=0.14$}\\
 P'(32)&=48+2.8(32)-0.14(32)^{1.8}\\
 &=48+89.6-0.14(512)=65.92. &&\why{$32^{1.8}=\bigl(2^5\bigr)^{1.8}=2^9=512$}
 \end{align*}
 The weekly profit is changing at $65.92$ dollars per kilogram.
\end{enumerate}
\end{examplesolution}

\practiceheading
% (a)-(b) new: a variation of the example, part for part.
\question[long]{practice}{c:ratesvar}{}{Indicate units of measure in each answer.
\begin{enumerate}[label=(\alph*)]
 \item \calc\ A cylindrical tank of radius $3$ meters is being filled with water. The height of the water in the tank is modeled by the function $h$, where $h(t)$ is measured in meters and $t$ is measured in hours. The rate at which the height of the water is rising is given by $h'(t)=0.5+0.2\cos(0.1t)$. At what rate is the volume of water in the tank changing at time $t=5$?
 \item A bakery's weekly revenue from selling $x$ kilograms of bread is modeled by $R(x)=60x+0.8x^2-0.02x^{2.5}-150$, where $R(x)$ is in dollars and $0\le x\le60$. Find the rate, in dollars per kilogram, at which the weekly revenue is changing when the bakery sells $25$ kilograms of bread.
\end{enumerate}}
\workingspace

\question{practice}{c:tankw}{}{The number of gallons of water in a storage tank at time $t$, in minutes, is modeled by $w(t)=25-t^2$ for $0\le t\le5$. At what rate, in gallons per minute, is the amount of water in the tank changing at time $t=3$ minutes?}
\fourchoices{$66$}{$16$}{$-3$}{$-6$}
\workingspace

\question{practice}{c:wave}{}{The velocity $v$, in meters per second, of a certain type of wave is given by $v(h)=3\sqrt h$, where $h$ is the depth, in meters, of the water through which the wave moves. What is the rate of change, in meters per second per meter, of the velocity of the wave with respect to the depth of the water, when the depth is $2$ meters?}
\fourchoices{$-\dfrac{3}{4\sqrt2}$}{$\dfrac{3}{2\sqrt2}$}{$\dfrac{3}{\sqrt2}$}{$4\sqrt2$}
\workingspace

% AB 2016 FRQ 1(a): the second sentence (increasing or decreasing, with a reason) is left out; it is Unit 5.
% ---------------------------------------------------------------------
\Needspace{20\baselineskip}
\subsection{Derivatives on the calculator}
\begin{notesbox}{Notes}
\begin{itemize}
 \item Calculator in radian mode. Give answers correct to three decimal places; store long values (\texttt{STO}$\to$) instead of rounding them halfway.
 \item $f'(a)$ on the TI-84: \texttt{MATH 8} (\texttt{nDeriv}). The calculator differentiates any formula, composites included.
 \[
 \texttt{nDeriv(}f(\texttt{X})\texttt{,X,}a\texttt{)}\qquad\text{or, on a MathPrint screen,}\qquad \frac{d}{d\texttt{X}}\bigl(f(\texttt{X})\bigr)\Big|_{\texttt{X}=a}.
 \]
 Quicker: type $f$ into \texttt{Y1} once and use \texttt{nDeriv(Y1,X,}$a$\texttt{)} (\texttt{Y1} is under \texttt{VARS}, \texttt{Y-VARS}, \texttt{Function}).
 \item On the test, write the derivative, the value and the units, e.g.\ $r'(4)=-39.932$ gallons per hour per hour. For a piecewise function, use the piece that contains $a$.
 \item Solving $f'(x)=k$: graph
 \[
 \texttt{Y2}=\texttt{nDeriv(Y1,X,X)},\qquad \texttt{Y3}=k
 \]
 with the window set to the interval, then \texttt{2nd TRACE} (\texttt{CALC}) \texttt{5:intersect}. Or graph $f'(x)-k$ and use \texttt{2:zero}. If $f'(x)$ is given, type it in directly. To count the solutions, count the crossings on the interval (zoom in where the graphs nearly touch).
 \item Which equation to solve:
 \begin{itemize}
  \item tangent parallel to a line of slope $m$: $f'(x)=m$;\quad horizontal tangent: $f'(x)=0$;
  \item tangent to $f$ at $x=2$ parallel to the tangent to $g$ at $x=a$: $g'(a)=f'(2)$;
  \item tangent at $\bigl(k,f(k)\bigr)$ through the origin, or with $y$-intercept $b$. The tangent is $y=f(k)+f'(k)(x-k)$; put $x=0$:
  \[
  y\text{-intercept}=f(k)-kf'(k),\qquad \text{through the origin: } f(k)-kf'(k)=0;
  \]
  \item slope where the graph crosses the $x$-axis: solve $f(x)=0$ first (\texttt{2:zero}); the root is stored in \texttt{X}, so \texttt{nDeriv(Y1,X,X)} on the home screen gives the slope there.
 \end{itemize}
\end{itemize}
\end{notesbox}

\Needspace{0.45\textheight}
\exampletitle{c:ex:calc}{Derivatives on the calculator}
\calc Give each answer correct to three decimal places.
\begin{enumerate}[label=(\alph*)]
 \item Let $f(x)=x^3-2x^2+\sin\bigl(x^2\bigr)$. Find $f'(1.5)$. Of $x=-1$, $0$, $1$ and $2$, at which does the line tangent to the graph of $f$ have the greatest slope?
 \item Water flows into a tank at the rate
 \[
 r(t)=\begin{cases}40+30e^{-0.4t} & \text{for } 0\le t\le3\\ 20+15\ln\bigl(t^2-5\bigr) & \text{for } 3<t\le8,\end{cases}
 \]
 where $r(t)$ is measured in liters per hour and $t$ in hours. Find $r'(4)$. Indicate units of measure.
 \item The graph of $y=e^{\cos x}-2$ crosses the $x$-axis at one point in the interval $[0,2]$. What is the slope of the graph at this point?
 \item The derivative of the function $g$ is $g'(x)=2\cos\bigl(x^2\bigr)+1$.
 \begin{enumerate}[label=(\roman*)]
  \item For how many values of $x$ in $0\le x\le3$ is the line tangent to the graph of $g$ parallel to the line $y=2x+5$?
  \item Find all values of $x$ in $0<x<2$ at which the graph of $g$ has a horizontal tangent line.
 \end{enumerate}
 \item The line tangent to the graph of $y=\sin x$ at the point $(k,\sin k)$, where $\pi<k<2\pi$, passes through the origin. Find $k$.
 \item The line tangent to the graph of $p(x)=\sqrt{x^2+5}$ at $x=2$ is parallel to the line tangent to the graph of $q(x)=x^3-x$ at $x=a$, where $a>0$. Find $a$.
\end{enumerate}
\end{questionbox}
\begin{examplesolution}
\relax
\begin{enumerate}[label=(\alph*)]
 \item \scratch{\texttt{Y1}$=$\texttt{X\^{}3-2X\^{}2+sin(X\^{}2)}; then \texttt{nDeriv(Y1,X,1.5)}, and the same with $-1$, $0$, $1$, $2$.}
 \[
 f'(1.5)=-1.135.
 \]
 The slopes are $f'(-1)=5.919$, $f'(0)=0$, $f'(1)=0.081$, $f'(2)=1.385$. The greatest slope is at $x=-1$.
 \item $t=4$ is in the second piece, $3<t\le8$.
 \scratch{\texttt{nDeriv(20+15ln(X\^{}2-5),X,4)}.}
 \[
 r'(4)=10.909\ \text{liters per hour per hour}.
 \]
 The first piece at $t=4$ gives $-2.423$, a trap answer.
 \item First find the point, then the slope there.
 \steph{Step 1: Find where the graph crosses the $x$-axis.}
 \scratch{\texttt{Y1}$=$\texttt{e\^{}(cos(X))-2}, window $0\le x\le2$, \texttt{2:zero}.}
 \[
 e^{\cos x}-2=0\imp x=0.80495\ldots
 \]
 \steph{Step 2: The slope there.}
 \scratch{home screen: \texttt{nDeriv(Y1,X,X)} (\texttt{X} holds the root).}
 \[
 \frac{dy}{dx}\bigg|_{x=0.80495\ldots}=-1.442.
 \]
 The trap answer $0.805$ is the $x$-value, not the slope.
 \item \begin{enumerate}[label=(\roman*)]
  \item The line $y=2x+5$ has slope $2$, so solve $g'(x)=2$ on $[0,3]$.
  \scratch{\texttt{Y2}$=$\texttt{2cos(X\^{}2)+1}, \texttt{Y3}$=2$, window $0\le x\le3$; count the crossings.}
  \[
  2\cos\bigl(x^2\bigr)+1=2\imp x=1.023,\ 2.288,\ 2.707.
  \]
  There are three values.
  \item Horizontal tangent: $g'(x)=0$.
  \[
  2\cos\bigl(x^2\bigr)+1=0\imp x=1.447 \quad\text{in } 0<x<2.
  \]
  (The next zero, $x=2.047$, is outside the interval.)
 \end{enumerate}
\begin{center}
\begin{asy}
unitsize(1.9cm,0.8cm);
grid(0,3,-1,3);
axes(-0.2,3.3,-1.3,3.5);
xnums(new real[]{1,2,3}); ynums(new real[]{-1,1,2,3});
pcurve(new real(real x){return 2*cos(x^2)+1;}, 0, 3);
draw((0,2)--(3,2), dashed+linewidth(0.6pt));
label("$y=2$",(3,2),E,fontsize(9pt));
label("$y=g'(x)$",(1.15,3.1),E,fontsize(9pt));
closeddot((1.0233,2)); closeddot((2.2882,2)); closeddot((2.7075,2));
closeddot((1.4472,0));
\end{asy}
\end{center}
 \item Write the $y$-intercept of the tangent in terms of $k$, then solve.
 \steph{Step 1: The $y$-intercept of the tangent.}
 With $f(x)=\sin x$, the tangent at $(k,\sin k)$ is $y=\sin k+\cos k\,(x-k)$. At $x=0$:
 \[
 y\text{-intercept}=\sin k-k\cos k.
 \]
 \steph{Step 2: Through the origin.}
 \[
 \sin k-k\cos k=0\imp k=4.493 \quad\text{in } \pi<k<2\pi.
 \]
 \scratch{\texttt{Y1}$=$\texttt{sin(X)-Xcos(X)}, window $\pi\le x\le2\pi$, \texttt{2:zero}.}
 \item Parallel lines have equal slopes.
 \steph{Step 1: The slope of the tangent to $p$ at $x=2$.}
 \[
 p'(2)=0.6667 \qquad \Bigl(\text{exactly } \frac23\Bigr).
 \]
 \steph{Step 2: Parallel lines have equal slopes.}
 \begin{align*}
 q'(a)&=p'(2)\\
 3a^2-1&=\frac23 &&\why{$q'(x)=3x^2-1$}\\
 a&=0.745. &&\why{$a>0$}
 \end{align*}
\end{enumerate}
\end{examplesolution}

\practiceheading
\question{practice}{c:greatest}{}{\calc Let $f$ be the function given by $f(x)=x^4+\dfrac12x^3-5x^2+\tan\Bigl(\dfrac x2\Bigr)$. Of the following values of $x$, at which does the line tangent to the graph of $f$ have the greatest slope?}
\fourchoices{$x=-2$}{$x=-1$}{$x=0$}{$x=1$}
\workingspace

\question{practice}{c:etan}{}{\calc The graph of $y=e^{\tan x}-2$ crosses the $x$-axis at one point in the interval $[0,1]$. What is the slope of the graph at this point?}
\fourchoices{$0.606$}{$2.242$}{$2.961$}{$3.747$}
\workingspace

\question{practice}{c:rainpiece}{}{\calc A $1000$-gallon tank is used to collect rainwater. The tank is empty at time $t=0$. Rainwater enters the tank at a rate modeled by the piecewise-defined function $r$ defined by
\[
r(t)=\begin{cases}200e^{-0.21t} & \text{for } 0\le t\le2\\ 115+25\sin\bigl(0.2t^2\bigr) & \text{for } 2<t\le5,\end{cases}
\]
where $r(t)$ is measured in gallons per hour and $t$ is measured in hours. Find $r'(4)$. Indicate units of measure.}
\workingspace

\question{practice}{c:count}{}{\calc The derivative of the function $f$ is given by $f'(x)=1-5\cos\bigl(3x^2-2\bigr)$. For how many values of $x$ in the interval $0\le x\le2$ is the line tangent to the graph of $y=f(x)$ at the point $\bigl(x,f(x)\bigr)$ parallel to the line $y=1+4x$?}
\fourchoices{One}{Two}{Three}{Four}
\workingspace

\question{practice}{c:origin}{}{\calc Line $\ell$ is tangent to the graph of $y=\cos x$ at the point $(k,\cos k)$, where $0<k<\pi$. For what value of $k$ does line $\ell$ pass through the origin?}
\twochoices{$0.860$}{$1.571$}{$2.798$}{There is no such value of $k$.}
\workingspace

\question{practice}{c:parallel}{}{\calc Let $f$ be the function defined by $f(x)=\ln\bigl(x^2+1\bigr)$, and let $g$ be the function defined by $g(x)=x^5+x^3$. The line tangent to the graph of $f$ at $x=2$ is parallel to the line tangent to the graph of $g$ at $x=a$, where $a$ is a positive constant. What is the value of $a$?}
\fourchoices{$0.246$}{$0.430$}{$0.447$}{$0.790$}
\workingspace

% ---------------------------------------------------------------------
\Needspace{\ifshowspace 0.4\textheight\else 14\baselineskip\fi}
\subsection*{Exercise \thesection}\addcontentsline{toc}{subsection}{Exercise \thesection}
% Helper A: PROPOSED Exercise 2 questions (Section 2, 2.1-2.3). Lee chooses among them.
% AB 2009 Form B FRQ 3(b) (P): the stem's "twice differentiable, with f''(x) > 0" is dropped (f'' is Unit 3);
% the graph shows the shape of the curve.
\question{exercise}{a:ex2:count}{}{A continuous function $f$ is defined on the closed interval $-4\le x\le6$. The graph of $f$ consists of a line segment and a curve that is tangent to the $x$-axis at $x=3$, as shown. For how many values of $a$, $-4\le a<6$, is the average rate of change of $f$ on the interval $[a,6]$ equal to $0$? Give a reason for your answer.}
\begin{center}
\begin{asy}
unitsize(0.75cm);
axes(-4.7,7,-1.3,2.7);
xnums(new real[]{-3,-2,-1,1,2,3,4,5,6});
draw((-4,0),(0,2.5)--(0,-2.5)); label("$-4$",(-4,0),2N,fontsize(8pt));
ynums(new real[]{1,2});
seg((-4,-2/3),(0,2));
curve(new real(real x){return 2/9*(x-3)^2;}, 0, 3);
curve(new real(real x){return 1/9*(x-3)^2;}, 3, 6);
closeddot((-4,-2/3)); closeddot((6,1));
label("$(6,1)$",(6,1),N,fontsize(8pt));
\end{asy}
\par Graph of $f$
\end{center}
\workingspace

% Unit 2 bank Q2 is part (a); Unit 2 bank Q126 is part (e); (b)-(d) new ((d) is the d/dx(f(a)) tidbit).
\question[long]{exercise}{a:ex2:mixed}{}{Let $f(x)=x^4-5x$.
\begin{enumerate}[label=(\alph*)]
 \item Find the average rate of change of $f$ on the closed interval $[0,3]$.
 \item Find $\displaystyle\lim_{h\to0}\frac{f(1+h)-f(1)}{h}$.
 \item Find $\displaystyle\lim_{x\to2}\frac{f(x)-6}{x-2}$.
 \item Find $\dfrac{d}{dx}\bigl(f(2)\bigr)$.
 \item If $g(x)=\tan x$, find $\displaystyle\lim_{x\to\frac{\pi}{4}}\frac{g(x)-g\bigl(\frac{\pi}{4}\bigr)}{x-\frac{\pi}{4}}$.
\end{enumerate}}
\workingspace

% New: which limits equal f'(2) (step -h, reversed order, and a symmetric quotient that gives 2f'(2)).
\starquestion{exercise}{a:ex2:forms}{}{Let $f$ be a differentiable function with $f(2)=3$ and $f'(2)=5$. Which of the following limits are equal to $5$?
\begin{statements}
 \item $\displaystyle\lim_{h\to0}\frac{f(2+h)-3}{h}$
 \item $\displaystyle\lim_{x\to2}\frac{3-f(x)}{2-x}$
 \item $\displaystyle\lim_{h\to0}\frac{f(2+h)-f(2-h)}{h}$
\end{statements}}
\fourchoices{I only}{I and II only}{I and III only}{I, II and III}
\workingspace

% Helper B: PROPOSED Exercise 2 questions (Lee chooses among them).
% b:ex2:slopefield  (a) is AB 1998 FRQ 4(a); (b), (c) added (product rule, d/dx(g(1)) tidbit).
%                   AB 1998 FRQ 4(b) is in the question bank (4.6), so no tangent line to f is asked.
% b:ex2:xn          AB 2004 Form B FRQ 6(b) (tangent at a general power, triangle area).
% b:ex2:order       AB 2013 MC 88 (Intl), choice "f'(0) < f'(-1) < f'(1)" dropped.
% b:ex2:fastest     AB/BC 2022 MC 1 (Part B, Intl).
\question[long]{exercise}{b:ex2:slopefield}{}{Let $f$ be a function with $f(1)=4$ such that for all points $(x,y)$ on the graph of $f$ the slope is given by $\dfrac{3x^2+1}{2y}$.
\begin{enumerate}[label=(\alph*)]
 \item Find the slope of the graph of $f$ at the point where $x=1$.
 \item Let $g(x)=x^2f(x)$. Write an equation of the line tangent to the graph of $g$ at $x=1$.
 \item Find $\dfrac{d}{dx}\bigl(g(1)\bigr)$.
\end{enumerate}}
\workingspace

\starquestion{exercise}{b:ex2:xn}{}{Let $\ell$ be the line tangent to the graph of $y=x^n$ at the point $(1,1)$, where $n>1$, as shown. Let $T$ be the triangular region bounded by $\ell$, the $x$-axis, and the line $x=1$. Show that the area of $T$ is $\dfrac{1}{2n}$.}
\begin{center}
\begin{asy}
unitsize(2.4cm);
axes(-0.1,1.45,-0.15,1.55);
path tri=(2/3,0)--(1,0)--(1,1)--cycle;
fill(tri,gray(0.85));
curve(new real(real x){return x^3;},0,1.13);
seg((2/3,0),(1.18,1.54));
draw((1,0)--(1,1.3),linewidth(0.6pt));
closeddot((1,1));
label("$(1,1)$",(1,1),E,fontsize(9pt));
label("$T$",(0.89,0.25),fontsize(9pt));
label("$\ell$",(1.18,1.54),E,fontsize(9pt));
label("$y=x^n$",(0.6,0.216),NW,fontsize(9pt));
label("$O$",(0,0),SW,fontsize(8pt));
\end{asy}
\end{center}
\workingspace

\question{exercise}{b:ex2:order}{}{The graph of a twice-differentiable function $f$ is shown. Which of the following is true?}
\begin{center}
\begin{asy}
unitsize(1.05cm);
pair[] P={(-2.781,2.111),(-2.632,1.822),(-2.483,1.588),(-2.334,1.376),(-2.185,1.193),(-2.036,1.021),(-1.888,0.866),(-1.739,0.724),(-1.590,0.598),(-1.441,0.491),(-1.292,0.404),(-1.143,0.343),(-0.994,0.324),(-0.845,0.342),(-0.696,0.388),(-0.547,0.456),(-0.398,0.532),(-0.249,0.614),(-0.101,0.706),(0.048,0.820),(0.197,0.932),(0.346,1.018),(0.495,1.047),(0.644,1.009),(0.793,0.913),(0.942,0.790),(1.091,0.672),(1.240,0.574),(1.389,0.519),(1.538,0.489),(1.687,0.466),(1.835,0.449),(1.984,0.439),(2.133,0.433),(2.282,0.432),(2.431,0.436),(2.580,0.443),(2.729,0.451),(2.878,0.465),(3.027,0.484)};
guide g; for (pair p: P) g=g..p; draw(g,linewidth(0.9pt));
axes(-2.95,3.3,-0.6,2.6);
xnums(new real[]{-2,-1,1,2});
label("$y=f(x)$",(1.7,0.95),E,fontsize(9pt));
draw((1.7,0.95)..(1.15,0.75),linewidth(0.4pt),EndArrow(3));
label("$O$",(0,0),SW,fontsize(8pt));
\end{asy}
\end{center}
\twochoices{$f'(-1)<f'(1)<f'(0)$}{$f'(-1)<f'(0)<f'(1)$}{$f'(1)<f'(-1)<f'(0)$}{$f'(1)<f'(0)<f'(-1)$}
\workingspace

% Proposed Exercise 2 questions from helper C (Lee chooses).
\question[long]{exercise}{c:xshoppers}{}{\calc A store is having a $12$-hour sale. The total number of shoppers who have entered the store $t$ hours after the sale begins is modeled by the function $S$ defined by $S(t)=0.5t^4-16t^3+144t^2$ for $0\le t\le12$.
\begin{enumerate}[label=(\alph*)]
 \item At what rate are shoppers entering the store $3$ hours after the start of the sale?
 \item Find $\dfrac{d}{dt}\bigl(S(3)\bigr)$.
 \item Find all times $t$, for $0<t<12$, at which the rate at which shoppers enter the store equals the average rate of change of $S$ over $0\le t\le12$.
\end{enumerate}}
\workingspace

% AB 2007 Form B FRQ 3(a): the interpretation of W'(20) is left out (it stays in the question bank, 4.1).
\question[long]{exercise}{c:xwindchill}{}{\calc The wind chill is the temperature, in degrees Fahrenheit $({}^\circ\mathrm{F})$, a human feels based on the air temperature, in degrees Fahrenheit, and the wind velocity $v$, in miles per hour $(\mathrm{mph})$. If the air temperature is $32^\circ\mathrm{F}$, then the wind chill is given by $W(v)=55.6-22.1v^{0.16}$ and is valid for $5\le v\le60$.
\begin{enumerate}[label=(\alph*)]
 \item Find $W'(20)$. Indicate units of measure.
 \item Find the average rate of change of $W$ over the interval $5\le v\le60$. Find the value of $v$ at which the instantaneous rate of change of $W$ is equal to the average rate of change of $W$ over the interval $5\le v\le60$.
\end{enumerate}}
\workingspace

% (a) is AB 2010 Form B FRQ 2(a) (g'' left out of the stem: not needed); (b) is AB 2025 FRQ 2(D), with g renamed h.
\question[long]{exercise}{c:xhoriz}{}{\calc
\begin{enumerate}[label=(\alph*)]
 \item \srctag{AB 2010 Form B FRQ}\enspace The function $g$ is defined for $x>0$ with $g(1)=2$ and $g'(x)=\sin\Bigl(x+\dfrac1x\Bigr)$. Find all values of $x$ in the interval $0.12\le x\le1$ at which the graph of $g$ has a horizontal tangent line.
 \item \srctag{AB 2025 FRQ}\enspace Let $f(x)=x^2-2x$ and $h(x)=x+\sin(\pi x)$. It can be shown that $h'(x)=1+\pi\cos(\pi x)$. Find the value of $x$, for $0<x<1$, at which the line tangent to the graph of $f$ is parallel to the line tangent to the graph of $h$.
\end{enumerate}}
\workingspace
\clearpage
\section{Differentiability and Continuity}\label{sec:diff}
% =====================================================================
% REMOVED AP QUESTIONS (Lee, Nov 2026). These were cut when the subsection "Is a piecewise function
% differentiable?" was folded into 3.2. They are NOT anywhere else: not in this worksheet and not in
% the question bank. Similar questions that stay serve the same purpose (shown after the arrow).
% - BC 2024 MC 15 = AB 2014 MC 23 = Q26, BC 2022 Intl MC 9 = Q95, Q110 (corner: values match, slopes
%   differ)  ->  Example 16(a), P78 (AB 2003 MC 20).
% - AB 1988 MC 27 = U1-Q9 = Q18, AB 2019 Alt MC 15, 2022 Intl MC 1 = Q94 (a join that looks like a
%   corner but is smooth)  ->  Example 16(c), P75(c), P80 (Q88, the join at x = -1).
% - BC 2024 Alt MC 2, BC 2025 MC 5 (three pieces)  ->  P80 (Q88).
% - AB 2013 Intl MC 14 = Q61 (two unknowns, x + b and ax^2)  ->  Example 18(a), P93, BC 2008 MC 25 in Exercise 3.
% (AB 1998 MC 81 = Q1 was moved to the question bank, 5.2, because of "absolute minimum".)
% =====================================================================
% =====================================================================
% CED 2.4. Helper D: 3.1 Where a derivative does not exist, 3.2 Differentiable implies continuous.
% ---------------------------------------------------------------------
\Needspace{20\baselineskip}
\subsection{Differentiability from a graph}
\begin{notesbox}{Notes}
\begin{minipage}{0.62\linewidth}
Graphically, $f'(a)$ exists when the graph is \textbf{smooth} at $a$ and does not have a vertical tangent at $a$.
\end{minipage}\hfill
\begin{minipage}{0.34\linewidth}\centering
\begin{asy}
unitsize(0.5cm);
real f(real x){return 0.6+0.9*x-0.12*x^2;}
real fp(real x){return 0.9-0.24*x;}
axes(-0.3,5.4,-0.3,3.4);
curve(f,0.2,5);
real a=2.2;
draw((a-1.6,f(a)-1.6*fp(a))--(a+1.6,f(a)+1.6*fp(a)), rgb(1,0.5,0)+linewidth(1.1pt));
closeddot((a,f(a))); xlab(a,"$a$");
\end{asy}
\par\small smooth at $a$: one tangent line
\end{minipage}

\medskip
$f'(a)$ does not exist at a corner, a cusp or a vertical tangent, or where $f$ is not continuous at $a$.
\begin{center}
\begin{tabular}{*{3}{>{\centering\arraybackslash}p{0.2\linewidth}}}
\begin{asy}
unitsize(0.42cm);
axes(-0.3,4.2,-0.3,3.3);
seg((0.4,2.8),(2,0.9)); seg((2,0.9),(3.7,2.3));
closeddot((2,0.9)); xlab(2,"$a$");
\end{asy}
&
\begin{asy}
unitsize(0.42cm);
axes(-0.3,4.2,-0.3,3.3);
pcurve(new real(real x){return 0.7+1.25*abs(x-2)^(2/3);}, 0.3, 3.7);
closeddot((2,0.7)); xlab(2,"$a$");
\end{asy}
&
\begin{asy}
unitsize(0.42cm);
axes(-0.3,4.2,-0.3,3.3);
real cr(real t){return t>=0 ? t^(1/3) : -(-t)^(1/3);}
pcurve(new real(real x){return 1.6+0.95*cr(x-2);}, 0.3, 3.7);
closeddot((2,1.6)); xlab(2,"$a$");
\end{asy}
\\
a corner & a cusp & a vertical tangent
\end{tabular}
\end{center}
\end{notesbox}

\Needspace{0.45\textheight}
% (a) new, the graph types of AB 1998 MC 13, AB 2003 MC 13, Q87 and Q190 in one graph;
% (b) new, the method of BC 2018 Alt FRQ 3(c), AB 2005 FRQ 5(b)-(c) and AB 2009 Form B FRQ 3(a).
\exampletitle{d:ex:dne}{Differentiability from a graph}
\begin{enumerate}[label=(\alph*)]
 \item The graph of the function $f$, defined for $-5\le x\le6$, is shown. The graph has a horizontal tangent at $x=-2$ and a vertical tangent at $x=4$.
 \begin{center}
 \begin{asy}
 unitsize(0.62cm);
 real cr(real t){return t>=0 ? t^(1/3) : -(-t)^(1/3);}
 grid(-5,6,-2,5);
 axes(-5.5,6.8,-2.5,5.6);
 xnums(new real[]{-5,-4,-3,-2,-1,1,2,3,4,5,6}); ynums(new real[]{-2,-1,1,2,3,4,5});
 seg((-5,4),(-4,2));
 curve(new real(real x){return 3-0.25*(x+2)^2;}, -4, 1);
 pcurve(new real(real x){return -2+abs(x-2)^(2/3);}, 1, 3);
 pcurve(new real(real x){return 1+2*cr(x-4);}, 3, 6);
 closeddot((-5,4)); opendot((1,0.75)); closeddot((1,-1));
 opendot((5,3)); closeddot((5,4.5)); closeddot((6,1+2*cr(2)));
 label("$O$",(0,0),SW,fontsize(8pt));
 \end{asy}
 \par\vspace{2pt}
 {\small Graph of $f$}
 \end{center}
 \begin{enumerate}[label=(\roman*)]
  \item At which values of $x$, $-5<x<6$, is $f$ not differentiable?
  \item At which of the values in (i) is $f$ continuous?
 \end{enumerate}
 \item The continuous function $f$ is defined on $-4\le x\le4$. Its graph consists of two line segments and a semicircle, as shown.
 \begin{center}
 \begin{asy}
 unitsize(0.62cm);
 grid(-4,4,-2,3);
 axes(-4.5,4.8,-2.5,3.6);
 xnums(new real[]{-4,-3,-2,-1,1,2,3,4}); ynums(new real[]{-2,-1,1,2,3});
 seg((-4,-2),(-2,2)); seg((-2,2),(0,0));
 draw(arc((2,0),2,180,0), linewidth(0.9pt));
 closeddot((-4,-2)); closeddot((-2,2)); closeddot((4,0));
 label("$(-2,2)$",(-2,2),N,fontsize(8pt),Fill(white));
 label("$O$",(0,0),SW,fontsize(8pt));
 \end{asy}
 \par\vspace{2pt}
 {\small Graph of $f$}
 \end{center}
 \begin{enumerate}[label=(\roman*)]
  \item Is $f$ differentiable at $x=-2$? Is $f$ differentiable at $x=0$? Answer yes or no.
  \item Write a piecewise-defined function for $f'(x)$ for $-4<x<0$, $x\ne-2$.
  \item For each of $\displaystyle\lim_{x\to0^-}f'(x)$ and $\displaystyle\lim_{x\to0^+}f'(x)$, find the value or state that it does not exist.
 \end{enumerate}
\end{enumerate}
\end{questionbox}
\begin{examplesolution}
\relax
\begin{enumerate}[label=(\alph*)]
 \item
 \begin{enumerate}[label=(\roman*)]
  \item $f$ is not differentiable at
  \begin{itemize}
   \item $x=-4$: a corner (slope $-2$ on the left, slope $+1$ on the right);
   \item $x=1$: a jump ($f(1)=-1$, but the graph comes in to the open dot at height $0.75$);
   \item $x=2$: a cusp;
   \item $x=4$: a vertical tangent;
   \item $x=5$: a hole ($f(5)=4.5$, but the graph comes in to height $3$ from both sides).
  \end{itemize}
  At $x=-2$ the tangent is horizontal, so $f'(-2)=0$ exists. At $x=3$ the two curves join smoothly.
  \item Continuous at $x=-4$, $2$ and $4$ (corner, cusp, vertical tangent). The jump at $x=1$ and the hole at $x=5$ are discontinuities.
 \end{enumerate}
 \item
 \begin{enumerate}[label=(\roman*)]
  \item At $x=-2$: no, the graph has a corner (slope $2$ on the left, $-1$ on the right).\\
  At $x=0$: no, on the right the semicircle has a vertical tangent at $(0,0)$.
  \item The slopes of the segments: $\dfrac{2-(-2)}{-2-(-4)}=2$ and $\dfrac{0-2}{0-(-2)}=-1$.
  \[
  f'(x)=\begin{cases}2 & \text{for } -4<x<-2\\ -1 & \text{for } -2<x<0\end{cases}
  \]
  \item $\displaystyle\lim_{x\to0^-}f'(x)=-1$, the slope of the segment.

  $\displaystyle\lim_{x\to0^+}f'(x)=\DNEinf$: at $(0,0)$ the semicircle has a vertical tangent, and its slope grows without bound as $x\to0^+$.
 \end{enumerate}
\end{enumerate}
\end{examplesolution}

\practiceheading
% (a)-(b) new: a variation of the example, part for part.
\question[long]{practice}{d:pr:dnevar}{}{%
\begin{enumerate}[label=(\alph*)]
 \item The graph of the function $f$, defined for $-5\le x\le6$, is shown below. The graph has a vertical tangent at $x=-3$ and a horizontal tangent at $x=3$.
 \begin{enumerate}[label=(\roman*)]
  \item At which values of $x$, $-5<x<6$, is $f$ not differentiable?
  \item At which of the values in (i) is $f$ continuous?
 \end{enumerate}
\end{enumerate}}
\begin{center}
\begin{asy}
unitsize(0.62cm);
real cr(real t){return t>=0 ? t^(1/3) : -(-t)^(1/3);}
grid(-5,6,-1,4);
axes(-5.5,6.8,-1.5,4.6);
xnums(new real[]{-5,-4,-3,-2,-1,1,2,3,4,5,6}); ynums(new real[]{-1,1,2,3,4});
pcurve(new real(real x){return 2+cr(x+3);}, -5, 1);
curve(new real(real x){return 2+cr(4)-2+0.5*(x-3)^2;}, 1, 4);
pcurve(new real(real x){return -0.5+abs(x-5)^(2/3);}, 4, 6);
closeddot((-5,2-cr(2))); opendot((-1,2+cr(2))); closeddot((-1,1));
closeddot((4,cr(4)+0.5)); opendot((4,0.5)); closeddot((6,0.5));
label("$O$",(0,0),SW,fontsize(8pt));
\end{asy}
\par\vspace{2pt}
{\small Graph of $f$}
\end{center}
\begin{enumerate}[label=(\alph*),start=2]
 \item The continuous function $f$ is defined on $-5\le x\le3$. Its graph consists of two line segments and a semicircle, as shown below.
 \begin{enumerate}[label=(\roman*)]
  \item Is $f$ differentiable at $x=-3$? Is $f$ differentiable at $x=-1$? Answer yes or no.
  \item Write a piecewise-defined function for $f'(x)$ for $-5<x<-1$, $x\ne-3$.
  \item For each of $\displaystyle\lim_{x\to-1^-}f'(x)$ and $\displaystyle\lim_{x\to-1^+}f'(x)$, find the value or state that it does not exist.
 \end{enumerate}
\end{enumerate}
\begin{center}
\begin{asy}
unitsize(0.62cm);
grid(-5,3,-1,4);
axes(-5.5,3.8,-1.5,4.6);
xnums(new real[]{-5,-4,-3,-2,-1,1,2,3}); ynums(new real[]{-1,1,2,3,4});
seg((-5,0),(-3,3)); seg((-3,3),(-1,2));
draw(arc((1,2),2,180,360), linewidth(0.9pt));
closeddot((-5,0)); closeddot((-3,3)); closeddot((-1,2)); closeddot((3,2));
label("$O$",(0,0),SW,fontsize(8pt));
\end{asy}
\par\vspace{2pt}
{\small Graph of $f$}
\end{center}
\workingspace

% Unit 2 bank Q63 = AB 1998 MC 13 (P). Choice "0, 1, and 3 only" dropped.
\Needspace{20\baselineskip}
\question{practice}{d:pr:vtan}{}{The graph of the function $f$ shown in the figure has a vertical tangent at the point $(2,0)$ and horizontal tangents at the points $(1,-1)$ and $(3,1)$. For what values of $x$, $-2<x<4$, is $f$ not differentiable?}
\begin{center}
\begin{asy}
unitsize(0.75cm);
axes(-2.5,4.6,-2.4,2.6);
xnums(new real[]{-2,-1,1,3,4}); ynums(new real[]{-2,1,2});
draw((2,0),(0,2.5)--(0,-2.5)); label("$2$",(2,0),2SE,fontsize(8pt));
draw((0,-1),(2.5,0)--(-2.5,0)); label("$-1$",(0,-1),2NW,fontsize(8pt));
draw((-2,0.86)..(-1.7,0.21)..(-1.25,-0.32)..(-0.8,-0.67)..(-0.35,-0.9)..(0,-1), linewidth(0.9pt));
draw(arc((1,0),1,180,360), linewidth(0.9pt));
draw(arc((3,0),1,180,0), linewidth(0.9pt));
closeddot((0,0)); opendot((0,-1));
label("$O$",(0,0),SW,fontsize(8pt));
\end{asy}
\end{center}
\fourchoices{$0$ only}{$0$ and $2$ only}{$1$ and $3$ only}{$0$, $1$, $2$, and $3$}
\workingspace

% Unit 2 bank Q14 = AB 2003 MC 13 = BC 2003 MC 13 (P). Choice (E) e dropped.
\Needspace{20\baselineskip}
\question{practice}{d:pr:contnotdiff}{}{The graph of a function $f$ is shown. At which value of $x$ is $f$ continuous, but not differentiable?}
\begin{center}
\begin{asy}
unitsize(0.8cm);
axes(-1.2,5.8,-1.7,3.3);
xlab(1,"$a$"); xlab(2,"$b$"); xlab(3,"$c$"); xlab(5.04,"$e$");
draw((4,0),(0,2.5)--(0,-2.5)); label("$d$",(3.72,-0.3),fontsize(9pt));
draw((-1,-1.6)..(-0.55,-1.14)..(-0.13,-0.53)..(0.21,0.12)..(0.55,1.02)..(0.8,1.95)..(1,2.8), linewidth(0.9pt));
draw((1,2.8)..(1.39,1.80)..(1.73,1.06)..(2,0.52)..(2.23,0.12)..(2.57,-0.31)..(2.91,-0.54)..(3.24,-0.47)..(3.58,-0.19)..(3.83,0.18)..(4,0.55), linewidth(0.9pt));
draw((4,-0.5)..(4.34,-0.48)..(4.67,-0.31)..(5.01,0)..(5.32,0.42), linewidth(0.9pt));
opendot((2,0.52)); closeddot((2,1.65)); closeddot((4,0.55)); opendot((4,-0.5));
label("$O$",(0,0),SW,fontsize(8pt));
\end{asy}
\par\vspace{2pt}
{\small Graph of $f$}
\end{center}
\fourchoices{$a$}{$b$}{$c$}{$d$}
\workingspace

% Unit 2 bank Q190.
\Needspace{20\baselineskip}
\question{practice}{d:pr:semicircles}{}{The figure shows the graph of a continuous piecewise-defined function $f$. The graph consists of a line segment and two semicircles. For which of the following values of $x$ is $f'(x)$ undefined?}
\begin{center}
\begin{asy}
unitsize(0.5cm);
axes(-5,8,-3.5,5);
draw((-1,0),(0,2.5)--(0,-2.5)); draw((1,0),(0,2.5)--(0,-2.5)); draw((3,0),(0,2.5)--(0,-2.5));
label("$r$",(-1,0),2S,fontsize(9pt)); label("$s$",(1,0),3S,fontsize(9pt)); label("$t$",(3,0),2S,fontsize(9pt));
seg((-4,-3),(-1,2));
draw(arc((1,2),2,180,360), linewidth(0.9pt));
draw(arc((5,2),2,180,0), linewidth(0.9pt));
closeddot((-1,2)); closeddot((1,0)); closeddot((3,2));
label("$O$",(0,0),SW,fontsize(8pt));
\end{asy}
\par\vspace{2pt}
{\small Graph of $f$}
\end{center}
\fourchoices{$r$ only}{$t$ only}{$r$ and $s$}{$r$ and $t$}
\workingspace

% Unit 2 bank Q87 (graph redrawn: same features).
\Needspace{24\baselineskip}
\question{practice}{d:pr:falsereason}{}{The figure shows the graph of a function $f$, which has a vertical tangent at $x=2$ and a horizontal tangent at $x=4$. Which of the following statements is false?}
\begin{center}
\begin{asy}
unitsize(0.7cm);
axes(-2.2,8.2,-0.6,4.2);
xnums(new real[]{-1,1,2,3,4,5,6,7});
real cr(real t){return t>=0 ? t^(1/3) : -(-t)^(1/3);}
seg((-2,0.46),(-1,0.46));
pcurve(new real(real x){real v=(2-x)/3; return 2+3*(cr(v)-v);}, -1, 2);
pcurve(new real(real x){real u=(x-2)/2; return 2*(1-cr(u))^2;}, 2, 4);
curve(new real(real x){return 0.71*(x-4)^2;}, 4, 6);
seg((6,2.84),(7.8,2.84));
opendot((-1,0.46)); closeddot((-1,2)); opendot((6,2.84));
label("$O$",(0,0),SW,fontsize(8pt));
\end{asy}
\par\vspace{2pt}
{\small Graph of $f$}
\end{center}
\begin{enumerate}[label=(\Alph*)]
 \item $f$ is not differentiable at $x=-1$ because the graph of $f$ has a jump discontinuity at $x=-1$.
 \item $f$ is not differentiable at $x=2$ because the graph of $f$ has a vertical tangent at $x=2$.
 \item $f$ is not differentiable at $x=4$ because the graph of $f$ has a horizontal tangent at $x=4$.
 \item $f$ is not differentiable at $x=6$ because the graph of $f$ has a removable discontinuity at $x=6$.
\end{enumerate}
\workingspace

% BC 2021 MC 13 = Unit 2 bank Q189 (P).
\Needspace{26\baselineskip}
\question{practice}{d:pr:whichgraph}{}{Let $g$ be a function such that $\displaystyle\lim_{x\to2}g(x)=g(2)$ and $\displaystyle\lim_{x\to2^-}g'(x)\ne\lim_{x\to2^+}g'(x)$. Which of the following could be the graph of $g$?}
\par\noindent\begin{minipage}{\linewidth}
\begin{center}
\begin{minipage}[t]{0.45\linewidth}\centering
(A)\par
\begin{asy}
unitsize(0.5cm);
axes(-0.8,5,-3,2.8);
xnums(new real[]{1,2,3,4}); ynums(new real[]{-2,-1,1,2});
seg((-0.41,-2.82),(2,2)); seg((2,1),(4.72,-1.72));
opendot((2,2)); closeddot((2,1));
\end{asy}
\end{minipage}\hfill
\begin{minipage}[t]{0.45\linewidth}\centering
(B)\par
\begin{asy}
unitsize(0.5cm);
axes(-0.8,5,-3,2.8);
xnums(new real[]{1,2,3,4}); ynums(new real[]{-2,-1,1,2});
seg((-0.41,-2.82),(2,2)); seg((2,2),(4.41,-2.82));
\end{asy}
\end{minipage}\par\medskip
\begin{minipage}[t]{0.45\linewidth}\centering
(C)\par
\begin{asy}
unitsize(0.5cm);
axes(-0.8,5,-3,2.8);
xnums(new real[]{1,2,3,4}); ynums(new real[]{-2,-1,1,2});
curve(new real(real x){return 1+0.5*(x-2)^3;}, 0.08, 3.52);
\end{asy}
\end{minipage}\hfill
\begin{minipage}[t]{0.45\linewidth}\centering
(D)\par
\begin{asy}
unitsize(0.5cm);
axes(-0.8,5,-3,2.8);
xnums(new real[]{1,2,3,4}); ynums(new real[]{-2,-1,1,2});
curve(new real(real x){return 1+0.5*(x-2)^3;}, 0.08, 2);
curve(new real(real x){return 0.5+0.5*(x-2)^3;}, 2, 3.5);
opendot((2,1)); closeddot((2,0.5));
\end{asy}
\end{minipage}
\end{center}
\end{minipage}
\workingspace

% BC 2003 MC 76 = Unit 1 bank Q106 (P). Choice (B) "f(x) is defined for 0 < x < a" dropped.
\Needspace{20\baselineskip}
\question{practice}{d:pr:limfprime}{}{The graph of the function $f$ is shown. Which of the following statements must be false?}
\begin{center}
\begin{asy}
unitsize(1.1cm);
axes(-1.05,3.2,-0.75,2.4);
xlab(1,"$a$");
seg((-0.8,-0.7),(1,2)); seg((1,2),(2.8,-0.7));
opendot((1,2)); closeddot((1,1));
label("$O$",(0,0),SW,fontsize(8pt));
\end{asy}
\par\vspace{2pt}
{\small Graph of $f$}
\end{center}
\twochoices{$f(a)$ exists.}{$f$ is not continuous at $x=a$.}{$\displaystyle\lim_{x\to a}f(x)$ exists.}{$\displaystyle\lim_{x\to a}f'(x)$ exists.}
\workingspace

% AB 2009 Form B FRQ 3(a) (P). The stem's sentence "On 0 < x < 6, f is twice differentiable, with f''(x) > 0" is left out (f'' is Unit 3);
% the justification reads from the graph that f decreases on 0 < x < 3.
\Needspace{22\baselineskip}
\question{practice}{d:pr:onesided}{}{A continuous function $f$ is defined on the closed interval $-4\le x\le6$. The graph of $f$ consists of a line segment and a curve that is tangent to the $x$-axis at $x=3$, as shown in the figure. Is $f$ differentiable at $x=0$? Use the definition of the derivative with one-sided limits to justify your answer.}
\begin{center}
\begin{asy}
unitsize(0.7cm);
axes(-4.6,6.9,-1.2,2.9);
xnums(new real[]{-3,-2,-1,1,2,3,4,5,6}); ynums(new real[]{1,2});
draw((-4,0),(0,2.5)--(0,-2.5)); label("$-4$",(-4,0),2NW,fontsize(8pt));
seg((-4,-2/3),(0,2));
curve(new real(real x){return 2/9*(x-3)^2;}, 0, 3);
curve(new real(real x){return 1/9*(x-3)^2;}, 3, 6);
closeddot((-4,-2/3)); closeddot((0,2)); closeddot((6,1));
label("$(6,1)$",(6,1),N,fontsize(8pt));
label("$O$",(0,0),SW,fontsize(8pt));
\end{asy}
\par\vspace{2pt}
{\small Graph of $f$}
\end{center}
\workingspace

% ---------------------------------------------------------------------
\Needspace{20\baselineskip}
\subsection{Differentiability from a formula}
\begin{notesbox}{Notes}
$f$ is differentiable at $a$ when
\[
\lim_{h\to0}\frac{f(a+h)-f(a)}{h}\ \text{ exists,}
\]
or, the same thing,
\[
\lim_{x\to a}\frac{f(x)-f(a)}{x-a}\ \text{ exists.}
\]
For a piecewise function, $f(a)$ comes from the piece that contains $a$; use the left formula for the left-hand limit and the right formula for the right-hand limit. If the pieces do not meet, one of the two limits blows up, even when the pieces have the same slope.
\end{notesbox}

\Needspace{0.45\textheight}
% (a) new (one-sided difference quotients at a corner, the |x - a| type of AB 1998 MC 81 and Q30);
% (b) new (sqrt|x - a|, the type of AB 2012 Intl MC 11).
\exampletitle{d:ex:dneformula}{Differentiability from a formula}
\begin{enumerate}[label=(\alph*)]
 \item Let $f(x)=\bigl|x^2-4\bigr|$. Is $f$ differentiable at $x=2$? Justify your answer.
 \item Let $g(x)=1+\sqrt{|x-1|}$. Is $g$ continuous at $x=1$? Is $g$ differentiable at $x=1$? Justify your answers.
 \item Let $f(x)=\begin{cases}x^2+1 & \text{for } x\le1\\ 2x & \text{for } x>1.\end{cases}$\\[4pt]
 Is $f$ differentiable at $x=1$? Justify your answer.
 \item \srctag{BC 2014}\enspace The function $f$ is defined by
 \[
 f(x)=\begin{cases} 2x+5 & \text{for } x<-1\\ -x^2+6 & \text{for } x\ge-1.\end{cases}
 \]
 Is $f$ differentiable at $x=-1$? Justify your answer.
\end{enumerate}
\end{questionbox}
\begin{examplesolution}
\relax
\begin{enumerate}[label=(\alph*)]
 \item We need to check whether $\displaystyle\lim_{h\to0}\frac{f(2+h)-f(2)}{h}$ exists, where $f(2)=0$. Split it into the two one-sided limits.
 \twosides{%
 \begin{align*}
 &\lim_{h\to0^-}\frac{f(2+h)-f(2)}{h}\\
 &=\lim_{h\to0^-}\frac{\bigl|(2+h)^2-4\bigr|-0}{h}\\
 &=\lim_{h\to0^-}\frac{\bigl|4h+h^2\bigr|}{h}\\
 &=\lim_{h\to0^-}\frac{\bigl|h(h+4)\bigr|}{h}\\
 &=\lim_{h\to0^-}\frac{-h(h+4)}{h}\\
 &=\lim_{h\to0^-}(-4-h)=-4
 \end{align*}}{%
 \begin{align*}
 &\lim_{h\to0^+}\frac{f(2+h)-f(2)}{h}\\
 &=\lim_{h\to0^+}\frac{\bigl|(2+h)^2-4\bigr|-0}{h}\\
 &=\lim_{h\to0^+}\frac{\bigl|4h+h^2\bigr|}{h}\\
 &=\lim_{h\to0^+}\frac{\bigl|h(h+4)\bigr|}{h}\\
 &=\lim_{h\to0^+}\frac{h(h+4)}{h}\\
 &=\lim_{h\to0^+}(4+h)=4
 \end{align*}}
 (For $h$ close to $0$, $h+4>0$, so $h(h+4)$ has the sign of $h$: negative on the left, so the absolute value changes its sign; positive on the right.)

 $-4\ne4$, so $f'(2)$ does not exist: $f$ is not differentiable at $x=2$. The graph has a corner at $(2,0)$:
 \begin{center}
 \begin{asy}
 unitsize(1.3cm,0.4cm);
 axes(-0.4,3.6,-0.6,4.8);
 xnums(new real[]{1,2,3});
 curve(new real(real x){return abs(x^2-4);}, 0, 2.9);
 pen org=rgb(1,0.5,0)+linewidth(1.1pt);
 draw((1.45,2.2)--(2,0)--(2.55,2.2), org);
 closeddot((2,0));
 label("slope $-4$",(1.55,0.8),W,fontsize(8pt));
 label("slope $4$",(2.45,0.8),E,fontsize(8pt));
 label("$O$",(0,0),SW,fontsize(8pt));
 \end{asy}
 \end{center}
 \item Continuity: $g(1)=1+0=1$ and $\displaystyle\lim_{x\to1}g(x)=1+\sqrt0=1=g(1)$, so $g$ is continuous at $x=1$.

 Differentiability: we need to check whether $\displaystyle\lim_{h\to0}\frac{g(1+h)-g(1)}{h}$ exists. Split it into the two one-sided limits.
 \twosides{%
 \begin{align*}
 &\lim_{h\to0^-}\frac{g(1+h)-g(1)}{h}\\
 &=\lim_{h\to0^-}\frac{1+\sqrt{|h|}-1}{h}=\lim_{h\to0^-}\frac{\sqrt{-h}}{h}\\
 &=\lim_{h\to0^-}\frac{\sqrt{-h}}{-\bigl(\sqrt{-h}\bigr)^2}=\lim_{h\to0^-}\frac{-1}{\sqrt{-h}}\\
 &=\DNEninf
 \end{align*}}{%
 \begin{align*}
 &\lim_{h\to0^+}\frac{g(1+h)-g(1)}{h}\\
 &=\lim_{h\to0^+}\frac{1+\sqrt{|h|}-1}{h}=\lim_{h\to0^+}\frac{\sqrt h}{h}\\
 &=\lim_{h\to0^+}\frac{\sqrt h}{\bigl(\sqrt h\bigr)^2}=\lim_{h\to0^+}\frac{1}{\sqrt h}\\
 &=\DNEinf
 \end{align*}}
 The one-sided limits are not finite, so $g'(1)$ does not exist: $g$ is continuous but not differentiable at $x=1$ (a cusp).
 \item We need to check whether $\displaystyle\lim_{h\to0}\frac{f(1+h)-f(1)}{h}$ exists, where $f(1)=1^2+1=2$ (the piece $x\le1$). Split it into the two one-sided limits. For $h<0$, $1+h<1$: use $x^2+1$. For $h>0$, $1+h>1$: use $2x$.
 \twosides{%
 \begin{align*}
 &\lim_{h\to0^-}\frac{f(1+h)-f(1)}{h}\\
 &=\lim_{h\to0^-}\frac{(1+h)^2+1-2}{h}\\
 &=\lim_{h\to0^-}\frac{2h+h^2}{h}\\
 &=\lim_{h\to0^-}(2+h)=2
 \end{align*}}{%
 \begin{align*}
 &\lim_{h\to0^+}\frac{f(1+h)-f(1)}{h}\\
 &=\lim_{h\to0^+}\frac{2(1+h)-2}{h}\\
 &=\lim_{h\to0^+}\frac{2h}{h}\\
 &=\lim_{h\to0^+}2=2
 \end{align*}}
 Both one-sided limits are $2$, so $f$ is differentiable at $x=1$, and $f'(1)=2$.
 \item $f(-1)=-(-1)^2+6=5$, from $-x^2+6$.
 \begin{align*}
 \lim_{x\to-1^-}\frac{f(x)-f(-1)}{x-(-1)}&=\lim_{x\to-1^-}\frac{(2x+5)-5}{x+1} &&\why{left formula}\\
 &=\lim_{x\to-1^-}\frac{2x}{x+1}=\DNEinf, &&\why{top $\to-2$, bottom $\to0^-$}\\
 \lim_{x\to-1^+}\frac{f(x)-f(-1)}{x-(-1)}&=\lim_{x\to-1^+}\frac{\bigl(-x^2+6\bigr)-5}{x+1}=\lim_{x\to-1^+}(1-x)=2.
 \end{align*}
 The left limit does not exist, so $f$ is not differentiable at $x=-1$. The pieces do not meet ($2(-1)+5=3\ne5$): the graph jumps. (The trap answer $2$ compares only the slopes of the pieces, which are both $2$.)
 \begin{center}
 \begin{asy}
void ocurve(real f(real), real a, real b){guide g; for(int i=0;i<=100;++i){real x=a+(b-a)*i/100; g=g..(x,f(x));} draw(g, rgb(1,0.5,0)+linewidth(0.9pt));}
 unitsize(0.5cm);
 axes(-3.6,2.6,-1.6,6.6);
 xnums(new real[]{-2,-1,1,2}); ynums(new real[]{3,5});
 curve(new real(real x){return 2*x+5;}, -3.2, -1);
 ocurve(new real(real x){return -x^2+6;}, -1, 2.3);
 opendot((-1,3)); closeddot((-1,5));
 \end{asy}
 \end{center}
\end{enumerate}
\end{examplesolution}

\practiceheading
% (a)-(b) new: a variation of the example, part for part.
\question[long]{practice}{d:pr:dnevarf}{}{%
\begin{enumerate}[label=(\alph*)]
 \item Let $f(x)=\bigl|9-x^2\bigr|$. Is $f$ differentiable at $x=-3$? Justify your answer.
 \item Let $g(x)=\sqrt{|x+2|}-3$. Is $g$ continuous at $x=-2$? Is $g$ differentiable at $x=-2$? Justify your answers.
 \item Let $f(x)=\begin{cases}x^2-x & \text{for } x\le3\\ 5x-9 & \text{for } x>3.\end{cases}$\\[4pt]
 Is $f$ differentiable at $x=3$? Justify your answer.
 \item The function $f$ is defined by
 \[
 f(x)=\begin{cases} 3x+1 & \text{for } x<2\\ x^2-x+1 & \text{for } x\ge2.\end{cases}
 \]
 Is $f$ differentiable at $x=2$? Justify your answer.
\end{enumerate}}
\workingspace

% Unit 2 bank Q30. Choice "2" dropped.
\question{practice}{d:pr:absderiv}{}{If $f(x)=2+|x-3|$ for all $x$, then the value of the derivative $f'(x)$ at $x=3$ is}
\fourchoices{$-1$}{$0$}{$1$}{nonexistent}
\workingspace

% AB 2003 MC 20 = Q13 (P). Five choices trimmed to four: "II only" dropped.
\question{practice}{e:pr:ab2003}{}{%
\[
f(x)=\begin{cases} x+2 & \text{if } x\le3\\ 4x-7 & \text{if } x>3\end{cases}
\]
Let $f$ be the function given above. Which of the following statements are true about $f$?
\begin{statements}
 \item $\displaystyle\lim_{x\to3}f(x)$ exists.
 \item $f$ is continuous at $x=3$.
 \item $f$ is differentiable at $x=3$.
\end{statements}}
\fourchoices{None}{I only}{I and II only}{I, II, and III}
\workingspace

% BC 2008 MC 6 = U1-Q22 = Q100 (P). Five choices trimmed to four: "III only" dropped.
\question{practice}{e:pr:removable}{}{%
\[
f(x)=\begin{cases} \dfrac{x^2-4}{x-2} & \text{if } x\ne2\\[8pt] 1 & \text{if } x=2\end{cases}
\]
Let $f$ be the function defined above. Which of the following statements about $f$ are true?
\begin{statements}
 \item $f$ has a limit at $x=2$.
 \item $f$ is continuous at $x=2$.
 \item $f$ is differentiable at $x=2$.
\end{statements}}
\fourchoices{I only}{II only}{I and II only}{I, II, and III}
\workingspace

% Q88 (bank only).
\question{practice}{e:pr:q88}{}{%
\[
f(x)=\begin{cases} \dfrac{2}{x} & \text{for } x<-1\\[6pt] x^2-3 & \text{for } -1\le x\le2\\ 4x-3 & \text{for } x>2\end{cases}
\]
Let $f$ be the function defined above. At what values of $x$, if any, is $f$ not differentiable?}
\twochoices{$x=-1$ only}{$x=2$ only}{$x=-1$ and $x=2$}{$f$ is differentiable for all values of $x$.}
\workingspace

% ---------------------------------------------------------------------
\Needspace{20\baselineskip}
\subsection{Differentiable implies continuous}
\begin{notesbox}{Notes}
\[
f \text{ is differentiable at } a \imp f \text{ is continuous at } a, \text{ so } \lim_{x\to a}f(x)=f(a).
\]
\begin{itemize}
 \item The converse is false: $|x|$ is continuous at $0$ but not differentiable there.
 \item Not continuous at $a$ $\imp$ not differentiable at $a$.
 \item $f'(a)$ existing does not give $\displaystyle\lim_{x\to a}f'(x)=f'(a)$.
\end{itemize}
\[
f \text{ is twice differentiable} \imp f' \text{ is differentiable} \imp f' \text{ is continuous},
\]
so the Intermediate Value Theorem applies to $f'$.
\end{notesbox}

\Needspace{0.45\textheight}
% (a) new, the type of U1-Q11 = AB 2017 Alt MC 90 = Q19 (P) and BC 2008 Form B MC 79 = AB 2018 Practice MC 79 = Q72 (P);
% (b) new, the type of BC 2019 Alt MC 9 = Q86 (P) and BC 2012 MC 83 = AB 2018 Practice MC 83;
% (c) new: a limit where differentiability gives both f'(3) (the definition) and continuity (the other factor); (d) new, the type of AB 2013 Intl FRQ 3(c) = BC 2013 Alt FRQ 3(c) (P) and AB 2022 FRQ 4(b) (P).
\exampletitle{f:ex:diffcont}{Differentiable implies continuous}
\begin{enumerate}[label=(\alph*)]
 \item The function $f$ is differentiable at $x=2$, with $f(2)=5$ and $f'(2)=-3$. Which of the following statements must be true?
 \begin{statements}
  \item $\displaystyle\lim_{x\to2}f(x)=5$
  \item $\displaystyle\lim_{h\to0}\frac{f(2+h)-5}{h}=-3$
  \item $\displaystyle\lim_{x\to2}f'(x)=-3$
 \end{statements}
 \item The function $g$ is continuous at $x=1$. Which of the following statements must be true?
 \begin{statements}
  \item $\displaystyle\lim_{x\to1^-}g(x)=\lim_{x\to1^+}g(x)=g(1)$
  \item $g'(1)$ exists.
  \item If $g$ is not differentiable at $x=1$, then $\displaystyle\lim_{x\to1}\frac{g(x)-g(1)}{x-1}$ does not exist.
 \end{statements}
 \item The function $f$ is differentiable. The table gives values of $f$ and $f'$ at selected values of $x$. Find $\displaystyle\lim_{x\to3}\frac{\bigl(f(x)\bigr)^2-16}{x-3}$, or state that the limit does not exist.
 \begin{center}
 \renewcommand{\arraystretch}{1.3}
 \begin{tabular}{|c|c|c|c|}\hline
 $x$ & $1$ & $3$ & $5$\\ \hline
 $f(x)$ & $-2$ & $4$ & $7$\\ \hline
 $f'(x)$ & $6$ & $-2$ & $1$\\ \hline
 \end{tabular}
 \end{center}
 \item The twice-differentiable function $W$ models the volume of water in a tank at time $t$, where $W(t)$ is measured in liters and $t$ is measured in minutes. The table gives values of $W'(t)$ at selected times.
 \begin{center}
 \renewcommand{\arraystretch}{1.3}
 \begin{tabular}{|c|c|c|c|c|}\hline
 $t$ (minutes) & $0$ & $4$ & $9$ & $15$\\ \hline
 $W'(t)$ (liters per minute) & $1.2$ & $2.5$ & $3.1$ & $1.8$\\ \hline
 \end{tabular}
 \end{center}
 \begin{enumerate}[label=(\roman*)]
  \item Explain why there must be at least one time $t$, other than $t=4$, such that $W'(t)=2.5$ liters per minute.
  \item Is there a time $t$, $0\le t\le4$, for which $W'(t)=2$? Justify your answer.
 \end{enumerate}
\end{enumerate}
\end{questionbox}
\begin{examplesolution}
\relax
\begin{enumerate}[label=(\alph*)]
 \item I and II only.
 \begin{itemize}
  \item I: $f$ is differentiable at $2$ $\imp$ $f$ is continuous at $2$ $\imp$ $\displaystyle\lim_{x\to2}f(x)=f(2)=5$. True.
  \item II: $5=f(2)$, so this is the definition of $f'(2)$. True.
  \item III: this says $f'$ is continuous at $2$, which is not given. Could be false.
 \end{itemize}
 \item I and III only.
 \begin{itemize}
  \item I: this is what ``continuous at $1$'' means. True.
  \item II: continuous does not give differentiable: $g(x)=|x-1|$ is continuous at $1$, but $g'(1)$ does not exist. Could be false.
  \item III: ``not differentiable at $1$'' means exactly that this limit, the definition of $g'(1)$, does not exist. True.
 \end{itemize}
 \item $f(3)=4$, so $\bigl(f(x)\bigr)^2-16$ is a difference of squares with $16=\bigl(f(3)\bigr)^2$. Factor it:
 \begin{align*}
 \lim_{x\to3}\frac{\bigl(f(x)\bigr)^2-16}{x-3}&=\lim_{x\to3}\frac{\bigl(f(x)-4\bigr)\bigl(f(x)+4\bigr)}{x-3}\\
 &=\lim_{x\to3}\frac{f(x)-f(3)}{x-3}\cdot\lim_{x\to3}\bigl(f(x)+4\bigr)\\
 &=f'(3)\cdot\bigl(f(3)+4\bigr) &&\why{definition of $f'(3)$; $f$ is continuous at $3$}\\
 &=(-2)(4+4)=-16.
 \end{align*}
 The second limit uses continuity: $f$ is differentiable at $3$, so $\displaystyle\lim_{x\to3}f(x)=f(3)=4$.
 \item
 \begin{enumerate}[label=(\roman*)]
  \item $W$ is twice differentiable $\imp$ $W'$ is differentiable $\imp$ $W'$ is continuous on $[9,15]$.
  \[
  W'(15)=1.8<2.5<3.1=W'(9).
  \]
  By the Intermediate Value Theorem there is a time $t$ in $(9,15)$ with $W'(t)=2.5$; this $t$ is not $4$. (The interval $[0,4]$ gives only $t=4$ itself.)
\begin{center}
\begin{asy}
unitsize(0.4cm,0.85cm);
axes(-0.8,17,-0.35,3.8,"$t$","$W'(t)$");
xnums(new real[]{4,9,15});
ynums(new real[]{1,2,3});
draw((0,2.5)--(16,2.5), dashed+gray(0.4));
label("$2.5$",(16,2.5),E,fontsize(8pt));
pair[] P={(0,1.2),(4,2.5),(9,3.1),(15,1.8)};
draw(P[0]..P[1]..P[2]..P[3], linewidth(0.9pt));
for (pair p: P) closeddot(p);
\end{asy}
\par\vspace{2pt}{\small One possible graph of $W'$: it is continuous, so it crosses the line $2.5$ between $t=9$ and $t=15$.}
\end{center}
  \item Yes. $W'$ is continuous on $[0,4]$ (as in (i)), and
  \[
  W'(0)=1.2<2<2.5=W'(4).
  \]
  By the Intermediate Value Theorem there is a time $t$ in $(0,4)$ with $W'(t)=2$.
 \end{enumerate}
\end{enumerate}
\end{examplesolution}

\practiceheading
% (a)-(d) new: a variation of the example, part for part.
\question[long]{practice}{f:pr:diffcontvar}{}{%
\begin{enumerate}[label=(\alph*)]
 \item The function $f$ is differentiable at $x=-1$, with $f(-1)=2$ and $f'(-1)=4$. Which of the following statements must be true?
 \begin{statements}
  \item $\displaystyle\lim_{x\to-1^+}f(x)=2$
  \item $\displaystyle\lim_{x\to-1}\frac{f(x)-2}{x+1}=4$
  \item $\displaystyle\lim_{x\to-1}f'(x)=4$
 \end{statements}
 \item The function $g$ is continuous at $x=-2$. Which of the following statements must be true?
 \begin{statements}
  \item $\displaystyle\lim_{x\to-2}g(x)$ exists.
  \item $g'(-2)$ exists.
  \item If $g'(-2)$ does not exist, then $\displaystyle\lim_{h\to0}\frac{g(-2+h)-g(-2)}{h}$ does not exist.
 \end{statements}
 \item The function $f$ is differentiable. The table gives values of $f$ and $f'$ at selected values of $x$. Find $\displaystyle\lim_{x\to2}\frac{\bigl(f(x)\bigr)^2-1}{x-2}$, or state that the limit does not exist.
 \begin{center}
 \renewcommand{\arraystretch}{1.3}
 \begin{tabular}{|c|c|c|c|}\hline
 $x$ & $0$ & $2$ & $4$\\ \hline
 $f(x)$ & $3$ & $-1$ & $5$\\ \hline
 $f'(x)$ & $-4$ & $2$ & $3$\\ \hline
 \end{tabular}
 \end{center}
 \item The twice-differentiable function $P$ models the population of a town, where $P(t)$ is measured in people and $t$ is measured in years. The table gives values of $P'(t)$ at selected times.
 \begin{center}
 \renewcommand{\arraystretch}{1.3}
 \begin{tabular}{|c|c|c|c|c|}\hline
 $t$ (years) & $0$ & $2$ & $5$ & $8$\\ \hline
 $P'(t)$ (people per year) & $40$ & $65$ & $90$ & $55$\\ \hline
 \end{tabular}
 \end{center}
 \begin{enumerate}[label=(\roman*)]
  \item Explain why there must be at least one time $t$, other than $t=2$, such that $P'(t)=65$ people per year.
  \item Is there a time $t$, $0\le t\le2$, for which $P'(t)=50$? Justify your answer.
 \end{enumerate}
\end{enumerate}}
\workingspace

% Unit 1 bank Q11 = AB 2017 Alt MC 90 = Unit 2 bank Q19 (P).
\question{practice}{f:pr:mustbe}{}{For any function $f$, which of the following statements must be true?
\begin{statements}
\item If $f$ is defined at $x=a$, then $\displaystyle\lim_{x\to a}f(x)=f(a)$.
\item If $f$ is continuous at $x=a$, then $\displaystyle\lim_{x\to a}f(x)=f(a)$.
\item If $f$ is differentiable at $x=a$, then $\displaystyle\lim_{x\to a}f(x)=f(a)$.
\end{statements}}
\fourchoices{III only}{I and II only}{II and III only}{I, II, and III}
\workingspace

% BC 2008 Form B MC 79 = AB 2018 Practice MC 79 = Unit 2 bank Q72 (P). Choice (A) "lim f(x) = 8" dropped.
\question{practice}{f:pr:couldfalse}{}{If $f$ is a differentiable function such that $f(3)=8$ and $f'(3)=5$, which of the following statements could be false?}
\twochoices{$\displaystyle\lim_{x\to3^+}f(x)=\lim_{x\to3^-}f(x)$}{$\displaystyle\lim_{x\to3}\frac{f(x)-8}{x-3}=5$}{$\displaystyle\lim_{h\to0}\frac{f(3+h)-8}{h}=5$}{$\displaystyle\lim_{x\to3}f'(x)=5$}
\workingspace

% BC 2019 Alt MC 9 = Unit 2 bank Q86 (P). Choice (D) "the integral of f from 0 to 5 does not exist" (Unit 6)
% replaced by "The graph of f has a corner at x = 5." (Lee to confirm).
\question{practice}{f:pr:notdiff}{}{The function $f$ is not differentiable at $x=5$. Which of the following statements must be true?}
\twochoices{$f$ is not continuous at $x=5$.}{$\displaystyle\lim_{x\to5}f(x)$ does not exist.}{$\displaystyle\lim_{x\to5}\frac{f(x)-f(5)}{x-5}$ does not exist.}{The graph of $f$ has a corner at $x=5$.}
\workingspace

% BC 2012 MC 83 = AB 2018 Practice MC 83. Choice (A) "f(3) < lim f(x)" dropped.
\question{practice}{f:pr:cont3}{}{If the function $f$ is continuous at $x=3$, which of the following must be true?}
\begin{enumerate}[label=(\Alph*)]
 \item $\displaystyle\lim_{x\to3^-}f(x)\ne\lim_{x\to3^+}f(x)$
 \item $f(3)=\displaystyle\lim_{x\to3^-}f(x)=\lim_{x\to3^+}f(x)$
 \item The derivative of $f$ at $x=3$ exists.
 \item The derivative of $f$ is positive for $x<3$ and negative for $x>3$.
\end{enumerate}
\workingspace

% Unit 2 bank Q89.
\question{practice}{f:pr:defined}{}{Which of the following statements about the function $f$, if true, cannot be used to conclude that $f$ is defined at $x=1$?}
\twochoices{$\displaystyle\lim_{x\to1}f(x)$ exists.}{$f$ is continuous at $x=1$.}{$f$ is differentiable at $x=1$.}{The line tangent to the graph of $f$ at $x=1$ exists.}
\workingspace

% BC 2017 Alt MC 6.
\Needspace{24\baselineskip}
\question{practice}{f:pr:graphhole}{}{The graph of $y=f(x)$ is shown. Which of the following is true?}
\begin{center}
\begin{asy}
unitsize(0.6cm);
axes(-0.5,6.6,-3,3.1);
xlab(4,"$a$");
curve(new real(real x){return 2.4*tanh((x-3)/1.2);}, -0.3, 6.2);
opendot((4,2.4*tanh(1/1.2))); closeddot((4,0.6));
label("$O$",(0,0),SW,fontsize(8pt));
\end{asy}
\par\vspace{2pt}
{\small Graph of $f$}
\end{center}
\twochoices{$\displaystyle\lim_{h\to0}\frac{f(a+h)-f(a)}{h}$ exists.}{$\displaystyle\lim_{x\to a^+}f(x)\ne\lim_{x\to a^-}f(x)$}{$\displaystyle\lim_{x\to a}f(x)\ne f(a)$}{$\displaystyle\lim_{x\to a}f(x)$ does not exist.}
\workingspace

% AB 2026 FRQ 6(A). The stem's "twice differentiable" kept (nothing here needs f'').
\question{practice}{f:pr:limtable}{}{The function $f$ is twice differentiable. The table gives values of $f$ and its derivative $f'$ at selected values of $x$. Find $\displaystyle\lim_{x\to2}\frac{f(x)}{x}$, or state that the limit does not exist.
\begin{center}
\renewcommand{\arraystretch}{1.3}
\begin{tabular}{|c|c|c|c|c|}\hline
$x$ & $0$ & $2$ & $3$ & $6$\\ \hline
$f(x)$ & $-1$ & $3$ & $8$ & $5$\\ \hline
$f'(x)$ & $-5$ & $4$ & $9$ & $-2$\\ \hline
\end{tabular}
\end{center}}
\workingspace

% ---------------------------------------------------------------------
% Helper E: 3.3 Is a piecewise function differentiable?
% Unit 2 bank Q91 (a jump, not a corner).
\question{practice}{d:pr:jumppiece}{}{Let $f$ be the function defined below. Which of the following statements is true?
\[
f(x)=\begin{cases} x^2-20 & \text{for } x<5\\ -x^2+20 & \text{for } x\ge5\end{cases}
\]}
\begin{enumerate}[label=(\Alph*)]
 \item $f$ is not differentiable at $x=5$ because $f$ is not continuous at $x=5$.
 \item $f$ is not differentiable at $x=5$ because the graph of $f$ has a sharp corner at $x=5$.
 \item $f$ is not differentiable at $x=5$ because the graph of $f$ has a vertical tangent at $x=5$.
 \item $f$ is not differentiable at $x=5$ because $f$ is not defined at $x=5$.
\end{enumerate}
\workingspace

% ---------------------------------------------------------------------
% Helper G: 3.4 Making a piecewise function differentiable (built from helper E's draft).
% ---------------------------------------------------------------------
\Needspace{20\baselineskip}
\subsection{Making a piecewise function differentiable}
\begin{notesbox}{Notes}
For $f(x)=\begin{cases} g(x), & x<a\\ k, & x=a\\ h(x), & x>a,\end{cases}$ there are two ways to check that $f$ is differentiable at $x=a$ (and to find unknown constants).

\medskip
\textbf{Method 1: match the value and the slope.} Both are needed:
\begin{align*}
g(a)&=k=h(a) &&\why{the value matches: $f$ is continuous at $a$}\\
g'(a)&=h'(a) &&\why{the slopes match}
\end{align*}
Then $f'(a)=g'(a)=h'(a)$. The slope equation often has only one unknown: solve it first.

\medskip
\textbf{Method 2: the definition.} The two one-sided limits must be finite and equal:
\[
\lim_{x\to a^-}\frac{g(x)-k}{x-a}=\lim_{x\to a^+}\frac{h(x)-k}{x-a}.
\]
\begin{center}
\begin{asy}
unitsize(0.5cm);
axes(-0.4,6.4,-0.4,4.4);
xlab(3,"$a$");
draw((0.6,0.6)--(5.4,3.0), dashed+linewidth(0.5pt));
curve(new real(real x){return 1.8+0.5*(x-3)-0.15*(x-3)^2;}, 1.0, 3);
{guide g; for (int i=0; i<=100; ++i) {real x=3+(5.3-(3))*i/100; g=g..(x,1.8+0.5*(x-3)+0.2*(x-3)^2);} draw(g, rgb(1,0.5,0)+linewidth(0.9pt));}
closeddot((3,1.8));
\end{asy}
\par
same point, same tangent line
\end{center}
\end{notesbox}

\Needspace{0.45\textheight}
% Parts and sources: (a) new (Lee: the general three-case form, solved by both methods);
% (b) new (cumulative: e^x and cos x); (c) BC 2023 MC 16 = Q96 (P), asked without the letters.
% Removed on Lee's request: AB 2013 Intl MC 14 = Q61 (covered by (a)); AB 2006 Form B FRQ 3 (the ramp) moved to Example 9 (Section 2). FRQ part (c) needs "increasing" (Unit 5; already in the bank)
% and is left out; in (d) the request "and (iii)" (increasing) is dropped for the same reason.
% Requirement 3 stays in the stem as in the source (nothing asks about it).
% The requirements are numbered 1, 2, 3 (the source's (i), (ii), (iii)) so that they are not
% confused with the subparts (i), (ii), (iii).
% AB 2003 FRQ 6(c) (k sqrt(x+1)) needs the chain rule: banked (Unit 3); g:pr:ksqrt is its Unit 2 version.
\exampletitle{g:ex:makediff}{Matching the value and the slope}
\begin{enumerate}[label=(\alph*)]
 \item Let $f(x)=\begin{cases} ax^2+3 & \text{if } x<1\\ k & \text{if } x=1\\ bx+1 & \text{if } x>1.\end{cases}$ Find the values of $a$, $b$ and $k$ for which $f$ is differentiable at $x=1$.
 \item Let $f(x)=\begin{cases} ae^x+b & \text{if } x\le0\\ 2\cos x+3x & \text{if } x>0.\end{cases}$ Find the values of $a$ and $b$ for which $f$ is differentiable at $x=0$.
 \item Let $f(x)=\begin{cases} x^2-2x, & x\le4\\ g(x), & x>4.\end{cases}$ If $f$ is differentiable, which of the following could define $g(x)$: $\ 6x$,\ \ $x^{3/2}$,\ \ $x^2-8$,\ \ $\dfrac18x^3$?
\end{enumerate}
\end{questionbox}
\begin{examplesolution}
\relax
\begin{enumerate}[label=(\alph*)]
 \item \textbf{Method 1.} Match the value and the slope at $x=1$.
 \begin{align*}
 a(1)^2+3&=k=b(1)+1 &&\why{the value: $a+3=k$ and $b+1=k$}\\
 2a(1)&=b &&\why{the slopes: $g'(x)=2ax$, $h'(x)=b$}
 \end{align*}
 So $a+3=b+1=2a+1$, which gives $a=2$, then $b=4$ and $k=5$.

 \textbf{Method 2.} Use the definition with $f(1)=k$.
 \twosides{%
 \begin{align*}
 &\lim_{x\to1^-}\frac{ax^2+3-k}{x-1}
 \end{align*}
 The bottom tends to $0$, so the top must too: $a+3-k=0$, $k=a+3$.
 \begin{align*}
 &=\lim_{x\to1^-}\frac{ax^2-a}{x-1}\\
 &=\lim_{x\to1^-}\frac{a(x-1)(x+1)}{x-1}=2a
 \end{align*}}{%
 \begin{align*}
 &\lim_{x\to1^+}\frac{bx+1-k}{x-1}
 \end{align*}
 The top must tend to $0$: $b+1-k=0$, $k=b+1$.
 \begin{align*}
 &=\lim_{x\to1^+}\frac{bx-b}{x-1}\\
 &=\lim_{x\to1^+}\frac{b(x-1)}{x-1}=b
 \end{align*}}
 Equal: $2a=b$. These are the same three equations, so again $a=2$, $b=4$, $k=5$.
 \item Slopes first: they give $a$ at once.
 \begin{align*}
 \lim_{x\to0^-}f'(x)&=\lim_{x\to0^-}ae^x=a,\qquad \lim_{x\to0^+}f'(x)=\lim_{x\to0^+}(-2\sin x+3)=3,\\
 a&=3, &&\why{the slopes match}\\
 ae^0+b&=2\cos0+3(0) &&\why{the values match}\\
 3+b&=2\imp b=-1.
 \end{align*}
 So $a=3$ and $b=-1$.
 \item At $x=4$ the left piece has value and slope
 \[
 4^2-2(4)=8,\qquad \frac{d}{dx}\bigl(x^2-2x\bigr)=2x-2=6.
 \]
 So $g$ needs $g(4)=8$ and $g'(4)=6$. Check each choice:
 \begin{center}
 \renewcommand{\arraystretch}{1.9}
 \begin{tabular}{|c|c|c|c|}\hline
 $g(x)$ & $g(4)$ & $g'(x)$ & $g'(4)$\\ \hline
 $6x$ & $24$ & $6$ & $6$\\ \hline
 $x^{3/2}$ & $8$ & $\dfrac32x^{1/2}$ & $3$\\ \hline
 $x^2-8$ & $8$ & $2x$ & $8$\\ \hline
 $\dfrac18x^3$ & $8$ & $\dfrac38x^2$ & $6$\\ \hline
 \end{tabular}
 \end{center}
 Only $g(x)=\dfrac18x^3$ matches both. ($6x$ matches only the slope; $x^{3/2}$ and $x^2-8$ match only the value.)
\end{enumerate}
\end{examplesolution}

\practiceheading
% (a)-(d) new: a variation of the example, part for part ((d) is a ramp with new numbers).
% Part (d) is typeset after the question argument so that its graph is not inside a macro argument.
\question[long]{practice}{g:pr:makevar}{}{%
\begin{enumerate}[label=(\alph*)]
 \item Let $f(x)=\begin{cases} ax^2-1 & \text{if } x<2\\ k & \text{if } x=2\\ bx+3 & \text{if } x>2.\end{cases}$ Find the values of $a$, $b$ and $k$ for which $f$ is differentiable at $x=2$.
 \item Let $f(x)=\begin{cases} a\sin x+b & \text{if } x\le\pi\\ x+\cos x & \text{if } x>\pi.\end{cases}$ Find the values of $a$ and $b$ for which $f$ is differentiable at $x=\pi$.
 \item Let $f(x)=\begin{cases} x^2+3, & x\le1\\ g(x), & x>1.\end{cases}$ If $f$ is differentiable, which of the following could define $g(x)$: $\ 4x$,\ \ $x^3+3$,\ \ $4\sqrt x$,\ \ $2x^2+2$?
\end{enumerate}}
\workingspace

% New: the type of AB 2003 FRQ 6(c) (both equations contain both unknowns) with sqrt(x) in place of
% sqrt(x+1); the source needs the chain rule and is banked (Unit 3).
\question{practice}{g:pr:ksqrt}{}{Let $g$ be the function defined by
\[
g(x)=\begin{cases} k\sqrt x & \text{for } 0\le x\le4\\ mx+2 & \text{for } 4<x\le6,\end{cases}
\]
where $k$ and $m$ are constants. If $g$ is differentiable at $x=4$, what are the values of $k$ and $m$?}
\workingspace

% ---------------------------------------------------------------------
\Needspace{\ifshowspace 0.5\textheight\else 16\baselineskip\fi}
\subsection*{Exercise \thesection}\addcontentsline{toc}{subsection}{Exercise \thesection}
% Exercise 3 (Lee): mixed questions from 3.1-3.4, in subsection order; the past-paper ones were moved here from the practice.
% BC 2018 Alt FRQ 3(c) (P). Only part (c) (the stem's f, an integral, belongs to part (d), banked in Unit 6).
\Needspace{24\baselineskip}
\question{exercise}{d:pr:gsemicircle}{}{The function $g$ is defined on the closed interval $[-4,8]$. The graph of $g$ consists of two linear pieces and a semicircle, as shown in the figure. For each of $\displaystyle\lim_{x\to0^-}g'(x)$ and $\displaystyle\lim_{x\to0^+}g'(x)$, find the value or state that it does not exist.}
\begin{center}
\begin{asy}
unitsize(0.55cm);
grid(-4,8,-3,3);
axes(-4.6,8.8,-3.6,3.9);
xnums(new real[]{-4,-3,-2,-1,1,2,3,4,5,7,8}); ynums(new real[]{-3,-2,-1,1,2,3});
label("$6$",(6,0),NE,fontsize(8pt),Fill(white));
seg((-4,2),(0,0));
draw(arc((3,0),3,180,360), linewidth(0.9pt));
seg((6,3),(8,3));
closeddot((-4,2)); closeddot((6,0)); opendot((6,3)); closeddot((8,3));
label("$O$",(0,0),SW,fontsize(8pt),Fill(white));
\end{asy}
\par\vspace{2pt}
{\small Graph of $g$}
\end{center}
\workingspace

% AB 2005 FRQ 5(b), and the Unit 2 part of 5(c): the source asks for the acceleration a(t);
% here it asks for v'(t) (the acceleration version is banked, Unit 4, 4.2).
\Needspace{22\baselineskip}
\question{exercise}{d:pr:velocity}{}{A car is traveling on a straight road. For $0\le t\le24$ seconds, the car's velocity $v(t)$, in meters per second, is modeled by the piecewise-linear function defined by the graph.}
\begin{center}
\begin{asy}
unitsize(0.3cm,0.16cm);
axes(-1,27,-2,25,"$t$","$v(t)$");
xnums(new real[]{4,8,12,16,20,24}); ynums(new real[]{5,10,15,20});
draw((0,0)--(4,20)--(16,20)--(24,0), linewidth(0.9pt));
closeddot((4,20)); closeddot((16,20));
label("$(4,20)$",(4,20),N,fontsize(8pt)); label("$(16,20)$",(16,20),N,fontsize(8pt));
label("$O$",(0,0),SW,fontsize(8pt));
label("Time (seconds)",(12,-6),fontsize(8pt));
\end{asy}
\end{center}
\begin{enumerate}[label=(\alph*)]
 \item For each of $v'(4)$ and $v'(20)$, find the value or explain why it does not exist. Indicate units of measure.
 \item For $0<t<24$, write a piecewise-defined function for $v'(t)$.
\end{enumerate}
\workingspace

% Unit 2 bank Q98 (the symmetric difference quotient at a corner). Choice (A) "lim (f(x) - f(0)) = 0" dropped.
\Needspace{20\baselineskip}
\question{exercise}{d:pr:symdq}{}{The graph of the even function $y=f(x)$ consists of $4$ line segments, as shown. Which of the following statements about $f$ is false?}
\begin{center}
\begin{asy}
unitsize(0.45cm);
axes(-8.4,9,-3,4.5);
xnums(new real[]{-5,-3,3,5}); ynums(new real[]{-2,-1,1,2,3});
draw((-8,-2)--(-3,3)--(0,0)--(3,3)--(8,-2), linewidth(0.9pt));
closeddot((-3,3)); closeddot((3,3));
label("$(-3,3)$",(-3,3),N,fontsize(8pt)); label("$(3,3)$",(3,3),N,fontsize(8pt));
label("$O$",(0,0),SW,fontsize(8pt));
\end{asy}
\end{center}
\twochoices{$\displaystyle\lim_{x\to0}\frac{f(x)-f(0)}{x}=0$}{$\displaystyle\lim_{x\to0}\frac{f(x)-f(-x)}{2x}=0$}{$\displaystyle\lim_{x\to2}\frac{f(x)-f(2)}{x-2}=1$}{$\displaystyle\lim_{x\to3}\frac{f(x)-f(3)}{x-3}$ does not exist.}
\workingspace

% Proposed Exercise 3 questions from helper D (new; Lee chooses).
\starquestion[long]{exercise}{d:ex3:symq}{}{Let $f(x)=x\,|x-2|$.
\begin{enumerate}[label=(\alph*)]
 \item Is $f$ continuous at $x=2$? Justify your answer.
 \item Use one-sided limits of the difference quotient to show that $f$ is not differentiable at $x=2$.
 \item Evaluate $\displaystyle\lim_{h\to0}\frac{f(2+h)-f(2-h)}{2h}$. Explain why your answer does not contradict (b).
\end{enumerate}}
\workingspace

\question[long]{exercise}{d:ex3:powers}{}{For each function, use one-sided limits of the difference quotient to decide whether the function is differentiable at $x=0$. If it is not, say whether the graph has a corner, a cusp or a vertical tangent at $x=0$.
\begin{enumerate}[label=(\alph*)]
 \item $f(x)=\sqrt[3]{x}$
 \item $g(x)=\sqrt[3]{x^2}$
 \item $k(x)=x\,|x|$
\end{enumerate}}
\workingspace

% AB 2012 Intl MC 11 = BC 2012 MC 11 = Unit 2 bank Q10 (P). Choice (D) "lim f(x) is not 0" dropped.
\question{exercise}{d:pr:rootabs}{}{Let $f$ be the function defined by $f(x)=\sqrt{|x-2|}$ for all $x$. Which of the following statements is true?}
\begin{enumerate}[label=(\Alph*)]
 \item $f$ is continuous but not differentiable at $x=2$.
 \item $f$ is differentiable at $x=2$.
 \item $f$ is not continuous at $x=2$.
 \item $x=2$ is a vertical asymptote of the graph of $f$.
\end{enumerate}
\workingspace

% AB 2002 FRQ 6(d) (P). The stem's last sentence, "The second derivative of f has the property that
% f''(x) > 0 for -1.5 <= x <= 1.5", is dropped (f'' is Unit 3; part (d) does not use it).
\question{exercise}{e:pr:samefn}{}{Let $f$ be a function that is differentiable for all real numbers. The table below gives the values of $f$ and its derivative $f'$ for selected points $x$ in the closed interval $-1.5\le x\le1.5$.
\begin{center}
\renewcommand{\arraystretch}{1.3}
\begin{tabular}{|c|c|c|c|c|c|c|c|}\hline
$x$ & $-1.5$ & $-1.0$ & $-0.5$ & $0$ & $0.5$ & $1.0$ & $1.5$\\ \hline
$f(x)$ & $-1$ & $-4$ & $-6$ & $-7$ & $-6$ & $-4$ & $-1$\\ \hline
$f'(x)$ & $-7$ & $-5$ & $-3$ & $0$ & $3$ & $5$ & $7$\\ \hline
\end{tabular}
\end{center}
Let $g$ be the function given by $g(x)=\begin{cases} 2x^2-x-7 & \text{for } x<0\\ 2x^2+x-7 & \text{for } x\ge0.\end{cases}$

The graph of $g$ passes through each of the points $\bigl(x,f(x)\bigr)$ given in the table above. Is it possible that $f$ and $g$ are the same function? Give a reason for your answer.}
\workingspace

% Helper F: PROPOSED Exercise 3 questions (from 3.2). Lee chooses among them.
% New: continuity from differentiability inside a limit (not 0/0) next to a derivative limit (Section 2),
% then the IVT on f (f continuous because differentiable) and the IVT on f' (f' continuous because twice differentiable).
\starquestion[long]{exercise}{f:ex3:table}{}{The function $f$ is twice differentiable. The table gives values of $f$ and its derivative $f'$ at selected values of $x$.
\begin{center}
\renewcommand{\arraystretch}{1.3}
\begin{tabular}{|c|c|c|c|c|}\hline
$x$ & $1$ & $2$ & $4$ & $6$\\ \hline
$f(x)$ & $3$ & $5$ & $5$ & $2$\\ \hline
$f'(x)$ & $4$ & $-1$ & $2$ & $-3$\\ \hline
\end{tabular}
\end{center}
\begin{enumerate}[label=(\alph*)]
 \item Find $\displaystyle\lim_{x\to4}\frac{f(x)-5}{x-2}$, or state that the limit does not exist.
 \item Find $\displaystyle\lim_{x\to4}\frac{f(x)-5}{x-4}$, or state that the limit does not exist.
 \item Must there be a value $c$, $1<c<2$, for which $f(c)=4$? Justify your answer.
 \item Explain why there must be at least three values of $c$, $1<c<6$, for which $f'(c)=0$.
\end{enumerate}}
\workingspace

% New: the proof in the notes run backwards: a limit of the form (f(x)-5)/(x-1) that exists forces lim f(x) = 5,
% but not f(1) = 5; continuity then makes it the derivative (Section 2, limits that are derivatives).
\question{exercise}{f:ex3:limexists}{}{The function $f$ is defined for all $x$, and $\displaystyle\lim_{x\to1}\frac{f(x)-5}{x-1}=3$.
\begin{enumerate}[label=(\alph*)]
 \item Which of the following statements must be true?
 \begin{statements}
  \item $\displaystyle\lim_{x\to1}f(x)=5$
  \item $f(1)=5$
  \item $f'(1)=3$
 \end{statements}
 \fourchoices{I only}{I and II only}{I and III only}{I, II, and III}
 \item Suppose, in addition, that $f$ is continuous at $x=1$. Find $f(1)$ and $f'(1)$.
\end{enumerate}}
\workingspace

% AB 2013 Intl FRQ 3(c) = BC 2013 Alt FRQ 3(c) (P). The stem keeps "twice-differentiable" (it gives W' continuous).
\question{exercise}{f:pr:ivtW}{}{The twice-differentiable function $W$ models the volume of water in a reservoir at time $t$, where $W(t)$ is measured in gigaliters (GL) and $t$ is measured in days. The table gives values of $W'(t)$ sampled at various times during the time interval $0\le t\le30$ days. At time $t=30$, the reservoir contains $125$ gigaliters of water. Explain why there must be at least one time $t$, other than $t=10$, such that $W'(t)=0.7$ GL/day.
\begin{center}
\renewcommand{\arraystretch}{1.3}
\begin{tabular}{|c|c|c|c|c|}\hline
$t$ (days) & $0$ & $10$ & $22$ & $30$\\ \hline
$W'(t)$ (GL per day) & $0.6$ & $0.7$ & $1.0$ & $0.5$\\ \hline
\end{tabular}
\end{center}}
\workingspace

% AB 2022 FRQ 4(b) (P).
\question{exercise}{f:pr:ivtr}{}{An ice sculpture melts in such a way that it can be modeled as a cone that maintains a conical shape as it decreases in size. The radius of the base of the cone is given by a twice-differentiable function $r$, where $r(t)$ is measured in centimeters and $t$ is measured in days. The table gives selected values of $r'(t)$, the rate of change of the radius, over the time interval $0\le t\le12$. Is there a time $t$, $0\le t\le3$, for which $r'(t)=-6$? Justify your answer.
\begin{center}
\renewcommand{\arraystretch}{1.3}
\begin{tabular}{|c|c|c|c|c|c|}\hline
$t$ (days) & $0$ & $3$ & $7$ & $10$ & $12$\\ \hline
$r'(t)$ (centimeters per day) & $-6.1$ & $-5.0$ & $-4.4$ & $-3.8$ & $-3.5$\\ \hline
\end{tabular}
\end{center}}
\workingspace

% Helper E: PROPOSED Exercise 3 questions (from 3.3 and 3.4). Lee chooses among them.
% New: four constants, two joins; cumulative (cos x); the answer pieces are the tangent lines at the joins.
\starquestion{exercise}{e:ex3:fourconst}{}{Let
\[
f(x)=\begin{cases} ax+b & \text{if } x<0\\ \cos x+x & \text{if } 0\le x\le\pi\\ cx+d & \text{if } x>\pi.\end{cases}
\]
Find the values of the constants $a$, $b$, $c$ and $d$ for which $f$ is differentiable for all $x$.}
\workingspace

% New: the join point itself is unknown; the slope equation gives two candidates and the value
% equation rejects one. Part (b) is cumulative (ln x): the tangent line to y = ln x through the origin.
\starquestion{exercise}{g:ex3:joinpoint}{}{%
\begin{enumerate}[label=(\alph*)]
 \item Let $f(x)=\begin{cases} x^3 & \text{if } x\le a\\ 12x-16 & \text{if } x>a.\end{cases}$ Find all values of the constant $a$ for which $f$ is differentiable for all $x$.
 \item Let $g(x)=\begin{cases} \ln x & \text{if } 0<x\le a\\ mx & \text{if } x>a,\end{cases}$ where $a>0$ and $m$ are constants. Find the values of $a$ and $m$ for which $g$ is differentiable for all $x>0$.
\end{enumerate}}
\workingspace

% BC 2008 MC 25 = Q65 = Q66 (P). Five choices trimmed to four: "4" dropped.
\question{exercise}{g:pr:cd}{}{%
\[
f(x)=\begin{cases} cx+d & \text{for } x\le2\\ x^2-cx & \text{for } x>2\end{cases}
\]
Let $f$ be the function defined above, where $c$ and $d$ are constants. If $f$ is differentiable at $x=2$, what is the value of $c+d$?}
\fourchoices{$-4$}{$-2$}{$0$}{$2$}
\workingspace
% =====================================================================
\ifshowsolutions
\clearpage
\section{Detailed solutions}
% =====================================================================
\Needspace{14\baselineskip}
\subsection*{Practice}\addcontentsline{toc}{subsection}{Practice}
\Needspace{12\baselineskip}
\subsubsection*{Section 1: Derivative Rules}
\solitem{pr:power}{%
\begin{enumerate}[label=(\alph*)]
 \item
 \begin{align*}
 y&=8x^{\frac13}+2x^{-3}-\frac74x^{-1}+\sqrt3 &&\why{rewrite every term as $cx^n$}\\
 \frac{dy}{dx}&=8\cdot\frac13x^{-\frac23}+2(-3)x^{-4}-\frac74(-1)x^{-2}+0 &&\why{power rule; $\sqrt3$ is a constant}\\
 &=\frac{8}{3x^{\frac23}}-\frac{6}{x^4}+\frac{7}{4x^2}.
 \end{align*}
 \item
 \begin{align*}
 y&=4x^4-4x^2+1 &&\why{expand}\\
 \frac{dy}{dx}&=16x^3-8x.
 \end{align*}
 \item
 \begin{align*}
 f(x)&=\frac{x^3}{x^2}-\frac{2x}{x^2}+\frac{5}{x^2}=x-2x^{-1}+5x^{-2} &&\why{split}\\
 f'(x)&=1+2x^{-2}-10x^{-3}=1+\frac{2}{x^2}-\frac{10}{x^3}\\
 f'(2)&=1+\frac24-\frac{10}{8}=\frac14.
 \end{align*}
 (The trap answer is $f(2)=\dfrac{8-4+5}{4}=\dfrac94$.)
\end{enumerate}
}

\solitem{pr:fifths}{%
\begin{align*}
y&=\frac12x^{\frac45}-3x^{-5} &&\why{rewrite as powers}\\
\frac{dy}{dx}&=\frac12\cdot\frac45x^{-\frac15}-3(-5)x^{-6} &&\why{power rule: $\dfrac45-1=-\dfrac15$, $-5-1=-6$}\\
&=\frac{2}{5x^{\frac15}}+\frac{15}{x^6}.
\end{align*}
The answer is (A). Choice (B) lowers the power $-5$ to $-4$ instead of $-6$.
}

\solitem{pr:negpower}{%
\begin{align*}
f'(x)&=4(-2)x^{-3}+\frac14\cdot2x+0=-\frac{8}{x^3}+\frac x2 &&\why{power rule}\\
f'(2)&=-\frac88+\frac22=-1+1=0.
\end{align*}
The answer is (C).
}

\solitem{pr:basicfns}{%
\begin{enumerate}[label=(\alph*)]
 \item $e^2$ is a constant.
 \begin{align*}
 f'(x)&=4(-\sin x)-5\cos x+2e^x-0=-4\sin x-5\cos x+2e^x\\
 f'(0)&=-4(0)-5(1)+2(1)=-3. &&\why{$\sin0=0$, $\cos0=1$, $e^0=1$}
 \end{align*}
 \item
 \begin{align*}
 y&=3\ln x+4x^{-1}-\ln2 &&\why{$\ln2$ is a constant}\\
 \frac{dy}{dx}&=\frac3x-4x^{-2}=\frac3x-\frac{4}{x^2}.
 \end{align*}
 \item $2^x$ is an exponential and $e$ is a constant.
 \[
 \frac{dy}{dx}=2^x\ln2-2x.
 \]
 \item Split the $\ln$ first.
 \begin{align*}
 f(x)&=\ln5+\ln x+4x^{\frac12} &&\why{$\ln(5x)=\ln5+\ln x$}\\
 f'(x)&=0+\frac1x+2x^{-\frac12}=\frac1x+\frac{2}{\sqrt x} &&\why{$\ln5$ is a constant}\\
 f'(4)&=\frac14+\frac22=\frac54.
 \end{align*}
 \item
 \begin{align*}
 f'(x)&=-3\sin x-2\cos x\\
 f'\Bigl(\frac{\pi}{6}\Bigr)&=-3\cdot\frac12-2\cdot\frac{\sqrt3}{2}=-\frac32-\sqrt3.
 \end{align*}
\end{enumerate}
}

\solitem{pr:product}{%
\begin{enumerate}[label=(\alph*)]
 \item
 \begin{align*}
 \frac{dy}{dx}&=\frac{d}{dx}\bigl(x^4\bigr)\cdot\sin x+\frac{d}{dx}(\sin x)\cdot x^4 &&\why{product rule}\\
 &=4x^3\sin x+x^4\cos x.
 \end{align*}
 \item
 \begin{align*}
 f'(x)&=\frac{d}{dx}\bigl(e^x\bigr)\cdot\bigl(x^2+2x\bigr)+\frac{d}{dx}\bigl(x^2+2x\bigr)\cdot e^x &&\why{product rule}\\
 &=e^x\bigl(x^2+2x\bigr)+(2x+2)e^x=e^x\bigl(x^2+4x+2\bigr) &&\why{factor out $e^x$}\\
 f'(1)&=e(1+4+2)=7e.
 \end{align*}
 \item
 \begin{align*}
 \frac{dy}{dx}&=\frac{d}{dx}\bigl(x^{\frac13}\bigr)\cdot\ln x+\frac{d}{dx}(\ln x)\cdot x^{\frac13} &&\why{product rule}\\
 &=\frac{1}{3x^{\frac23}}\ln x+\frac1x\cdot x^{\frac13}=\frac{\ln x}{3x^{\frac23}}+\frac{1}{x^{\frac23}} &&\why{$\dfrac{x^{\frac13}}{x}=x^{-\frac23}$}\\
 &=\frac{\ln x+3}{3x^{\frac23}}.
 \end{align*}
 \item Two polynomials: expand first.
 \begin{align*}
 y&=2x^3-x^2+6x-3 &&\why{expand}\\
 \frac{dy}{dx}&=6x^2-2x+6.
 \end{align*}
 \item Three factors: three terms.
 \begin{align*}
 f'(x)&=\frac{d}{dx}\bigl(x^3\bigr)\cdot2^x\ln x+\frac{d}{dx}\bigl(2^x\bigr)\cdot x^3\ln x+\frac{d}{dx}(\ln x)\cdot x^3\,2^x\\
 &=3x^2\,2^x\ln x+2^x\ln2\cdot x^3\ln x+\frac1x\cdot x^3\,2^x\\
 f'(1)&=0+0+1\cdot2=2. &&\why{$\ln1=0$}
 \end{align*}
\end{enumerate}
}

\solitem{pr:quotient}{%
\begin{enumerate}[label=(\alph*)]
 \item
 \begin{align*}
 \frac{dy}{dx}&=\frac{\dfrac{d}{dx}(x)\cdot e^x-\dfrac{d}{dx}\bigl(e^x\bigr)\cdot x}{\bigl(e^x\bigr)^2} &&\why{quotient rule}\\
 &=\frac{e^x-e^x\,x}{\bigl(e^x\bigr)^2}=\frac{e^x(1-x)}{e^x\cdot e^x}=\frac{1-x}{e^x}. &&\why{cancel one $e^x$}
 \end{align*}
 \item
 \begin{align*}
 \frac{dy}{dx}&=\frac{\dfrac{d}{dx}(\cos x)\cdot(1-\sin x)-\dfrac{d}{dx}(1-\sin x)\cdot\cos x}{(1-\sin x)^2} &&\why{quotient rule}\\
 &=\frac{(-\sin x)(1-\sin x)-(-\cos x)\cos x}{(1-\sin x)^2}\\
 &=\frac{-\sin x+\sin^2x+\cos^2x}{(1-\sin x)^2}=\frac{1-\sin x}{(1-\sin x)^2} &&\why{$\sin^2x+\cos^2x=1$}\\
 &=\frac{1}{1-\sin x}.
 \end{align*}
 \item One term on the bottom: split.
 \begin{align*}
 y&=\frac{4x^3}{x^2}-\frac{x^{\frac12}}{x^2}=4x-x^{-\frac32} &&\why{split}\\
 \frac{dy}{dx}&=4+\frac32x^{-\frac52}.
 \end{align*}
 \item $2^x$ is not a power of $x$, so use the quotient rule.
 \begin{align*}
 f'(x)&=\frac{\dfrac{d}{dx}\bigl(2^x\bigr)\cdot x-\dfrac{d}{dx}(x)\cdot2^x}{x^2} &&\why{quotient rule}\\
 &=\frac{2^x\ln2\cdot x-2^x}{x^2}=\frac{2^x(x\ln2-1)}{x^2}\\
 f'(2)&=\frac{4(2\ln2-1)}{4}=2\ln2-1.
 \end{align*}
 \item
 \begin{align*}
 f'(x)&=\frac{\dfrac{d}{dx}(\ln x)\cdot(2x-1)-\dfrac{d}{dx}(2x-1)\cdot\ln x}{(2x-1)^2} &&\why{quotient rule}\\
 &=\frac{\dfrac1x(2x-1)-2\ln x}{(2x-1)^2}\\
 f'(1)&=\frac{1\cdot1-0}{1^2}=1. &&\why{$\ln1=0$}
 \end{align*}
\end{enumerate}
}

\solitem{pr:quotcubic}{%
\begin{align*}
f'(x)&=\frac{\dfrac{d}{dx}(5-x)\cdot\bigl(x^3+2\bigr)-\dfrac{d}{dx}\bigl(x^3+2\bigr)\cdot(5-x)}{\bigl(x^3+2\bigr)^2} &&\why{quotient rule}\\
&=\frac{(-1)\bigl(x^3+2\bigr)-3x^2(5-x)}{\bigl(x^3+2\bigr)^2}\\
&=\frac{-x^3-2-15x^2+3x^3}{\bigl(x^3+2\bigr)^2}=\frac{2x^3-15x^2-2}{\bigl(x^3+2\bigr)^2}.
\end{align*}
The answer is (C).
}

\solitem{pr:quotlinear}{%
\begin{align*}
f'(x)&=\frac{\dfrac{d}{dx}(3x-2)\cdot(2x+3)-\dfrac{d}{dx}(2x+3)\cdot(3x-2)}{(2x+3)^2} &&\why{quotient rule}\\
&=\frac{3(2x+3)-2(3x-2)}{(2x+3)^2}=\frac{6x+9-6x+4}{(2x+3)^2}=\frac{13}{(2x+3)^2}.
\end{align*}
The answer is (C). Choice (A) has the numerator in the wrong order; choice (D) adds instead of subtracting.
}

\solitem{pr:quotfg}{%
\begin{align*}
h'(x)&=\frac{f'(x)g(x)-g'(x)f(x)}{\bigl(g(x)\bigr)^2} &&\why{quotient rule}\\
h'(1)&=\frac{(6)(6)-(4)(3)}{6^2}=\frac{24}{36}=\frac23. &&\why{$f(1)=3$, $f'(1)=2+4=6$, $g(1)=6$, $g'(1)=4$}
\end{align*}
The answer is (B).
}

\solitem{pr:trigsix}{%
\begin{enumerate}[label=(\alph*)]
 \item
 \begin{align*}
 \dfrac{d}{dx}(\cot x)&=\dfrac{d}{dx}\left(\frac{\cos x}{\sin x}\right)=\frac{\dfrac{d}{dx}(\cos x)\cdot\sin x-\dfrac{d}{dx}(\sin x)\cdot\cos x}{\sin^2x} &&\why{quotient rule}\\
 &=\frac{(-\sin x)\sin x-\cos x\cos x}{\sin^2x}=\frac{-\bigl(\sin^2x+\cos^2x\bigr)}{\sin^2x}\\
 &=\frac{-1}{\sin^2x}=-\csc^2x.\\[4pt]
 \dfrac{d}{dx}(\csc x)&=\dfrac{d}{dx}\left(\frac{1}{\sin x}\right)=\frac{\dfrac{d}{dx}(1)\cdot\sin x-\dfrac{d}{dx}(\sin x)\cdot1}{\sin^2x} &&\why{quotient rule}\\
 &=\frac{0\cdot\sin x-\cos x}{\sin^2x}=-\frac{1}{\sin x}\cdot\frac{\cos x}{\sin x}=-\csc x\cot x.
 \end{align*}
 \item One rule for each term.
 \[
 \frac{dy}{dx}=-4\csc^2x-3\csc x\cot x-\sec x\tan x+2\cos x+\frac{4}{\sqrt x}-\frac3x.
 \]
 \item
 \begin{align*}
 \frac{dy}{dx}&=\frac{d}{dx}\bigl(e^x\bigr)\cdot\sec x+\frac{d}{dx}(\sec x)\cdot e^x &&\why{product rule}\\
 &=e^x\sec x+\sec x\tan x\,e^x=e^x\sec x(1+\tan x).
 \end{align*}
 \item Simplify first.
 \begin{align*}
 f(x)&=\frac{\cos x}{\sin x}\cdot\sin x=\cos x\\
 f'(x)&=-\sin x\\
 f'\Bigl(\frac{\pi}{6}\Bigr)&=-\sin\frac{\pi}{6}=-\frac12.
 \end{align*}
 \item
 \begin{align*}
 \frac{dy}{dx}&=\frac{d}{dx}\bigl(5^x\bigr)\cdot\tan x+\frac{d}{dx}(\tan x)\cdot5^x &&\why{product rule}\\
 &=5^x\ln5\,\tan x+5^x\sec^2x=5^x\bigl(\ln5\,\tan x+\sec^2x\bigr).
 \end{align*}
 \item
 \begin{align*}
 f'(x)&=-\csc x\cot x+\csc^2x\\
 f'\Bigl(\frac{\pi}{6}\Bigr)&=-2\sqrt3+2^2=4-2\sqrt3. &&\why{$\csc\dfrac{\pi}{6}=2$, $\cot\dfrac{\pi}{6}=\sqrt3$}
 \end{align*}
\end{enumerate}
}

\solitem{pr:tanminuscot}{%
\[
\frac{dy}{dx}=\sec^2x-\bigl(-\csc^2x\bigr)=\sec^2x+\csc^2x.
\]
The answer is (D). Choice (C) forgets the minus sign in the derivative of $\cot x$.
}

\solitem{pr:derivecot}{%
\begin{itemize}
 \item (A), (B): the derivative of $\dfrac{1}{\tan x}$ is not $\dfrac{1}{\text{derivative of }\tan x}$.
 \item (C): the quotient rule with top $1$ (derivative $0$) and bottom $\tan x$ (derivative $\sec^2x$), in the correct order. Correct.
 \item (D): the numerator has a plus sign instead of a minus sign.
\end{itemize}
(Check: $-\dfrac{\sec^2x}{\tan^2x}=-\dfrac{1}{\cos^2x}\cdot\dfrac{\cos^2x}{\sin^2x}=-\dfrac{1}{\sin^2x}=-\csc^2x$.) The answer is (C).
}

\solitem{pr:costan}{%
Simplify first.
\begin{align*}
\cos x\tan x&=\cos x\cdot\frac{\sin x}{\cos x}=\sin x\\
\frac{d}{dx}(\cos x\tan x)&=\frac{d}{dx}(\sin x)=\cos x.
\end{align*}
The answer is (B).
}

\solitem{pr:tablerules}{%
\begin{enumerate}[label=(\alph*)]
 \item Use the row $x=0$; the middle term is a product of three factors.
 \begin{align*}
 h'(x)&=4g'(x)-\Bigl[e^xf(x)\cos x+f'(x)e^x\cos x-\sin x\,e^xf(x)\Bigr]-3^x\ln3\\
 h'(0)&=4(-4)-\bigl[(1)(-2)(1)+(3)(1)(1)-0\bigr]-\ln3\\
 &=-16-1-\ln3=-17-\ln3.
 \end{align*}
 \item Use the row $x=1$.
 \begin{itemize}
  \item Top $f(x)g(x)$: value $(-1)(3)=-3$; derivative $f'(1)g(1)+g'(1)f(1)=(4)(3)+(-2)(-1)=14$.
  \item Bottom $2\sqrt x+\ln x$: value $2+0=2$; derivative $\dfrac{1}{\sqrt x}+\dfrac1x=1+1=2$ at $x=1$.
 \end{itemize}
 \[
 h'(1)=\frac{(14)(2)-(2)(-3)}{2^2}=\frac{28+6}{4}=\frac{17}{2}.
 \]
 \item Use the row $x=0$.
 \begin{itemize}
  \item Top $f(x)-\sin x$: value $-2-0=-2$; derivative $f'(0)-\cos0=3-1=2$.
  \item Bottom $\tan x+g(x)$: value $0+1=1$; derivative $\sec^20+g'(0)=1-4=-3$.
 \end{itemize}
 \[
 h'(0)=\frac{(2)(1)-(-3)(-2)}{1^2}=2-6=-4.
 \]
\end{enumerate}
}
\Needspace{12\baselineskip}
\subsubsection*{Section 2: What the Derivative Means}
\solitem{a:pr:defvar}{%
\begin{enumerate}[label=(\alph*)]
 \item Replace every $x$ by $x+h$. For $f'(-1)$: $f(-1)=-2-1=-3$.
 \begin{align*}
 f'(x)&=\lim_{h\to0}\frac{2(x+h)^3+(x+h)-\bigl(2x^3+x\bigr)}{h},\\
 f'(-1)&=\lim_{h\to0}\frac{2(-1+h)^3+(-1+h)+3}{h}=\lim_{x\to-1}\frac{2x^3+x+3}{x+1}.
 \end{align*}
 (Check: $f'(x)=6x^2+1$, so $f'(-1)=7$.)
 \item
 \begin{align*}
 f'(9)&=\lim_{x\to9}\frac{\sqrt x-3}{x-9} &&\why{$f(9)=3$}\\
 &=\lim_{x\to9}\frac{\sqrt x-3}{\bigl(\sqrt x-3\bigr)\bigl(\sqrt x+3\bigr)}=\lim_{x\to9}\frac{1}{\sqrt x+3}=\frac16.
 \end{align*}
 \item
 \begin{align*}
 f'(x)&=\lim_{h\to0}\frac{\dfrac{1}{x+h+2}-\dfrac{1}{x+2}}{h}\\
 &=\lim_{h\to0}\frac{\dfrac{(x+2)-(x+h+2)}{(x+h+2)(x+2)}}{h} &&\why{common denominator}\\
 &=\lim_{h\to0}\frac{-h}{h\,(x+h+2)(x+2)}=\lim_{h\to0}\frac{-1}{(x+h+2)(x+2)} &&\why{cancel $h$}\\
 &=-\frac{1}{(x+2)^2},\\
 f'(2)&=-\frac{1}{16}.
 \end{align*}
 (Check with the quotient rule: $\dfrac{0\cdot(x+2)-1\cdot1}{(x+2)^2}=-\dfrac{1}{(x+2)^2}$.)
\end{enumerate}
}

\solitem{a:pr:defexpr}{%
By the definition, $f'(x)=\displaystyle\lim_{h\to0}\frac{f(x+h)-f(x)}{h}$, and $f(x+h)$ replaces every $x$ by $x+h$:
\[
f(x+h)=3(x+h)^2+2(x+h).
\]
The answer is (C). Choice (A) adds $h$ outside $f$ (it is $f(x)+h$, not $f(x+h)$); (B) and (D) send $x\to0$ instead of $h\to0$.
}

\solitem{a:pr:derivlimits}{%
\begin{enumerate}[label=(\alph*)]
 \item $f(x)=\cos x$ and $a=\dfrac{\pi}{5}$:\quad the limit is $f'\bigl(\frac{\pi}{5}\bigr)=-\sin\dfrac{\pi}{5}$.
 \item The definition of $f'(3)$:\quad $f'(3)=\dfrac13$. (The trap answer $\ln3$ is $f(3)$.)
 \item $2=\sqrt4$, so $f(x)=\sqrt x$ and $a=4$:
 \[
 \lim_{h\to0}\frac{\sqrt{4+h}-\sqrt4}{h}=f'(4)=\frac{1}{2\sqrt4}=\frac14.
 \]
 \item The step is $-h$, so the bottom must be $-h$ too: $f(x)=\ln x$, $a=3$.
 \[
 \lim_{h\to0}\frac{\ln(3-h)-\ln3}{h}=-\lim_{h\to0}\frac{\ln\bigl(3+(-h)\bigr)-\ln3}{-h}=-f'(3)=-\frac13.
 \]
 \item The step on the top is $h^2$; multiply the top and the bottom by $h$ to make the bottom $h^2$:
 \[
 \lim_{h\to0}\frac{\sin\bigl(\pi+h^2\bigr)-\sin\pi}{h}=\lim_{h\to0}h\cdot\frac{\sin\bigl(\pi+h^2\bigr)-\sin\pi}{h^2}=0\cdot f'(\pi)=0\cdot\cos\pi=0.
 \]
 \item As $h\to0$ the quotient tends to $\dfrac{d}{dx}(\ln x)=\dfrac1x$, so for $h$ very close to $0$ it is approximately $\dfrac1x$. When $x=2e$ the limit is $\dfrac{1}{2e}$.
\end{enumerate}
}

\solitem{a:pr:dqtable}{%
$f(4)=3$, so $6-2f(x)=-2\bigl(f(x)-f(4)\bigr)$ and $4-x=-(x-4)$.
\begin{align*}
\lim_{x\to4}\frac{6-2f(x)}{4-x}&=\lim_{x\to4}\frac{-2\bigl(f(x)-f(4)\bigr)}{-(x-4)}=2\lim_{x\to4}\frac{f(x)-f(4)}{x-4}\\
&=2f'(4) &&\why{$f$ is differentiable at $x=4$}\\
&=2(14)=28.
\end{align*}
}

\solitem{a:pr:evallim}{%
\begin{enumerate}[label=(\alph*)]
 \item
 \begin{align*}
 &\lim_{h\to0}\frac{(x+h)^2+4(x+h)-x^2-4x}{h}\\
 &=\lim_{h\to0}\frac{x^2+2xh+h^2+4x+4h-x^2-4x}{h} &&\why{expand}\\
 &=\lim_{h\to0}\frac{2xh+h^2+4h}{h}=\lim_{h\to0}(2x+h+4) &&\why{cancel $h$}\\
 &=2x+4.
 \end{align*}
 It is $f'(x)$ for $f(x)=x^2+4x$.
 \item This is $f'(2)$ for $f(x)=x^3+2x$.
 \begin{align*}
 &\lim_{h\to0}\frac{(2+h)^3+2(2+h)-8-4}{h}\\
 &=\lim_{h\to0}\frac{8+12h+6h^2+h^3+4+2h-12}{h} &&\why{$(2+h)^3=8+12h+6h^2+h^3$}\\
 &=\lim_{h\to0}\frac{14h+6h^2+h^3}{h}=\lim_{h\to0}\bigl(14+6h+h^2\bigr)\\
 &=14.
 \end{align*}
 (Check: $f'(x)=3x^2+2$, $f'(2)=14$.)
\end{enumerate}
}

\solitem{a:pr:dqgraph}{%
With $h=0.001$, very close to $0$,
\[
\frac{\cos(x+0.001)-\cos x}{0.001}\approx\frac{d}{dx}(\cos x)=-\sin x.
\]
So the graph looks like $y=-\sin x$: $0$ at $x=0$, $-1$ at $x=\dfrac{\pi}{2}$ and $1$ at $x=-\dfrac{\pi}{2}$. The answer is (D). Graph (A) is $\cos x$ itself.
}

\solitem{b:pr:tangent}{%
\begin{enumerate}[label=(\alph*)]
 \item
 \begin{align*}
 g(-1)&=-2+1-5=-6 &&\why{the point is $(-1,-6)$}\\
 g'(x)&=6x^2+2x\imp g'(-1)=6-2=4 &&\why{the slope}\\
 y+6&=4(x+1).
 \end{align*}
 \item The value is given; the slope comes from the formula for $h'$.
 \begin{align*}
 h'(4)&=\frac{4^2-2}{4}=\frac{14}{4}=\frac72 &&\why{substitute $x=4$}\\
 y-(-3)&=\frac72(x-4),\quad\text{i.e. } y+3=\frac72(x-4).
 \end{align*}
 \item The line tangent to $f$: point $(0,-1)$, slope $3$, so $y+1=3x$.

 For $g$: substitute $x=0$, $f(0)=-1$ and $f'(0)=3$.
 \begin{align*}
 g'(0)&=\bigl(3f(0)+f'(0)\bigr)\cos0=\bigl(3(-1)+3\bigr)(1)=0 &&\why{$\cos0=1$}\\
 y-5&=0(x-0),\quad\text{i.e. } y=5. &&\why{a horizontal tangent}
 \end{align*}
 \item The graph of $f'$ passes through $(0,-2)$, so $f'(0)=-2$ (the height of the graph of $f'$, not a slope). The point is $(0,3)$.
 \[
 y-3=-2(x-0),\quad\text{i.e. } y=3-2x.
 \]
 \item Requirement 2: value $1$ and slope $2$ at $x=2$.
 \begin{enumerate}[label=(\roman*)]
  \item $f'(x)=2ax$.
  \begin{align*}
  f(2)&=4a=1\imp a=\frac14 &&\why{the value}\\
  f'(2)&=4a=2\imp a=\frac12 &&\why{the slope}
  \end{align*}
  The two equations need different values of $a$, so no value of $a$ works.
  \item $g'(x)=3cx^2-\dfrac x2$.
  \begin{align*}
  g(2)&=8c-1=1\imp c=\frac14 &&\why{the value}\\
  g'(2)&=12c-1=2\imp c=\frac14 &&\why{the slope}
  \end{align*}
  So $c=\dfrac14$.
  \item $h'(x)=\dfrac{nx^{n-1}}{k}$.
  \begin{align*}
  h(2)&=\frac{2^n}{k}=1\imp k=2^n &&\why{the value}\\
  h'(2)&=\frac{n\cdot2^{n-1}}{2^n}=\frac n2=2\imp n=4 &&\why{the slope}
  \end{align*}
  So $n=4$ and $k=2^4=16$, i.e. $h(x)=\dfrac{x^4}{16}$.

  Requirement 1: $h(0)=0$ and $h'(0)=\dfrac{4\cdot0^3}{16}=0$.
 \end{enumerate}
 \item Substitute the point $x=-1$, $y=f(-1)=2$ into $\dfrac{dy}{dx}$.
 \begin{align*}
 \frac{dy}{dx}\Big|_{(-1,\,2)}&=\frac{-(-1)(2)^2}{2}=\frac{4}{2}=2\\
 y-2&=2(x+1).
 \end{align*}
 \item
 \begin{enumerate}[label=(\roman*)]
  \item
  \begin{align*}
  f'(-2)&=-3 &&\why{the slope of $y=-3x+1$}\\
  f(-2)&=-3(-2)+1=7. &&\why{the value of the line at $x=-2$}
  \end{align*}
  $f(0)$ need not be $1$: that is the value of the line at $x=0$, and the graph of $f$ only touches the line at $x=-2$.
  \item
  \begin{align*}
  f'(-1)&=\frac{11-3}{3-(-1)}=\frac{8}{4}=2\\
  y-3&=2(x+1).
  \end{align*}
 \end{enumerate}
\end{enumerate}
}

\solitem{b:pr:tanrules}{%
\begin{enumerate}[label=(\alph*),start=8]
 \item Slope: $y=x^{-2}$, so $\dfrac{dy}{dx}=-2x^{-3}=-\dfrac{2}{x^3}$; at $x=w$ the slope is $-\dfrac{2}{w^3}$.
 \begin{align*}
 \ell\colon\quad y-\frac{1}{w^2}&=-\frac{2}{w^3}(x-w)\\
 \text{at } R\ (y=0)\colon\quad -\frac{1}{w^2}&=-\frac{2}{w^3}(k-w)\imp \frac{w}{2}=k-w &&\why{multiply by $-\dfrac{w^3}{2}$}\\
 k&=\frac{3w}{2}.
 \end{align*}
 When $w=3$: $k=\dfrac92$.
 \item Check the point and the slope at $x=3$.
 \begin{itemize}
  \item Same point: $f(3)=4(9)-27=9$, and the line gives $18-3(3)=9$.
  \item Same slope: $f'(x)=8x-3x^2$, so $f'(3)=24-27=-3$, the slope of $y=18-3x$.
 \end{itemize}
 So $\ell$ passes through $\bigl(3,f(3)\bigr)$ with slope $f'(3)$: it is tangent to the graph of $f$ at $x=3$.
 \item (i) Parallel to $y=9x-1$: slope $9$.
 \begin{align*}
 \frac{dy}{dx}=3x^2+6x&=9\imp x^2+2x-3=0\imp(x+3)(x-1)=0\imp x=1\text{ or }x=-3.\\
 y(1)&=1+3-5=-1,\qquad y(-3)=-27+27-5=-5.
 \end{align*}
 The points are $(1,-1)$ and $(-3,-5)$.

 (ii) The graph of $g$ has slope $-1$, so the perpendicular slope is $-\dfrac{1}{-1}=1$.
 \[
 f'(x)=\frac{1}{2\sqrt x}=1\imp\sqrt x=\frac12\imp x=\frac14,\qquad f\Bigl(\frac14\Bigr)=\frac12.
 \]
 $P=\Bigl(\dfrac14,\dfrac12\Bigr)$.
 \item (i)
 \begin{align*}
 \frac{dy}{dx}=2x-5&=3\imp x=4\imp y=16-20+1=-3 &&\why{the point of tangency is $(4,-3)$}\\
 -3&=3(4)+k\imp k=-15. &&\why{the line passes through $(4,-3)$}
 \end{align*}
 (ii) The line through $(0,1)$ and $(2,-3)$ has slope $\dfrac{-3-1}{2-0}=-2$. Parallel: $g'(1)=-2$.
 \[
 g'(x)=2ax+3\imp g'(1)=2a+3=-2\imp a=-\frac52.
 \]
\end{enumerate}
}

\solitem{b:pr:twopoints}{%
Both points lie on the tangent line at $x=1$, so its slope is $f'(1)$.
\[
f'(1)=\frac{7-(-2)}{1-(-2)}=\frac{9}{3}=3.
\]
The answer is (B). Choice (C) is $f(1)$.
}

\solitem{b:pr:drawntangent}{%
$f'(3)$ is the slope of the tangent line. Read two grid points on the line: $(1,5)$ and $(5,3)$.
\[
f'(3)=\frac{3-5}{5-1}=-\frac12.
\]
The answer is (A). Choice (C) is $f(3)=4$, choice (D) is the $y$-intercept of the tangent line, and choice (B) is the negative reciprocal of the slope.
}

\solitem{b:pr:tanlinegraph}{%
The line crosses the $y$-axis at $-2$ and the $x$-axis at $-1$.
\begin{align*}
f(0)&=-2 &&\why{the line meets the graph of $f$ at $x=0$}\\
f'(0)&=\frac{0-(-2)}{-1-0}=-2. &&\why{the slope of the line}
\end{align*}
So $f'(0)=f(0)$. The answer is (C).
}

\solitem{b:pr:limitslope}{%
The slope of the tangent line at $x=-1$ is
\[
f'(-1)=\lim_{x\to-1}\frac{f(x)-f(-1)}{x-(-1)}=\lim_{x\to-1}\frac{f(x)-f(-1)}{x+1}=-\frac35.
\]
The answer is (A). Choice (C) is the limit of $f(x)$, a $y$-value, not a slope.
}

\solitem{b:pr:tangentk}{%
The line is $y=-x+k$, slope $-1$. Find where the curve has slope $-1$.
\begin{align*}
\frac{dy}{dx}=2x+3&=-1\imp x=-2\imp y=4-6+1=-1 &&\why{the point of tangency is $(-2,-1)$}\\
k&=x+y=-2+(-1)=-3. &&\why{the line passes through $(-2,-1)$}
\end{align*}
The answer is (A). Choice (B) is $x$ and choice (C) is $y$ at the point of tangency.
}

\solitem{b:pr:parallelb}{%
The line through $(0,-2)$ and $(3,4)$ has slope $\dfrac{4-(-2)}{3-0}=2$. Parallel: $g'(-1)=2$.
\[
g'(x)=2x+b\imp g'(-1)=-2+b=2\imp b=4.
\]
The answer is (D).
}

\solitem{b:pr:estimate}{%
\begin{enumerate}[label=(\alph*)]
 \item
 \begin{enumerate}[label=(\roman*)]
  \item $t=22$ is on the segment from $(20,15)$ to $(24,3)$, and $t=14$ on the segment from $(12,9)$ to $(16,15)$. At $t=2$, $f$ is the formula: differentiate.
  \begin{align*}
  f'(22)&=\frac{3-15}{24-20}=-3\text{ calories per minute per minute},\\
  f'(14)&=\frac{15-9}{16-12}=1.5\text{ calories per minute per minute},\\
  f'(t)&=-\frac34t^2+3t\imp f'(2)=-3+6=3\text{ calories per minute per minute}.
  \end{align*}
  \item From the graph: $f(14)=12$ (halfway along the segment from $(12,9)$ to $(16,15)$) and $f'(14)=1.5$ from (i).
  \begin{align*}
  k'(t)&=\frac{d}{dt}(t)\cdot f(t)+\frac{d}{dt}\bigl(f(t)\bigr)\cdot t=f(t)+f'(t)\,t &&\why{product rule}\\
  k'(14)&=12+(1.5)(14)=33.
  \end{align*}
 \end{enumerate}
 \item
 \begin{enumerate}[label=(\roman*)]
  \item At $b$ the graph rises, so $f'(b)>0$; $c$ is the peak, so $f'(c)=0$; at $d$ the graph falls, so $f'(d)<0$.
  \[
  f'(d)<f'(c)<f'(b).
  \]
  \item $f(a)<0$ (below the $x$-axis), $f'(c)=0$, and $f(d)-f(c)<0$ ($f(c)$ is the highest value). That leaves $f'(a)$ and the secant slope $\dfrac{f(b)-f(a)}{b-a}$, both positive. The graph gets less steep from $a$ to $b$, so the tangent at $a$ is steeper than the secant from $a$ to $b$ (from the grid, about $3.8$ against about $2.8$). The greatest is $f'(a)$.
  \item The secant from $\bigl(a,f(a)\bigr)\approx(-2.3,-2.8)$ to $\bigl(d,f(d)\bigr)\approx(1.1,1.8)$ has slope about $\dfrac{4.6}{3.3}\approx1.4$. The slope of $f$ goes down steadily from left to right (the graph keeps bending down), so it equals $1.4$ at exactly one point. At $b$ the slope is about $1.8$, more than $1.4$; at $c$ it is $0$. So there is $1$ point, between $b$ and $c$.
 \end{enumerate}
\end{enumerate}
}

\solitem{b:pr:tickets}{%
$5.5$ is between $4$ and $7$.
\[
L'(5.5)\approx\frac{L(7)-L(4)}{7-4}=\frac{150-126}{3}=8\text{ people per hour.}
\]
}

\solitem{b:pr:tea}{%
$3.5$ is between $2$ and $5$.
\[
H'(3.5)\approx\frac{H(5)-H(2)}{5-2}=\frac{52-60}{3}=-\frac83\approx-2.667\text{ degrees Celsius per minute.}
\]
}

\solitem{b:pr:orderslopes}{%
\begin{itemize}
 \item At $x=-2$ the graph rises: $f'(-2)>0$.
 \item At $x=0$ the graph is flat (a horizontal tangent): $f'(0)=0$.
 \item At $x=3$ the graph falls steeply: $f'(3)<0$.
\end{itemize}
So $f'(3)<f'(0)<f'(-2)$. The answer is (D).
}

\solitem{b:pr:countsecant}{%
The limit is $f'(c)$, so the question asks for the points where the tangent line is parallel to the secant through $\bigl(a,f(a)\bigr)$ and $\bigl(b,f(b)\bigr)$. The secant rises gently. Slide a line with this slope across the graph:
\begin{itemize}
 \item from the valley to the peak on the left of the $y$-axis, the slope goes up from $0$ and back to $0$, passing the secant slope twice: $2$ points;
 \item just right of $x=0$, the graph rises steeply and levels off at a peak: $1$ point;
 \item after the next valley, the graph rises more and more steeply up to $b$: $1$ point.
\end{itemize}
The graph falls everywhere else, and $x=0$ does not count ($f'(0)$ does not exist). That makes $4$ values of $c$. The answer is (D).
}

\solitem{b:pr:rulesgraph}{%
\begin{enumerate}[label=(\alph*)]
 \item The derivative of a line is its slope.
 \begin{itemize}
  \item $f$ passes through $(0,1)$ and $(1,2)$: slope $1$.
  \item $g$ passes through $(0,2)$ and $(3,0)$: slope $-\dfrac23$.
 \end{itemize}
 \[
 h'(x)=f'(x)+g'(x)=1-\frac23=\frac13.
 \]
 The answer is (C). Choice (D) is $h(x)=(x+1)+\Bigl(-\dfrac23x+2\Bigr)=\dfrac13x+3$.
 \item Read $f'(2)=1.5$ from the graph of $f'$ (the point $(2,1.5)$).
 \begin{align*}
 g'(x)&=f'(x)-\frac1x\\
 g'(2)&=f'(2)-\frac12=1.5-0.5=1.
 \end{align*}
\end{enumerate}
}

\solitem{a:pr:aroc}{%
\begin{enumerate}[label=(\alph*)]
 \item
 \begin{align*}
 \frac{f\bigl(\frac{\pi}{2}\bigr)-f(0)}{\dfrac{\pi}{2}-0}&=\frac{e^{\frac{\pi}{2}}\sin\frac{\pi}{2}-e^0\sin0}{\dfrac{\pi}{2}}=\frac{e^{\frac{\pi}{2}}}{\dfrac{\pi}{2}}=\frac{2e^{\frac{\pi}{2}}}{\pi}. &&\why{$\sin\frac{\pi}{2}=1$, $\sin0=0$}
 \end{align*}
 \item $R$ is itself a rate (gallons per hour), so the answer is in gallons per hour per hour.
 \begin{align*}
 \frac{R(8)-R(2)}{8-2}&=\frac{25.553-33.929}{6} &&\why{calculator, radians}\\
 &=-1.396\text{ gallons per hour per hour}.
 \end{align*}
 \item
 \begin{align*}
 \frac{f(b)-f(2)}{b-2}&=6 &&\why{average rate of change over $[2,b]$}\\
 \frac{b^2-3b+2}{b-2}&=6 &&\why{$f(2)=4-6=-2$}\\
 \frac{(b-1)(b-2)}{b-2}&=6\\
 b-1&=6 &&\why{$b\ne2$}\\
 b&=7.
 \end{align*}
\end{enumerate}
}

\solitem{a:pr:costwox}{%
Only the two endpoint values are needed (no derivative of $\cos(2x)$).
\[
\frac{\cos\pi-\cos0}{\dfrac{\pi}{2}-0}=\frac{-1-1}{\dfrac{\pi}{2}}=-\frac{4}{\pi}.
\]
The answer is (A).
}

\solitem{a:pr:fprimerow}{%
\begin{enumerate}[label=(\alph*)]
 \item Use the $f(x)$ row:\quad $\dfrac{f(12)-f(3)}{12-3}=\dfrac{15-(-2)}{9}=\dfrac{17}{9}$.
 \item Now the function is $f'$, so use the $f'(x)$ row:\quad $\dfrac{f'(12)-f'(3)}{12-3}=\dfrac{9-5}{9}=\dfrac49$.
 \item Compute each one from the $f(x)$ row.
 \begin{align*}
 [-2,3]&:\ \frac{-2-(-10)}{5}=\frac85, &\qquad [3,5]&:\ \frac{-1-(-2)}{2}=\frac12,\\
 [5,12]&:\ \frac{15-(-1)}{7}=\frac{16}{7}, &\qquad [12,15]&:\ \frac{36-15}{3}=7.
 \end{align*}
 The greatest is on $[12,15]$.
\end{enumerate}
}

\solitem{a:pr:arocgraph}{%
Read the endpoints: $f(-5)=0$ (the start of the quarter circle) and $f(0)=2$ (the segment from $(-2,0)$ to $(0,2)$).
\[
\frac{f(0)-f(-5)}{0-(-5)}=\frac{2-0}{5}=\frac25.
\]
}

\solitem{a:pr:arocpiece}{%
$x=1$ is in the first piece and $x=5$ in the second: $g(1)=1^2+1=2$ and $g(5)=13-5=8$.
\[
\frac{g(5)-g(1)}{5-1}=\frac{8-2}{4}=\frac32.
\]
The answer is (B). The graph has a corner at $x=3$, but the average rate of change needs only $g(1)$ and $g(5)$, so ``Undefined'' is a trap.
}

\solitem{a:pr:instant}{%
\begin{enumerate}[label=(\alph*)]
 \item The best estimate of $P'(4)$ is the average rate of change over the shortest interval containing $4$, here $[3.9,4.1]$. The answer is (D). Choice (A) is a number of fish, not a rate.
 \item Average rate of change: $f(5)=125-50+25-16=84$ and $f(0)=-16$, so
 \[
 \frac{f(5)-f(0)}{5-0}=\frac{84-(-16)}{5}=20.
 \]
 Instantaneous rate: $f'(x)=3x^2-4x+5$.
 \[
 3x^2-4x+5=20\imp3x^2-4x-15=0\imp(3x+5)(x-3)=0\imp x=3\text{ or }x=-\frac53.
 \]
 Reject $-\dfrac53$ (not in $[0,5]$). The answer is $x=3$.
\end{enumerate}
}

\solitem{a:pr:plane}{%
The rate at $t=9$ is $h'(9)$. The best estimate is the average rate of change over the shortest interval containing $t=9$ whose endpoints are known: $[7,11]$.
\[
h'(9)\approx\frac{30{,}763-16{,}859}{11-7}.
\]
The answer is (D). The other choices use the longer intervals $[0,15]$, $[5,15]$ and $[7,15]$.
}

\solitem{b:pr:avgrate}{%
\begin{align*}
\frac{f(4)-f(1)}{4-1}&=\frac{(4+\ln4)-(1+0)}{3}=1+\frac{\ln4}{3}\approx1.462 &&\why{average rate of change}\\
f'(c)&=1+\frac1c &&\why{instantaneous rate}\\
1+\frac1c&=1+\frac{\ln4}{3}\imp c=\frac{3}{\ln4}\approx2.164.
\end{align*}
The answer is (C). (The calculator is needed only for the decimal; you can also intersect $y=1+\dfrac1x$ with $y=1.462$.) Choice (B) solves $f(c)=1.462$, using the value instead of the slope.
}

\solitem{c:ratesvar}{%
\begin{enumerate}[label=(\alph*)]
 \item The radius is fixed at $3$ meters, so
 \[
 V=\pi(3)^2h=9\pi h\imp V'(t)=9\pi h'(t). \qquad\why{constant-multiple rule}
 \]
 \[
 V'(5)=9\pi\bigl(0.5+0.2\cos(0.5)\bigr)=19.100. \qquad\why{radian mode}
 \]
 The volume is changing at $19.100$ cubic meters per hour.
 \item The rate is $R'(25)$.
 \begin{align*}
 R'(x)&=60+1.6x-0.05x^{1.5} &&\why{power rule: $0.02\cdot2.5=0.05$}\\
 R'(25)&=60+40-0.05(125)=93.75. &&\why{$25^{1.5}=\bigl(\sqrt{25}\bigr)^3=125$}
 \end{align*}
 The weekly revenue is changing at $93.75$ dollars per kilogram.
\end{enumerate}
}

\solitem{c:tankw}{%
\begin{align*}
w'(t)&=-2t\\
w'(3)&=-6.
\end{align*}
The answer is (D). Choice (B) is $w(3)=25-9=16$.
}

\solitem{c:wave}{%
\begin{align*}
v(h)&=3h^{\frac12}\\
v'(h)&=\frac32h^{-\frac12}=\frac{3}{2\sqrt h}\\
v'(2)&=\frac{3}{2\sqrt2}.
\end{align*}
The answer is (B).
}

\solitem{c:greatest}{%
\scratch{\texttt{Y1}$=$\texttt{X\^{}4+(1/2)X\^{}3-5X\^{}2+tan(X/2)}; \texttt{nDeriv(Y1,X,}$a$\texttt{)} for each $a$.}
\begin{align*}
f'(-2)&=-4.287, &f'(-1)&=8.149, &f'(0)&=0.5, &f'(1)&=-3.851.
\end{align*}
The greatest slope is at $x=-1$. The answer is (B).
}

\solitem{c:etan}{%
\steph{Step 1: Where the graph crosses the $x$-axis.}
\[
e^{\tan x}-2=0\imp x=0.60611\ldots \qquad\why{\texttt{2:zero} on $[0,1]$}
\]
\steph{Step 2: The slope there.}
\[
\frac{dy}{dx}\bigg|_{x=0.60611\ldots}=2.961. \qquad\why{\texttt{nDeriv(Y1,X,X)}}
\]
The answer is (C). Choice (A) is the $x$-value, not the slope.
}

\solitem{c:rainpiece}{%
$t=4$ is in the second piece, $2<t\le5$, so use $r(t)=115+25\sin\bigl(0.2t^2\bigr)$.
\[
r'(4)=-39.932\ \text{gallons per hour per hour}. \qquad\why{\texttt{nDeriv(115+25sin(0.2X\^{}2),X,4)}}
\]
(The first piece at $t=4$ gives the trap answer $-18.132$.)
}

\solitem{c:count}{%
The line $y=1+4x$ has slope $4$, so solve $f'(x)=4$ on $[0,2]$.
\scratch{\texttt{Y1}$=$\texttt{1-5cos(3X\^{}2-2)}, \texttt{Y2}$=4$, window $0\le x\le2$; count the crossings.}
\[
1-5\cos\bigl(3x^2-2\bigr)=4\imp x=1.185,\ 1.422,\ 1.871.
\]
There are three values. The answer is (C).
}

\solitem{c:origin}{%
\steph{Step 1: The $y$-intercept of $\ell$.}
The tangent at $(k,\cos k)$ is $y=\cos k-\sin k\,(x-k)$. At $x=0$:
\[
y\text{-intercept}=\cos k+k\sin k.
\]
\steph{Step 2: Through the origin.}
\[
\cos k+k\sin k=0\imp k=2.798 \quad\text{in } 0<k<\pi.
\]
The answer is (C).
}

\solitem{c:parallel}{%
\steph{Step 1: The slope of the tangent to $f$ at $x=2$.}
\[
f'(2)=0.8. \qquad\why{\texttt{nDeriv(ln(X\^{}2+1),X,2)}}
\]
\steph{Step 2: Equal slopes.}
\begin{align*}
g'(a)&=f'(2)\\
5a^4+3a^2&=0.8 &&\why{$g'(x)=5x^4+3x^2$}\\
a&=0.447. &&\why{$a>0$}
\end{align*}
The answer is (C). Choice (D) solves $g(a)=0.8$ instead of $g'(a)=0.8$.
}
\Needspace{12\baselineskip}
\subsubsection*{Section 3: Differentiability and Continuity}
\solitem{d:pr:dnevar}{%
\begin{enumerate}[label=(\alph*)]
 \item
 \begin{enumerate}[label=(\roman*)]
  \item $f$ is not differentiable at
  \begin{itemize}
   \item $x=-3$: a vertical tangent;
   \item $x=-1$: a hole ($f(-1)=1$, but the graph comes in to the open dot from both sides);
   \item $x=1$: a corner (a gentle positive slope on the left, slope $-2$ on the right);
   \item $x=4$: a jump;
   \item $x=5$: a cusp.
  \end{itemize}
  At $x=3$ the tangent is horizontal, so $f'(3)=0$ exists.
  \item Continuous at $x=-3$, $1$ and $5$ (vertical tangent, corner, cusp).
 \end{enumerate}
 \item
 \begin{enumerate}[label=(\roman*)]
  \item At $x=-3$: no, the graph has a corner (slope $\frac32$ on the left, $-\frac12$ on the right).\\
  At $x=-1$: no, on the right the semicircle has a vertical tangent at $(-1,2)$.
  \item The slopes of the segments: $\dfrac{3-0}{-3-(-5)}=\dfrac32$ and $\dfrac{2-3}{-1-(-3)}=-\dfrac12$.
  \[
  f'(x)=\begin{cases}\dfrac32 & \text{for } -5<x<-3\\[6pt] -\dfrac12 & \text{for } -3<x<-1\end{cases}
  \]
  \item $\displaystyle\lim_{x\to-1^-}f'(x)=-\frac12$, the slope of the segment.

  $\displaystyle\lim_{x\to-1^+}f'(x)=\DNEninf$: the lower semicircle starts at $(-1,2)$ with a vertical tangent, going down.
 \end{enumerate}
\end{enumerate}
}

\solitem{d:pr:vtan}{%
At $x=0$: $f(0)=0$, but the left branch comes in to the open dot at $(0,-1)$. A jump, so $f$ is not continuous and not differentiable at $0$.

At $x=2$: a vertical tangent, so $f'(2)$ does not exist.

At $x=1$ and $x=3$: horizontal tangents, so $f'(1)=f'(3)=0$ exist. The answer is (B).
}

\solitem{d:pr:contnotdiff}{%
\begin{itemize}
 \item $a$: the two pieces meet at a sharp peak, a corner. Continuous, not differentiable.
 \item $b$: a hole with the point plotted above it. Not continuous.
 \item $c$: smooth. Differentiable.
 \item $d$: a jump. Not continuous.
\end{itemize}
The answer is (A).
}

\solitem{d:pr:semicircles}{%
\begin{itemize}
 \item At $r$ the segment meets the lower semicircle at its left end, where the semicircle has a vertical tangent: $f'(r)$ is undefined. (The slopes are $\frac53$ on the left and $\DNEninf$ on the right.)
 \item At $s$ the lower semicircle has its lowest point: a horizontal tangent, $f'(s)=0$.
 \item At $t$ both semicircles have vertical tangents: $f'(t)$ is undefined.
\end{itemize}
The answer is (D).
}

\solitem{d:pr:falsereason}{%
A horizontal tangent means $f'(4)=0$: $f$ is differentiable at $x=4$. Statement (C) is false. The answer is (C).

(A), (B) and (D) are true: a jump, a vertical tangent and a hole all make the derivative fail to exist.
}

\solitem{d:pr:whichgraph}{%
$\displaystyle\lim_{x\to2}g(x)=g(2)$ means $g$ is continuous at $2$. That rules out the jumps in (A) and (D).

Unequal one-sided limits of $g'$ mean the slopes on the two sides of $x=2$ differ. In (C) the curve is smooth, with the same slope on both sides. In (B) the two segments meet at $(2,2)$ with slopes $2$ and $-2$. The answer is (B).
}

\solitem{d:pr:limfprime}{%
\begin{itemize}
 \item (A): $f(a)$ is the filled dot. True.
 \item (B): $\displaystyle\lim_{x\to a}f(x)$ is the height of the open dot, not $f(a)$. True.
 \item (C): both sides approach the open dot. True.
 \item (D): $f'(x)$ is the positive slope of the left segment for $x<a$ and the negative slope of the right segment for $x>a$, so $\displaystyle\lim_{x\to a^-}f'(x)\ne\lim_{x\to a^+}f'(x)$ and $\displaystyle\lim_{x\to a}f'(x)$ does not exist. Must be false.
\end{itemize}
The answer is (D).
}

\solitem{d:pr:onesided}{%
No. $f(0)=2$. On $-4\le x\le0$ the graph is the segment through $(-3,0)$ and $(0,2)$, so $f(x)=\dfrac23(x+3)$ there.
\[
\lim_{h\to0^-}\frac{f(0+h)-f(0)}{h}=\lim_{h\to0^-}\frac{\frac23(h+3)-2}{h}=\lim_{h\to0^-}\frac{\frac23h}{h}=\frac23.
\]
On $0<x<3$ the graph decreases from $(0,2)$ to $(3,0)$, so $f(h)<f(0)$ for $0<h<3$, and
\[
\frac{f(0+h)-f(0)}{h}<0 \quad\text{for } 0<h<3.
\]
So $\displaystyle\lim_{h\to0^+}\frac{f(0+h)-f(0)}{h}$, if it exists, is $\le0$, and it is not $\dfrac23$. The one-sided limits are not equal, so $f'(0)$ does not exist: $f$ is not differentiable at $x=0$.
}

\solitem{d:pr:dnevarf}{%
\begin{enumerate}[label=(\alph*)]
 \item We need to check whether $\displaystyle\lim_{h\to0}\frac{f(-3+h)-f(-3)}{h}$ exists, where $f(-3)=0$. Split it into the two one-sided limits.
 \twosides{%
 \begin{align*}
 &\lim_{h\to0^-}\frac{f(-3+h)-f(-3)}{h}\\
 &=\lim_{h\to0^-}\frac{\bigl|9-(-3+h)^2\bigr|-0}{h}\\
 &=\lim_{h\to0^-}\frac{\bigl|6h-h^2\bigr|}{h}\\
 &=\lim_{h\to0^-}\frac{\bigl|h(6-h)\bigr|}{h}\\
 &=\lim_{h\to0^-}\frac{-h(6-h)}{h}\\
 &=\lim_{h\to0^-}(-6+h)=-6
 \end{align*}}{%
 \begin{align*}
 &\lim_{h\to0^+}\frac{f(-3+h)-f(-3)}{h}\\
 &=\lim_{h\to0^+}\frac{\bigl|9-(-3+h)^2\bigr|-0}{h}\\
 &=\lim_{h\to0^+}\frac{\bigl|6h-h^2\bigr|}{h}\\
 &=\lim_{h\to0^+}\frac{\bigl|h(6-h)\bigr|}{h}\\
 &=\lim_{h\to0^+}\frac{h(6-h)}{h}\\
 &=\lim_{h\to0^+}(6-h)=6
 \end{align*}}
 (For $h$ close to $0$, $6-h>0$, so $h(6-h)$ has the sign of $h$.)

 $-6\ne6$, so $f'(-3)$ does not exist: $f$ is not differentiable at $x=-3$ (a corner).
 \item Continuity: $g(-2)=0-3=-3$ and $\displaystyle\lim_{x\to-2}g(x)=\sqrt0-3=-3=g(-2)$, so $g$ is continuous at $x=-2$.

 Differentiability: we need to check whether $\displaystyle\lim_{h\to0}\frac{g(-2+h)-g(-2)}{h}$ exists. Split it into the two one-sided limits.
 \twosides{%
 \begin{align*}
 &\lim_{h\to0^-}\frac{g(-2+h)-g(-2)}{h}\\
 &=\lim_{h\to0^-}\frac{\sqrt{|h|}-3+3}{h}=\lim_{h\to0^-}\frac{\sqrt{-h}}{h}\\
 &=\lim_{h\to0^-}\frac{-1}{\sqrt{-h}}=\DNEninf
 \end{align*}}{%
 \begin{align*}
 &\lim_{h\to0^+}\frac{g(-2+h)-g(-2)}{h}\\
 &=\lim_{h\to0^+}\frac{\sqrt{|h|}-3+3}{h}=\lim_{h\to0^+}\frac{\sqrt h}{h}\\
 &=\lim_{h\to0^+}\frac{1}{\sqrt h}=\DNEinf
 \end{align*}}
 The one-sided limits are not finite, so $g$ is continuous but not differentiable at $x=-2$ (a cusp).
 \item We need to check whether $\displaystyle\lim_{h\to0}\frac{f(3+h)-f(3)}{h}$ exists, where $f(3)=9-3=6$. For $h<0$ use $x^2-x$; for $h>0$ use $5x-9$.
 \twosides{%
 \begin{align*}
 &\lim_{h\to0^-}\frac{f(3+h)-f(3)}{h}\\
 &=\lim_{h\to0^-}\frac{(3+h)^2-(3+h)-6}{h}\\
 &=\lim_{h\to0^-}\frac{5h+h^2}{h}\\
 &=\lim_{h\to0^-}(5+h)=5
 \end{align*}}{%
 \begin{align*}
 &\lim_{h\to0^+}\frac{f(3+h)-f(3)}{h}\\
 &=\lim_{h\to0^+}\frac{5(3+h)-9-6}{h}\\
 &=\lim_{h\to0^+}\frac{5h}{h}\\
 &=\lim_{h\to0^+}5=5
 \end{align*}}
 Both are $5$, so $f$ is differentiable at $x=3$, and $f'(3)=5$.
 \item $f(2)=4-2+1=3$, from $x^2-x+1$.
 \begin{align*}
 \lim_{x\to2^-}\frac{f(x)-f(2)}{x-2}&=\lim_{x\to2^-}\frac{(3x+1)-3}{x-2}\\
 &=\lim_{x\to2^-}\frac{3x-2}{x-2}=\DNEninf, &&\why{top $\to4$, bottom $\to0^-$}\\
 \lim_{x\to2^+}\frac{f(x)-f(2)}{x-2}&=\lim_{x\to2^+}\frac{\bigl(x^2-x+1\bigr)-3}{x-2}=\lim_{x\to2^+}\frac{(x-2)(x+1)}{x-2}\\
 &=\lim_{x\to2^+}(x+1)=3.
 \end{align*}
 The left limit does not exist, so $f$ is not differentiable at $x=2$. The pieces do not meet ($3(2)+1=7\ne3$), even though both pieces have slope $3$ there.
\end{enumerate}
}

\solitem{d:pr:absderiv}{%
$f(3)=2$ and $f(3+h)-f(3)=|h|$.
\[
\lim_{h\to0^-}\frac{|h|}{h}=\lim_{h\to0^-}\frac{-h}{h}=-1,\qquad \lim_{h\to0^+}\frac{|h|}{h}=\lim_{h\to0^+}\frac{h}{h}=1.
\]
$-1\ne1$, so $f'(3)$ does not exist (a corner). The answer is (D).
}

\solitem{e:pr:ab2003}{%
I and II: $\displaystyle\lim_{x\to3^-}(x+2)=5$ and $\displaystyle\lim_{x\to3^+}(4x-7)=5$, so the limit exists and equals $f(3)=5$. I and II are true.

III: $f(3)=5$.
\begin{align*}
\lim_{x\to3^-}\frac{f(x)-f(3)}{x-3}&=\lim_{x\to3^-}\frac{(x+2)-5}{x-3}=1,\\
\lim_{x\to3^+}\frac{f(x)-f(3)}{x-3}&=\lim_{x\to3^+}\frac{(4x-7)-5}{x-3}=\lim_{x\to3^+}\frac{4(x-3)}{x-3}=4.
\end{align*}
$1\ne4$: a corner, so III is false. The answer is (C), I and II only.
}

\solitem{e:pr:removable}{%
For $x\ne2$, $f(x)=\dfrac{(x-2)(x+2)}{x-2}=x+2$.

I: $\displaystyle\lim_{x\to2}f(x)=\lim_{x\to2}(x+2)=4$. True.

II: $f(2)=1\ne4$. False.

III: $f(2)=1$, so
\[
\lim_{x\to2}\frac{f(x)-f(2)}{x-2}=\lim_{x\to2}\frac{(x+2)-1}{x-2}=\lim_{x\to2}\frac{x+1}{x-2}:
\]
the top tends to $3$ and the bottom to $0$, so the quotient blows up on both sides. False.

The answer is (A), I only.
}

\solitem{e:pr:q88}{%
\steph{At $x=-1$: $f(-1)=1-3=-2$}
\begin{align*}
\lim_{x\to-1^-}\frac{f(x)-f(-1)}{x-(-1)}&=\lim_{x\to-1^-}\frac{\dfrac2x+2}{x+1}=\lim_{x\to-1^-}\frac{\dfrac{2(x+1)}{x}}{x+1}=\lim_{x\to-1^-}\frac2x=-2,\\
\lim_{x\to-1^+}\frac{f(x)-f(-1)}{x-(-1)}&=\lim_{x\to-1^+}\frac{\bigl(x^2-3\bigr)+2}{x+1}=\lim_{x\to-1^+}(x-1)=-2.
\end{align*}
Equal: differentiable at $x=-1$.
\steph{At $x=2$: $f(2)=4-3=1$}
\[
\lim_{x\to2^+}\frac{f(x)-f(2)}{x-2}=\lim_{x\to2^+}\frac{(4x-3)-1}{x-2}=\lim_{x\to2^+}\frac{4x-4}{x-2}=\DNEinf,
\]
since the top tends to $4$ and the bottom to $0$ through positive values. So $f$ is not differentiable at $x=2$ (the pieces do not meet: $4(2)-3=5\ne1$), although both pieces have slope $4$ there.

The answer is (B), $x=2$ only.
}

\solitem{f:pr:diffcontvar}{%
\begin{enumerate}[label=(\alph*)]
 \item I and II only.
 \begin{itemize}
  \item I: $f$ is differentiable at $-1$ $\imp$ $f$ is continuous at $-1$ $\imp$ $\displaystyle\lim_{x\to-1}f(x)=f(-1)=2$, and then the one-sided limit is $2$ as well. True.
  \item II: $2=f(-1)$ and $x+1=x-(-1)$, so this is the definition of $f'(-1)$. True.
  \item III: this says $f'$ is continuous at $-1$, which is not given. Could be false.
 \end{itemize}
 \item I and III only.
 \begin{itemize}
  \item I: continuous at $-2$ means $\displaystyle\lim_{x\to-2}g(x)=g(-2)$, so the limit exists. True.
  \item II: continuous does not give differentiable: $g(x)=|x+2|$ is continuous at $-2$, but $g'(-2)$ does not exist. Could be false.
  \item III: this limit is the definition of $g'(-2)$, so if $g'(-2)$ does not exist, the limit does not exist. True.
 \end{itemize}
 \item $f(2)=-1$, so $\bigl(f(x)\bigr)^2-1=\bigl(f(x)+1\bigr)\bigl(f(x)-1\bigr)$, and $f(x)+1=f(x)-f(2)$.
 \begin{align*}
 \lim_{x\to2}\frac{\bigl(f(x)\bigr)^2-1}{x-2}&=\lim_{x\to2}\frac{f(x)-f(2)}{x-2}\cdot\lim_{x\to2}\bigl(f(x)-1\bigr)\\
 &=f'(2)\cdot\bigl(f(2)-1\bigr) &&\why{definition of $f'(2)$; $f$ is continuous at $2$}\\
 &=(2)(-1-1)=-4.
 \end{align*}
 \item
 \begin{enumerate}[label=(\roman*)]
  \item $P$ is twice differentiable $\imp$ $P'$ is differentiable $\imp$ $P'$ is continuous on $[5,8]$.
  \[
  P'(8)=55<65<90=P'(5).
  \]
  By the Intermediate Value Theorem there is a time $t$ in $(5,8)$ with $P'(t)=65$; this $t$ is not $2$. (The interval $[0,2]$ gives only $t=2$ itself.)
  \item Yes. $P'$ is continuous on $[0,2]$ (as in (i)), and
  \[
  P'(0)=40<50<65=P'(2).
  \]
  By the Intermediate Value Theorem there is a time $t$ in $(0,2)$ with $P'(t)=50$.
 \end{enumerate}
\end{enumerate}
}

\solitem{f:pr:mustbe}{%
\begin{itemize}
 \item I: false. At a removable discontinuity $f(a)$ is defined, but $\displaystyle\lim_{x\to a}f(x)\ne f(a)$.
 \item II: true. This is the definition of continuity.
 \item III: true. Differentiable at $a$ $\imp$ continuous at $a$.
\end{itemize}
The answer is (C).
}

\solitem{f:pr:couldfalse}{%
\begin{itemize}
 \item (A): $f$ is differentiable at $3$ $\imp$ continuous at $3$, so both one-sided limits equal $f(3)=8$. Must be true.
 \item (B) and (C): $8=f(3)$, so both are the definition of $f'(3)=5$. Must be true.
 \item (D): this says $f'$ is continuous at $3$, which is not given. Could be false.
\end{itemize}
The answer is (D).
}

\solitem{f:pr:notdiff}{%
\begin{itemize}
 \item (A) and (B): $f(x)=|x-5|$ is not differentiable at $5$, but it is continuous there and $\displaystyle\lim_{x\to5}f(x)=0$. Not necessarily true.
 \item (C): ``not differentiable at $5$'' means exactly that this limit, the definition of $f'(5)$, does not exist. Must be true.
 \item (D): the graph could instead have a cusp, a vertical tangent, a jump or a hole at $x=5$. Not necessarily true.
\end{itemize}
The answer is (C).
}

\solitem{f:pr:cont3}{%
Continuous at $3$ means $\displaystyle\lim_{x\to3}f(x)=f(3)$, that is, both one-sided limits equal $f(3)$. The answer is (B).
\begin{itemize}
 \item (A) says the opposite of (B).
 \item (C): continuous does not give differentiable; $f(x)=|x-3|$ has no derivative at $3$.
 \item (D): $f(x)=x$ is continuous at $3$ with $f'(x)=1>0$ everywhere.
\end{itemize}
}

\solitem{f:pr:defined}{%
\begin{itemize}
 \item (A): a limit can exist at a hole, where $f(1)$ is not defined; for example $f(x)=\dfrac{x^2-1}{x-1}$ has $\displaystyle\lim_{x\to1}f(x)=2$, but $f(1)$ is not defined.
 \item (B): continuity at $1$ needs $f(1)$ (it says $\displaystyle\lim_{x\to1}f(x)=f(1)$).
 \item (C): $f'(1)=\displaystyle\lim_{x\to1}\dfrac{f(x)-f(1)}{x-1}$ uses $f(1)$. (Also, differentiable $\imp$ continuous.)
 \item (D): the tangent line passes through the point $\bigl(1,f(1)\bigr)$ on the graph.
\end{itemize}
The answer is (A).
}

\solitem{f:pr:graphhole}{%
At $x=a$ the graph has a hole (open dot) on the curve and the value $f(a)$ (filled dot) below it.
\begin{itemize}
 \item From both sides $f(x)$ approaches the height of the hole, so $\displaystyle\lim_{x\to a}f(x)$ exists ((B) and (D) are false), but it is not $f(a)$. So (C) is true.
 \item (A): $f$ is not continuous at $a$ $\imp$ $f$ is not differentiable at $a$, so the limit in (A), the definition of $f'(a)$, does not exist.
\end{itemize}
The answer is (C).
}

\solitem{f:pr:limtable}{%
$f$ is (twice) differentiable, so $f$ is continuous at $2$: $\displaystyle\lim_{x\to2}f(x)=f(2)=3$. The bottom tends to $2\ne0$, so substitute.
\[
\lim_{x\to2}\frac{f(x)}{x}=\frac{f(2)}{2}=\frac32.
\]
The value $f'(2)=4$ is not needed: the limit is not $0/0$.
}

\solitem{d:pr:jumppiece}{%
Check continuity first:
\begin{align*}
\lim_{x\to5^-}f(x)&=\lim_{x\to5^-}\bigl(x^2-20\bigr)=5,\\
f(5)=\lim_{x\to5^+}f(x)&=\lim_{x\to5^+}\bigl(-x^2+20\bigr)=-5.
\end{align*}
$5\ne-5$: a jump. $f$ is not continuous at $5$, so it is not differentiable there. The answer is (A).

Choice (B) is a trap: the slopes $2x=10$ and $-2x=-10$ differ, but the pieces do not meet, so the graph has no corner.
}

\solitem{g:pr:makevar}{%
\begin{enumerate}[label=(\alph*)]
 \item Match the value and the slope at $x=2$:
 \begin{align*}
 4a-1&=k=2b+3 &&\why{the value}\\
 4a&=b &&\why{the slopes: $g'(x)=2ax$, $h'(x)=b$}
 \end{align*}
 So $4a-1=2(4a)+3$, which gives $a=-1$, then $b=-4$ and $k=-5$. (Method 2 gives the same equations: the top of each one-sided quotient must tend to $0$, then $\displaystyle\lim_{x\to2^-}\frac{a(x^2-4)}{x-2}=4a$ and $\displaystyle\lim_{x\to2^+}\frac{b(x-2)}{x-2}=b$.)
 \item Slopes first.
 \begin{align*}
 \lim_{x\to\pi^-}f'(x)&=\lim_{x\to\pi^-}a\cos x=-a,\qquad \lim_{x\to\pi^+}f'(x)=\lim_{x\to\pi^+}(1-\sin x)=1,\\
 -a&=1\imp a=-1, &&\why{the slopes match}\\
 a\sin\pi+b&=\pi+\cos\pi &&\why{the values match}\\
 b&=\pi-1.
 \end{align*}
 So $a=-1$ and $b=\pi-1$.
 \item At $x=1$ the left piece has value $1+3=4$ and slope $2x=2$. So $g$ needs $g(1)=4$ and $g'(1)=2$.
 \begin{center}
 \renewcommand{\arraystretch}{1.9}
 \begin{tabular}{|c|c|c|c|}\hline
 $g(x)$ & $g(1)$ & $g'(x)$ & $g'(1)$\\ \hline
 $4x$ & $4$ & $4$ & $4$\\ \hline
 $x^3+3$ & $4$ & $3x^2$ & $3$\\ \hline
 $4\sqrt x$ & $4$ & $\dfrac{2}{\sqrt x}$ & $2$\\ \hline
 $2x^2+2$ & $4$ & $4x$ & $4$\\ \hline
 \end{tabular}
 \end{center}
 Only $g(x)=4\sqrt x$ matches both.
\end{enumerate}
}

\solitem{g:pr:ksqrt}{%
Here both equations contain both unknowns. $\dfrac{d}{dx}\bigl(k\sqrt x\bigr)=\dfrac{k}{2\sqrt x}$ and $\dfrac{d}{dx}(mx+2)=m$.
\begin{align*}
k\sqrt4&=4m+2\imp 2k=4m+2 &&\why{the values match}\\
\frac{k}{2\sqrt4}&=m\imp m=\frac k4 &&\why{the slopes match}\\
2k&=4\cdot\frac k4+2=k+2\imp k=2,\qquad m=\frac24=\frac12.
\end{align*}
So $k=2$ and $m=\dfrac12$. (Check: both pieces give $4$ at $x=4$, and both slopes are $\dfrac12$.)
}

\Needspace{14\baselineskip}
\subsection*{Exercises}\addcontentsline{toc}{subsection}{Exercises}
\Needspace{12\baselineskip}
\subsubsection*{Section 1: Derivative Rules}
\solitem{ex1:cubic}{%
\begin{align*}
f'(x)&=3x^2-2x+1\\
f'(2)&=12-4+1=9.
\end{align*}
The answer is (C). Choice (A) is $f(2)=8-4+2-1$.
}

\solitem{ex1:dfthree}{%
$f(2)=8+4=12$ is a number, and the derivative of a number is $0$:
\[
\frac{d}{dx}\bigl(f(2)\bigr)=\frac{d}{dx}(12)=0.
\]
The answer is (A). Choice (B) is $f(2)$; choice (C) is $f'(2)=3(2)^2+2=14$.
}

\solitem{ex1:expsin}{%
\begin{align*}
f'(x)&=\frac{d}{dx}\bigl(e^x\bigr)\cdot\sin x+\frac{d}{dx}(\sin x)\cdot e^x &&\why{product rule}\\
&=e^x\sin x+e^x\cos x=e^x(\sin x+\cos x).
\end{align*}
The answer is (C). Choice (A) is the product of the derivatives.
}

\solitem{ex1:cosln}{%
\begin{align*}
f'(x)&=\frac{\dfrac{d}{dx}(\cos x)\cdot\ln x-\dfrac{d}{dx}(\ln x)\cdot\cos x}{(\ln x)^2} &&\why{quotient rule}\\
&=\frac{(-\sin x)\ln x-\dfrac1x\cos x}{(\ln x)^2}\\
&=\frac{-x\sin x\ln x-\cos x}{x(\ln x)^2}. &&\why{multiply the top and the bottom by $x$}
\end{align*}
The answer is (B).
}

\solitem{ex1:cot}{%
\[
\frac{d}{dx}(\cot x)=-\csc^2x.
\]
The answer is (B).
}

\solitem{ex1:big}{%
\begin{enumerate}[label=(\alph*)]
 \item $e^2$ and $\ln3$ are constants.
 \begin{align*}
 y&=4x^{\frac12}-3x^{-2}+2\sin x-\tan x+5e^x-\ln x+2^x+e^2+\ln3\\
 \frac{dy}{dx}&=\frac{2}{\sqrt x}+\frac{6}{x^3}+2\cos x-\sec^2x+5e^x-\frac1x+2^x\ln2.
 \end{align*}
 \item Three factors: three terms.
 \begin{align*}
 \frac{dy}{dx}&=\frac{d}{dx}\bigl(x^3\bigr)\cdot e^x\cos x+\frac{d}{dx}\bigl(e^x\bigr)\cdot x^3\cos x+\frac{d}{dx}(\cos x)\cdot x^3e^x\\
 &=3x^2e^x\cos x+x^3e^x\cos x-x^3e^x\sin x.
 \end{align*}
 \item
 \begin{align*}
 f'(x)&=\frac{\dfrac{d}{dx}\bigl(e^x+\sin x\bigr)\cdot\bigl(x^2+\cos x\bigr)-\dfrac{d}{dx}\bigl(x^2+\cos x\bigr)\cdot\bigl(e^x+\sin x\bigr)}{\bigl(x^2+\cos x\bigr)^2} &&\why{quotient rule}\\
 &=\frac{\bigl(e^x+\cos x\bigr)\bigl(x^2+\cos x\bigr)-\bigl(2x-\sin x\bigr)\bigl(e^x+\sin x\bigr)}{\bigl(x^2+\cos x\bigr)^2}\\
 f'(0)&=\frac{(1+1)(0+1)-(0-0)(1+0)}{1^2}=2.
 \end{align*}
 \item
 \begin{align*}
 f'(x)&=\frac{d}{dx}(x)\cdot\ln x+\frac{d}{dx}(\ln x)\cdot x-3^x\ln3=\ln x+1-3^x\ln3 &&\why{product rule on $x\ln x$}\\
 f'(1)&=0+1-3\ln3=1-3\ln3.
 \end{align*}
 $f(1)=0-3=-3$ is a number, so $\dfrac{d}{dx}\bigl(f(1)\bigr)=0$.
\end{enumerate}
}

\solitem{ex1:frq}{%
\begin{align*}
g'(x)&=\frac{d}{dx}(2x-1)\cdot f(x)+\frac{d}{dx}\bigl(f(x)\bigr)\cdot(2x-1) &&\why{product rule}\\
&=2f(x)+f'(x)(2x-1)\\
g'(3)&=2f(3)+f'(3)\cdot5=2(7)+(-1)(5)=9.
\end{align*}
}

\solitem{ex1:threefns}{%
The top is $f(x)+g(x)$, with derivative $f'(x)+g'(x)$.
\begin{align*}
M'(x)&=\frac{\bigl(f'(x)+g'(x)\bigr)h(x)-h'(x)\bigl(f(x)+g(x)\bigr)}{\bigl(h(x)\bigr)^2} &&\why{quotient rule}\\
M'(0)&=\frac{(-2+5)(2)-(7)(-1+3)}{2^2}=\frac{6-14}{4}=-2.
\end{align*}
The answer is (A).
}

\solitem{ex1:bigtable}{%
\begin{enumerate}[label=(\alph*)]
 \item $g(1)=4$ is a constant, so its derivative is $0$. Use the row $x=0$.
 \begin{itemize}
  \item Top $f(x)g(x)$: value $(2)(3)=6$; derivative $f'(0)g(0)+g'(0)f(0)=(-1)(3)+(1)(2)=-1$.
  \item Bottom $e^x+x$: value $1+0=1$; derivative $e^x+1=2$ at $x=0$.
 \end{itemize}
 \[
 h'(0)=\frac{(-1)(1)-(2)(6)}{1^2}+0=-13.
 \]
 \item Use the row $x=1$.
 \begin{align*}
 k'(x)&=\Bigl[3^x\ln3\cdot f(x)+f'(x)\cdot3^x\Bigr]-\Bigl[\frac{1}{2\sqrt x}\cdot g(x)+g'(x)\cdot\sqrt x\Bigr] &&\why{product rule, twice}\\
 k'(1)&=\bigl[3\ln3\,(-1)+(2)(3)\bigr]-\Bigl[\frac12(4)+(-3)(1)\Bigr]\\
 &=-3\ln3+6-(-1)=7-3\ln3.
 \end{align*}
 \item $f(0)g(0)=(2)(3)=6$ is a number, so $\dfrac{d}{dx}\bigl(f(0)g(0)\bigr)=0$.
\end{enumerate}
}

\solitem{ex1:cuberoot}{%
The bottom is $x^{\frac13}$, with derivative $\dfrac13x^{-\frac23}$.
\begin{align*}
g'(x)&=\frac{f'(x)\,x^{\frac13}-\frac13x^{-\frac23}\cdot f(x)}{x^{\frac23}} &&\why{quotient rule; $\bigl(x^{\frac13}\bigr)^2=x^{\frac23}$}\\
g'(8)&=\frac{5\cdot2-\frac13\cdot\frac14\cdot2}{4}=\frac{10-\frac16}{4}=\frac{59}{24}. &&\why{$8^{\frac13}=2$, $8^{\frac23}=4$}
\end{align*}
The answer is (A).
}

\solitem{ex1:whichrule}{%
\begin{enumerate}[label=(\alph*)]
 \item Two polynomials: expand.
 \begin{align*}
 y&=x^4-1 &&\why{$(a+b)(a-b)=a^2-b^2$}\\
 \frac{dy}{dx}&=4x^3.
 \end{align*}
 \item One term on the bottom: split.
 \begin{align*}
 y&=x^2-4x^{-\frac12} &&\why{$\dfrac{x^{\frac12}}{x}=x^{-\frac12}$}\\
 \frac{dy}{dx}&=2x+2x^{-\frac32}=2x+\frac{2}{x^{\frac32}}.
 \end{align*}
 \item Two terms on the bottom: quotient rule.
 \begin{align*}
 \frac{dy}{dx}&=\frac{\dfrac{d}{dx}(x)\cdot\bigl(x^2+1\bigr)-\dfrac{d}{dx}\bigl(x^2+1\bigr)\cdot x}{\bigl(x^2+1\bigr)^2}\\
 &=\frac{1\cdot\bigl(x^2+1\bigr)-2x\cdot x}{\bigl(x^2+1\bigr)^2}=\frac{1-x^2}{\bigl(x^2+1\bigr)^2}.
 \end{align*}
 \item $\sin^2x=\sin x\cdot\sin x$: product rule.
 \begin{align*}
 \frac{dy}{dx}&=\frac{d}{dx}(\sin x)\cdot\sin x+\frac{d}{dx}(\sin x)\cdot\sin x\\
 &=\cos x\sin x+\cos x\sin x=2\sin x\cos x.
 \end{align*}
 \item A constant over a power: power rule, not the quotient rule.
 \begin{align*}
 y&=\frac52x^{-3}\\
 \frac{dy}{dx}&=\frac52(-3)x^{-4}=-\frac{15}{2x^4}.
 \end{align*}
\end{enumerate}
}
\Needspace{12\baselineskip}
\subsubsection*{Section 2: What the Derivative Means}
\solitem{a:ex2:count}{%
\steph{Step 1: Turn the question into an equation.}
\[
\frac{f(6)-f(a)}{6-a}=0\iff f(a)=f(6)=1.
\]
So count the values of $a$ with $-4\le a<6$ and $f(a)=1$.
\steph{Step 2: Look at each piece of the graph.}
\begin{itemize}
 \item On $[-4,0]$ the graph is the segment through $(-3,0)$ and $(0,2)$, so $f(x)=\dfrac23(x+3)$. Then $\dfrac23(a+3)=1$ gives $a=-\dfrac32$.
 \item On $(0,3)$, $f$ decreases from $2$ to $0$ and is continuous, so it takes the value $1$ exactly once (Intermediate Value Theorem).
 \item On $[3,6)$, $f$ increases from $0$ towards $1$ and stays below $1$; it reaches $1$ only at $x=6$, which is excluded.
\end{itemize}
\steph{Answer.} Two values of $a$.
}

\solitem{a:ex2:mixed}{%
$f'(x)=4x^3-5$.
\begin{enumerate}[label=(\alph*)]
 \item $\dfrac{f(3)-f(0)}{3-0}=\dfrac{(81-15)-0}{3}=\dfrac{66}{3}=22$.
 \item $f'(1)=4-5=-1$.
 \item The $6$ is a hidden $f(2)$: $f(2)=16-10=6$. So the limit is $f'(2)=4(8)-5=27$.
 \item $f(2)=6$ is a number, so $\dfrac{d}{dx}\bigl(f(2)\bigr)=\dfrac{d}{dx}(6)=0$. (It is not $f'(2)=27$.)
 \item The definition of $g'\bigl(\frac{\pi}{4}\bigr)$: $\sec^2\dfrac{\pi}{4}=\bigl(\sqrt2\bigr)^2=2$.
\end{enumerate}
}

\solitem{a:ex2:forms}{%
\begin{itemize}
 \item I: $3=f(2)$, so this is $\displaystyle\lim_{h\to0}\frac{f(2+h)-f(2)}{h}=f'(2)=5$. True.
 \item II: multiply the top and the bottom by $-1$:\quad $\dfrac{3-f(x)}{2-x}=\dfrac{f(x)-f(2)}{x-2}$, so the limit is $f'(2)=5$. True.
 \item III: add and subtract $f(2)$ on the top.
 \begin{align*}
 \lim_{h\to0}\frac{f(2+h)-f(2-h)}{h}&=\lim_{h\to0}\frac{f(2+h)-f(2)}{h}+\lim_{h\to0}\frac{f(2)-f(2-h)}{h}\\
 &=f'(2)+f'(2)=10. &&\why{the second limit: $k=-h$}
 \end{align*}
 Not $5$.
\end{itemize}
The answer is (B).
}

\solitem{b:ex2:slopefield}{%
\begin{enumerate}[label=(\alph*)]
 \item Substitute the point $x=1$, $y=f(1)=4$.
 \[
 f'(1)=\frac{3(1)^2+1}{2(4)}=\frac48=\frac12.
 \]
 \item Find $g(1)$ and $g'(1)$, using $f(1)=4$ and $f'(1)=\dfrac12$.
 \begin{align*}
 g(1)&=1^2\cdot f(1)=4\\
 g'(x)&=\frac{d}{dx}\bigl(x^2\bigr)\cdot f(x)+\frac{d}{dx}\bigl(f(x)\bigr)\cdot x^2=2xf(x)+f'(x)\,x^2 &&\why{product rule}\\
 g'(1)&=2(1)(4)+\frac12(1)=\frac{17}{2}\\
 y-4&=\frac{17}{2}(x-1).
 \end{align*}
 \item $g(1)=4$ is a number, and the derivative of a number is $0$:
 \[
 \frac{d}{dx}\bigl(g(1)\bigr)=\frac{d}{dx}(4)=0.
 \]
 It is not $g'(1)=\dfrac{17}{2}$.
\end{enumerate}
}

\solitem{b:ex2:xn}{%
\steph{Step 1: the tangent line.}
$\dfrac{dy}{dx}=nx^{n-1}$, so the slope at $x=1$ is $n\cdot1^{n-1}=n$.
\[
\ell\colon\quad y-1=n(x-1).
\]
\steph{Step 2: where $\ell$ meets the $x$-axis.}
\[
0-1=n(x-1)\imp x-1=-\frac1n\imp x=1-\frac1n.
\]
\steph{Step 3: the area.}
$T$ has vertices $\Bigl(1-\dfrac1n,0\Bigr)$, $(1,0)$ and $(1,1)$: a right triangle with base $1-\Bigl(1-\dfrac1n\Bigr)=\dfrac1n$ and height $1$.
\[
\text{Area of }T=\frac12\cdot\frac1n\cdot1=\frac{1}{2n}.
\]
}

\solitem{b:ex2:order}{%
\begin{itemize}
 \item At $x=-1$ the graph is at the bottom of a valley: $f'(-1)=0$.
 \item At $x=0$ the graph rises: $f'(0)>0$.
 \item At $x=1$ the graph falls: $f'(1)<0$.
\end{itemize}
So $f'(1)<f'(-1)<f'(0)$. The answer is (C).
}

\solitem{c:xshoppers}{%
\begin{enumerate}[label=(\alph*)]
 \item $S(t)$ counts the shoppers who have entered so far, so the rate at which they enter is $S'(t)$.
 \begin{align*}
 S'(t)&=2t^3-48t^2+288t\\
 S'(3)&=54-432+864=486.
 \end{align*}
 Shoppers are entering the store at $486$ shoppers per hour.
 \item $S(3)$ is a number ($S(3)=40.5-432+1296=904.5$), and the derivative of a number is $0$:
 \[
 \frac{d}{dt}\bigl(S(3)\bigr)=0.
 \]
 It is not $S'(3)$.
 \item Find the average rate, then solve $S'(t)=$ that value.
 \steph{Step 1: The average rate of change.}
 \[
 \frac{S(12)-S(0)}{12-0}=\frac{3456-0}{12}=288. \qquad\why{$S(12)=10368-27648+20736$}
 \]
 \steph{Step 2: Solve $S'(t)=288$.}
 \[
 2t^3-48t^2+288t=288\imp t=1.245,\ 7.666 \quad\text{in } 0<t<12.
 \]
 \scratch{\texttt{Y1}$=$\texttt{2X\^{}3-48X\^{}2+288X}, \texttt{Y2}$=288$, window $0\le x\le12$, \texttt{5:intersect} twice.}
 (The third solution, $t=15.089$, is outside the interval.)
\end{enumerate}
}

\solitem{c:xwindchill}{%
\begin{enumerate}[label=(\alph*)]
 \item
 \begin{align*}
 W'(v)&=-22.1(0.16)v^{-0.84}=-3.536v^{-0.84} &&\why{power rule}\\
 W'(20)&=-0.286.
 \end{align*}
 $W'(20)=-0.286$ degrees Fahrenheit per mph.
 \item Find the average rate, then solve $W'(v)=$ that value.
 \steph{Step 1: The average rate of change.}
 \[
 \frac{W(60)-W(5)}{60-5}=-0.254\ \text{degrees Fahrenheit per mph}.
 \]
 \steph{Step 2: Solve $W'(v)=-0.254$ (stored value).}
 \[
 -3.536v^{-0.84}=-0.25380\ldots\imp v=23.011.
 \]
 \scratch{store the average rate in \texttt{A}; \texttt{Y1}$=$\texttt{-3.536X\^{}(-0.84)}, \texttt{Y2}$=$\texttt{A}, \texttt{5:intersect}.}
\end{enumerate}
}

\solitem{c:xhoriz}{%
\begin{enumerate}[label=(\alph*)]
 \item Horizontal tangent: $g'(x)=0$.
 \[
 \sin\Bigl(x+\frac1x\Bigr)=0\imp x=0.163,\ 0.359 \quad\text{in } 0.12\le x\le1.
 \]
 \scratch{\texttt{Y1}$=$\texttt{sin(X+1/X)}, window $0.12\le x\le1$, \texttt{2:zero} twice.}
 (Check: $x+\dfrac1x=2\pi$ and $\pi$ there; $x=1$ gives $\sin2\ne0$.)
 \item Parallel tangent lines have equal slopes: $f'(x)=h'(x)$.
 \begin{align*}
 2x-2&=1+\pi\cos(\pi x) &&\why{$f'(x)=2x-2$}\\
 x&=0.676. &&\why{\texttt{5:intersect} on $0<x<1$}
 \end{align*}
\end{enumerate}
}

\Needspace{12\baselineskip}
\subsubsection*{Section 3: Differentiability and Continuity}
\solitem{d:pr:gsemicircle}{%
For $-4<x<0$ the graph is the segment from $(-4,2)$ to $(0,0)$:
\[
\lim_{x\to0^-}g'(x)=\frac{0-2}{0-(-4)}=-\frac12.
\]
For $0<x<6$ the graph is the lower semicircle with centre $(3,0)$ and radius $3$. At $(0,0)$ its tangent is vertical, and it goes down, so
\[
\lim_{x\to0^+}g'(x)=\DNEninf.
\]
}

\solitem{d:pr:velocity}{%
\begin{enumerate}[label=(\alph*)]
 \item $v'(4)$ does not exist. The graph has a corner at $(4,20)$:
 \[
 \lim_{h\to0^-}\frac{v(4+h)-v(4)}{h}=\frac{20-0}{4-0}=5,\qquad \lim_{h\to0^+}\frac{v(4+h)-v(4)}{h}=0,
 \]
 and $5\ne0$.

 $v'(20)$ is the slope of the segment from $(16,20)$ to $(24,0)$:
 \[
 v'(20)=\frac{0-20}{24-16}=-\frac52 \text{ meters per second per second}.
 \]
 \item The slopes of the three segments:
 \[
 v'(t)=\begin{cases}5 & \text{for } 0<t<4\\ 0 & \text{for } 4<t<16\\ -\dfrac52 & \text{for } 16<t<24\end{cases}
 \]
 $v'(t)$ is undefined at $t=4$ and $t=16$ (corners).
\end{enumerate}
}

\solitem{d:pr:symdq}{%
$f$ is even, so $f(-x)=f(x)$ and $\dfrac{f(x)-f(-x)}{2x}=0$ for every $x\ne0$: (B) is true.

At $x=0$ the graph has a corner (slope $-1$ on the left, $1$ on the right), so
\[
\lim_{x\to0^-}\frac{f(x)-f(0)}{x}=-1\ne1=\lim_{x\to0^+}\frac{f(x)-f(0)}{x},
\]
and the limit in (A) does not exist. (A) is false.

(C): near $x=2$ the slope is $\frac{3-0}{3-0}=1$. True. (D): $x=3$ is a corner (slopes $1$ and $-1$). True.

The answer is (A). The symmetric quotient in (B) has a limit even though $f'(0)$ does not exist.
}

\solitem{d:ex3:symq}{%
\begin{enumerate}[label=(\alph*)]
 \item Yes. $f(2)=2\cdot0=0$ and $\displaystyle\lim_{x\to2}x\,|x-2|=2\cdot0=0=f(2)$.
 \item $f(2+h)-f(2)=(2+h)\,|h|$.
 \twosides{%
 \begin{align*}
 \lim_{h\to0^-}\frac{(2+h)|h|}{h}&=\lim_{h\to0^-}\frac{(2+h)(-h)}{h}\\
 &=\lim_{h\to0^-}\bigl(-(2+h)\bigr)=-2
 \end{align*}}{%
 \begin{align*}
 \lim_{h\to0^+}\frac{(2+h)|h|}{h}&=\lim_{h\to0^+}\frac{(2+h)h}{h}\\
 &=\lim_{h\to0^+}(2+h)=2
 \end{align*}}
 $-2\ne2$, so $f'(2)$ does not exist (a corner).
 \item $f(2+h)=(2+h)|h|$ and $f(2-h)=(2-h)|-h|=(2-h)|h|$.
 \begin{align*}
 \lim_{h\to0}\frac{f(2+h)-f(2-h)}{2h}&=\lim_{h\to0}\frac{(2+h)|h|-(2-h)|h|}{2h}\\
 &=\lim_{h\to0}\frac{2h\,|h|}{2h}=\lim_{h\to0}|h|=0.
 \end{align*}
 No contradiction: this symmetric quotient is not the definition of $f'(2)$. It averages the slopes on the two sides, $\frac12(-2+2)=0$, and can have a limit at a corner.
\end{enumerate}
}

\solitem{d:ex3:powers}{%
\begin{enumerate}[label=(\alph*)]
 \item $\dfrac{f(h)-f(0)}{h}=\dfrac{h^{1/3}}{h}=\dfrac{1}{h^{2/3}}$, and $h^{2/3}=\bigl(\sqrt[3]{h}\bigr)^2>0$ on both sides:
 \[
 \lim_{h\to0^-}\frac{1}{h^{2/3}}=\DNEinf,\qquad \lim_{h\to0^+}\frac{1}{h^{2/3}}=\DNEinf.
 \]
 Not differentiable; both sides go to $\infty$: a vertical tangent.
 \item $\dfrac{g(h)-g(0)}{h}=\dfrac{h^{2/3}}{h}=\dfrac{1}{h^{1/3}}$, which has the sign of $h$:
 \[
 \lim_{h\to0^-}\frac{1}{h^{1/3}}=\DNEninf,\qquad \lim_{h\to0^+}\frac{1}{h^{1/3}}=\DNEinf.
 \]
 Not differentiable; the sides go to $-\infty$ and $\infty$: a cusp.
 \item $\dfrac{k(h)-k(0)}{h}=\dfrac{h\,|h|}{h}=|h|$, and $\displaystyle\lim_{h\to0^-}|h|=\lim_{h\to0^+}|h|=0$.
 Differentiable, with $k'(0)=0$ (the graph is smooth at the origin, although $|x|$ has a corner there).
\end{enumerate}
}

\solitem{d:pr:rootabs}{%
$f(2)=0$ and $\displaystyle\lim_{x\to2}\sqrt{|x-2|}=0$, so $f$ is continuous at $2$ (this rules out (C) and (D)). With $f(2+h)-f(2)=\sqrt{|h|}$,
\[
\lim_{h\to0^+}\frac{\sqrt{|h|}}{h}=\lim_{h\to0^+}\frac{1}{\sqrt h}=\DNEinf,
\]
so $f'(2)$ does not exist (a cusp). The answer is (A).
}

\solitem{e:pr:samefn}{%
\steph{Is $g$ differentiable at $x=0$?} $g(0)=-7$, from $2x^2+x-7$.
\begin{align*}
\lim_{x\to0^-}\frac{g(x)-g(0)}{x-0}&=\lim_{x\to0^-}\frac{\bigl(2x^2-x-7\bigr)+7}{x}=\lim_{x\to0^-}(2x-1)=-1,\\
\lim_{x\to0^+}\frac{g(x)-g(0)}{x-0}&=\lim_{x\to0^+}\frac{\bigl(2x^2+x-7\bigr)+7}{x}=\lim_{x\to0^+}(2x+1)=1.
\end{align*}
$-1\ne1$, so $g$ is not differentiable at $x=0$.
\steph{Answer} No. $f$ is differentiable for all real numbers (and $f'(0)=0$), but $g$ is not differentiable at $x=0$. So $f$ and $g$ cannot be the same function.
}

\solitem{f:ex3:table}{%
$f$ is twice differentiable $\imp$ $f$ and $f'$ are both differentiable $\imp$ $f$ and $f'$ are both continuous.
\begin{enumerate}[label=(\alph*)]
 \item $f$ is continuous at $4$: $\displaystyle\lim_{x\to4}f(x)=f(4)=5$. The bottom tends to $4-2=2\ne0$, so substitute.
 \[
 \lim_{x\to4}\frac{f(x)-5}{x-2}=\frac{f(4)-5}{4-2}=\frac{0}{2}=0.
 \]
 \item Now the bottom tends to $0$. The $5$ is a hidden $f(4)$, so this is the definition of $f'(4)$:
 \[
 \lim_{x\to4}\frac{f(x)-5}{x-4}=\lim_{x\to4}\frac{f(x)-f(4)}{x-4}=f'(4)=2.
 \]
 \item Yes. $f$ is continuous on $[1,2]$ and $f(1)=3<4<5=f(2)$. By the Intermediate Value Theorem there is a $c$ in $(1,2)$ with $f(c)=4$.
 \item $f'$ is continuous on $[1,6]$. Use the Intermediate Value Theorem with the value $0$ on three intervals that do not overlap:
 \begin{align*}
 f'(2)&=-1<0<4=f'(1) &&\why{a $c$ in $(1,2)$}\\
 f'(2)&=-1<0<2=f'(4) &&\why{a $c$ in $(2,4)$}\\
 f'(6)&=-3<0<2=f'(4) &&\why{a $c$ in $(4,6)$}
 \end{align*}
 The three intervals do not overlap, so these are three different values of $c$, $1<c<6$, with $f'(c)=0$.
\end{enumerate}
}

\solitem{f:ex3:limexists}{%
\textbf{(a)} I only.
\begin{itemize}
 \item I: write $f(x)-5$ as the given quotient times $x-1$ (the proof in the notes):
 \[
 \lim_{x\to1}\bigl(f(x)-5\bigr)=\lim_{x\to1}\frac{f(x)-5}{x-1}\cdot(x-1)=3\cdot0=0,
 \]
 so $\displaystyle\lim_{x\to1}f(x)=5$. True.
 \item II and III: the limit does not use the value $f(1)$. Take $f(x)=3x+2$ for $x\ne1$ and $f(1)=7$. Then $\dfrac{f(x)-5}{x-1}=\dfrac{3(x-1)}{x-1}=3$ for $x\ne1$, so the limit is $3$; but $f(1)=7\ne5$, and $f$ is not continuous at $1$, so $f'(1)$ does not exist. Both could be false.
\end{itemize}}
\begin{center}
\begin{asy}
unitsize(1.3cm,0.42cm);
axes(-0.6,2.6,-0.5,8.4);
xnums(new real[]{1,2});
ynums(new real[]{5,7});
curve(new real(real x){return 3x+2;}, -0.3, 1.9);
opendot((1,5)); closeddot((1,7));
label("$O$",(0,0),SW,fontsize(8pt));
\end{asy}
\par\vspace{2pt}{\small $f(x)=3x+2$ for $x\ne1$, $f(1)=7$: the limit is $3$, but $f(1)\ne5$.}
\end{center}
The answer is (A).

\textbf{(b)} $f$ is continuous at $1$, so $f(1)=\displaystyle\lim_{x\to1}f(x)=5$ by (a). Then the $5$ is $f(1)$, and the given limit is the definition of $f'(1)$:
\[
f'(1)=\lim_{x\to1}\frac{f(x)-f(1)}{x-1}=\lim_{x\to1}\frac{f(x)-5}{x-1}=3.
\]
So $f(1)=5$ and $f'(1)=3$.

\solitem{f:pr:ivtW}{%
$W$ is twice differentiable $\imp$ $W'$ is differentiable $\imp$ $W'$ is continuous on $[22,30]$.
\[
W'(30)=0.5<0.7<1.0=W'(22).
\]
By the Intermediate Value Theorem there is a time $t$ in $(22,30)$ with $W'(t)=0.7$ GL/day; this $t$ is not $10$. (The interval $[0,10]$ gives only $t=10$ itself. The interval $[10,22]$ does not work either: $0.7$ is an endpoint value there.)
}

\solitem{f:pr:ivtr}{%
Yes. $r$ is twice differentiable $\imp$ $r'$ is differentiable $\imp$ $r'$ is continuous on $[0,3]$.
\[
r'(0)=-6.1<-6<-5.0=r'(3).
\]
By the Intermediate Value Theorem there is a time $t$ in $(0,3)$ with $r'(t)=-6$.
}

\solitem{e:ex3:fourconst}{%
At each join the values and the slopes must match. The middle piece has $\dfrac{d}{dx}(\cos x+x)=-\sin x+1$.
\steph{At $x=0$}
\begin{align*}
a&=-\sin0+1=1 &&\why{the slopes match}\\
a(0)+b&=\cos0+0\imp b=1 &&\why{the values match}
\end{align*}
\steph{At $x=\pi$}
\begin{align*}
c&=-\sin\pi+1=1 &&\why{the slopes match}\\
c\pi+d&=\cos\pi+\pi\imp\pi+d=\pi-1\imp d=-1 &&\why{the values match}
\end{align*}
\steph{Answer} $a=1$, $b=1$, $c=1$, $d=-1$. (The outer pieces $y=x+1$ and $y=x-1$ are the tangent lines to $y=\cos x+x$ at $x=0$ and $x=\pi$.)
}

\solitem{g:ex3:joinpoint}{%
\begin{enumerate}[label=(\alph*)]
 \item Each piece is differentiable, so only the join $x=a$ matters. Slopes first.
 \begin{align*}
 3a^2&=12\imp a=2\text{ or }a=-2 &&\why{the slopes match}\\
 a=2:&\quad 2^3=8,\quad 12(2)-16=8 &&\why{the values match}\\
 a=-2:&\quad (-2)^3=-8,\quad 12(-2)-16=-40 &&\why{the values differ: a jump}
 \end{align*}
 So $a=2$ only. The line $y=12x-16$ is the tangent line to $y=x^3$ at $x=2$; at $x=-2$ the slopes match but the graph jumps.
\end{enumerate}
}
\begin{center}
 \begin{asy}
 unitsize(0.55cm,0.11cm);
 axes(-2.8,3.2,-28,30);
 xnums(new real[]{-2,-1,1,2}); ynums(new real[]{-20,-10,10,20});
 {guide g; for (int i=0; i<=100; ++i) {real x=-2.8+(3.05-(-2.8))*i/100; g=g..(x,x^3);} draw(g, gray(0.6)+dashed+linewidth(0.6pt));}
 curve(new real(real x){return x^3;},-2.8,2);
 draw((2,8)--(3.2,22.4), rgb(1,0.5,0)+linewidth(0.9pt));
 closeddot((2,8));
 label("$y=x^3$",(-2.4,-16.6),E,fontsize(8pt));
 label("$y=12x-16$",(2.75,12),E,fontsize(8pt));
 \end{asy}
 \end{center}
\begin{enumerate}[label=(\alph*),start=2]
 \item $\dfrac{d}{dx}(\ln x)=\dfrac1x$ and $\dfrac{d}{dx}(mx)=m$.
 \begin{align*}
 \frac1a&=m &&\why{the slopes match}\\
 \ln a&=ma=\frac1a\cdot a=1 &&\why{the values match}\\
 a&=e,\qquad m=\frac1e.
 \end{align*}
 So $a=e$ and $m=\dfrac1e$: the line $y=\dfrac xe$ is the tangent line to $y=\ln x$ at $(e,1)$, and it passes through the origin.
\end{enumerate}

\solitem{g:pr:cd}{%
\begin{align*}
c&=2(2)-c &&\why{slopes: $\tfrac{d}{dx}(cx+d)=c$, $\tfrac{d}{dx}\bigl(x^2-cx\bigr)=2x-c$}\\
c&=2,\\
2c+d&=2^2-2c &&\why{the values match}\\
4+d&=0\imp d=-4.
\end{align*}
So $c+d=2+(-4)=-2$. The answer is (B). Choice (A) is $d$ and (D) is $c$.
}
\fi
\end{document}